\chapter{状态、时钟与同步边界}

第一章讲的是组合逻辑：输入稳定后，输出由当前输入唯一决定。它能计算，却不能保存过去。真正的机器必须能把某个结果留下来，等下一个步骤继续使用；也必须能规定“什么时候更新状态”，否则反馈路径会让硬件在一个周期内连续变化，机器无法分步骤推进。

本章将反馈、锁存器、触发器、时钟、复位、寄存器、计数器和状态机合并讨论。读完后，应当能把“状态”理解成真实硬件边界：一些 bit 在时钟边界被提交，边界之间由组合逻辑准备下一值。后续 CPU 的 PC、指令寄存器、控制状态、流水线寄存器和中断入口，都将反复使用这一章的模型。

\begin{keyidea}
组合逻辑负责准备下一值，状态元件负责在边界保存下一值。同步硬件就是把这两件事反复组织起来。
\end{keyidea}

第一章已经把建立时间、保持时间、时钟到输出和时钟偏斜作为词汇预告。本章从这里正式接上：组合逻辑的输出并不是“最终正确即可”，它必须在目标触发器的采样窗口之前稳定；状态元件也不是理想变量，它采样后需要时间把新值送到 Q 端。同步设计的基本任务，就是把这些时间约定整理成一组可检查的时序关系：时钟到输出说明触发器被时钟触发后 Q 端多久开始并完成变化，它决定下一段组合逻辑最早什么时候开始传播；组合最大延迟说明 Q 端变化经过门和连线到达 D 端的最晚时间，它决定下一个时钟沿到来时结果是否已稳定；建立时间要求 D 端在采样沿之前提前稳定，否则目标触发器会采到过渡值；保持时间要求 D 端在采样沿之后继续稳定，否则新值会过早冲进采样窗口；时钟偏斜说明源触发器和目标触发器看到时钟沿的时间差，它会压缩有效周期、甚至破坏保持窗口。

本章的版面组织遵循同一顺序：先说明反馈为什么能保存一个 bit，再说明锁存器和触发器如何定义写入边界，随后把多个触发器扩展成寄存器、计数器、移位寄存器和状态机，最后把复位、跨边界输入、实验板接线和数据手册参数整合成完整的工程方法。读者不应把这些小节看成独立名词表，而应把它们看成“状态如何被创建、提交、传播和验证”的连续链条。

\section{一、反馈、锁存器与保存一个 bit}

保存一个 bit 的起点是反馈。两个反相器首尾相接，会形成两个稳定状态。若左边节点为 0，右边经过反相后为 1；右边的 1 再反相回来，正好维持左边为 0。反过来，左边为 1、右边为 0 也能稳定。只要没有外部强制改变，这个状态会保持。

该结构说明，“保存一个 bit”对应电路中某些节点在反馈作用下停留于稳定物理状态。一个稳定状态解释为 0，另一个稳定状态解释为 1。状态的本质，是电路有能力在输入不再变化时仍然维持某个内部结果。

从电路角度看，反馈不是抽象箭头，而是输出重新影响输入。两个反相器互相连接时，每一侧都在“支持”另一侧的当前电平。若某个扰动让左侧电压略有上升，右侧反相器会把右侧电压压低，右侧的低电平又通过另一个反相器继续支持左侧高电平。这个正反馈过程让电路迅速落入两个稳定点之一。数字电路把这两个稳定点命名为 0 和 1，但底层仍然是晶体管、电容、电阻和电源共同决定的物理平衡。

\begin{circuit}[x=1cm,y=1cm]
  \node[not gate US, draw=BookAccent, line width=0.55pt, scale=1.35] (inv1) at (1.6,0) {};
  \node[not gate US, draw=BookAccent, line width=0.55pt, scale=1.35] (inv2) at (6.4,0) {};
  \node[hw label, above=0.3cm of inv1] {反相器 A};
  \node[hw label, above=0.3cm of inv2] {反相器 B};
  \coordinate (tap1) at ($(inv1.output)+(0.75,0)$);
  \coordinate (tap2) at ($(inv2.output)+(0.75,0)$);
  \draw[hw bus] (inv1.output) -- (tap1) -- (inv2.input);
  \draw[hw return] (inv2.output) -- (tap2) |- (4.0,-1.15) -| ($(inv1.input)+(-0.65,0)$) -- (inv1.input);
  \node[BookTeal,font=\small\bfseries] at (3.95,0.45) {X};
  \node[BookTeal,font=\small\bfseries] at (7.7,0.35) {Y};
  \node[hw label, text width=8.5cm] at (4.0,-1.85) {两个反相器首尾相接后，任意一个稳定输出都会成为另一侧的输入；反馈线绕到下方，避免穿过门符号和标签。};
\end{circuit}
\diagramnote{两个反相器首尾相接形成双稳态：X=0, Y=1 稳定；X=1, Y=0 也稳定。}

真实锁存结构通常不会直接暴露为两个反相器，而是用交叉耦合的 NAND 门或 NOR 门构成 SR 锁存器。SR 的含义是 set 与 reset：set 请求把状态置为 1，reset 请求把状态置为 0。它已经能表达“外部命令改变内部状态”，但也引入一个重要约束：不能同时给出相互冲突的命令。若 set 和 reset 同时有效，输出可能进入不符合设计语义的状态；当冲突解除时，最终稳定在 0 还是 1 取决于门延迟和释放顺序。这就是第一类状态设计约定：不仅要说明哪个状态合法，还要说明哪些输入组合禁止出现。

把 SR 锁存器的四种输入组合逐行推一遍。两个输入 $\overline{S}$、$\overline{R}$ 都是低有效（为 0 才表示请求），以 NAND 实现为例：

\begin{longtable}{L{0.1\linewidth}L{0.1\linewidth}L{0.14\linewidth}L{0.56\linewidth}}
\toprule
$\overline{S}$ & $\overline{R}$ & Q & 逐行推导 \\
\midrule
0 & 0 & 禁用 & 两个 NAND 的输入都有 0，Q 和 $\overline{Q}$ 都被迫为 1，互补关系被破坏；这是非法输入，不是“存了 1” \\
0 & 1 & 1 & 上侧 NAND 因 $\overline{S}=0$ 输出 1；下侧 NAND 两输入全 1 输出 0，反馈回去维持 Q=1。set 生效 \\
1 & 0 & 0 & 与上一行对称：$\overline{R}=0$ 把 $\overline{Q}$ 拉成 1，反馈把 Q 拉成 0。reset 生效 \\
1 & 1 & 保持 & 无外部请求：若 Q=1，下侧输出 0 支持上侧输出 1；若 Q=0 则反过来。反馈环自己维持旧值 \\
\bottomrule
\end{longtable}

禁行还要单独说一句“释放之后发生什么”。若 $\overline{S}$、$\overline{R}$ 从 0,0 同时回到 1,1，两个 NAND 的旧输出都是 1，于是两个门都打算把输出改成 0——谁先翻转，反馈就支持谁，最终状态取决于两个门的实际延迟差。这就是为什么正文说“稳定在 0 还是 1 取决于门延迟和释放顺序”：它不是玄学，是禁行退出瞬间的一场不可预测的竞争。可靠设计要么不产生 0,0，要么让两个输入不同时释放。

\begin{circuit}[x=1cm,y=1cm]
  \node[nand gate US, draw=BookAccent, line width=0.55pt, logic gate inputs=nn, scale=1.28] (n1) at (3.0,1.25) {};
  \node[nand gate US, draw=BookAccent, line width=0.55pt, logic gate inputs=nn, scale=1.28] (n2) at (3.0,-1.25) {};
  \draw[hw bus] (0.4,1.45) node[left] {$\overline{S}$} -- (n1.input 1);
  \draw[hw bus] (0.4,-1.45) node[left] {$\overline{R}$} -- (n2.input 2);
  \coordinate (q) at (5.25,1.25);
  \coordinate (qn) at (5.25,-1.25);
  \draw[hw bus] (n1.output) -- (q) node[right,BookTeal,font=\small\bfseries,xshift=0.8cm] {Q};
  \draw[hw bus] (n2.output) -- (qn) node[right,BookTeal,font=\small\bfseries,xshift=0.8cm] {$\overline{Q}$};
  \draw[hw return] (q) -- ++(0.62,0) |- (0.78,-0.98) -- (n2.input 1);
  \draw[hw return] (qn) -- ++(0.62,0) |- (0.78,0.98) -- (n1.input 2);
  \node[hw label, text width=6.9cm] at (3.0,-2.35) {交叉反馈只连接到另一只 NAND 的第二输入，输入线、输出线和反馈线各走独立通道。};
\end{circuit}
\diagramnote{由两个交叉耦合的 NAND 门构成的 SR 锁存器。$\overline{S}$ 与 $\overline{R}$ 低有效；当两者都为高时，反馈维持原状态。}

\topic{专题：NOR 版 SR 锁存器——同一结构，另一套输入约定}

上面用 NAND 搭出的 SR 锁存器不是唯一的搭法。把两个 NAND 换成两个 NOR，得到的仍是交叉耦合反馈环，存储机制完全一样；变化只发生在输入约定上。NAND 版的 $\overline{S}$、$\overline{R}$ 是低有效：输入为 0 才表示请求，无请求时两者都停在 1。NOR 版的 S、R 是高有效：输入为 1 才表示请求，无请求时两者都停在 0。极性相反不是设计风格问题，而是门本身的函数决定的：NAND 只要有一个输入为 0 输出就是 1，所以“请求”只能靠 0 来表达；NOR 只要有一个输入为 1 输出就是 0，所以“请求”只能靠 1 来表达。

把 NOR 版逐行推一遍。设上侧 NOR 输出 Q，下侧 NOR 输出 $\overline{Q}$，S 接下侧、R 接上侧。S=1, R=0 时，下侧 NOR 因 S=1 输出 0，上侧 NOR 两个输入全为 0 输出 1，Q 被写成 1，置位生效；S=0, R=1 时对称，Q 被写成 0，复位生效；S=0, R=0 时，两个 NOR 各有一个输入来自对方输出：若 Q=1，下侧输出 0 正好支持上侧输出 1；若 Q=0 则反过来，反馈环维持旧值；S=1, R=1 时，两个 NOR 都被迫输出 0，Q 与 $\overline{Q}$ 的互补关系被破坏，这就是 NOR 版的禁用组合。两版放在一起对照：

\begin{longtable}{L{0.12\linewidth}L{0.12\linewidth}L{0.15\linewidth}L{0.15\linewidth}L{0.34\linewidth}}
\toprule
S（NOR） & R（NOR） & $\overline{S}$（NAND） & $\overline{R}$（NAND） & 锁存器行为 \\
\midrule
0 & 0 & 1 & 1 & 无请求：反馈环保持旧值 \\
1 & 0 & 0 & 1 & 置位：Q 被写成 1 \\
0 & 1 & 1 & 0 & 复位：Q 被写成 0 \\
1 & 1 & 0 & 0 & 禁用：set 与 reset 同时请求，互补输出被破坏 \\
\bottomrule
\end{longtable}

表面看，两版的禁用组合一个是 11、一个是 00，似乎完全不同；实际上它们表示同一件事——set 和 reset 被同时请求。禁用组合指向的从来不是某个二进制图案，而是“两条相互冲突的命令同时有效”这个语义。退出禁用组合时的竞争也一样：NOR 版从 11 同时回到 00，与 NAND 版从 00 同时回到 11 一样，两个门的旧输出都是 0，于是都打算把输出改成 1，谁先翻转反馈就支持谁，最终状态同样取决于门延迟差和释放顺序。

这个对照的意义在于把“结构”和“约定”分开。存储能力来自交叉耦合反馈，它在 NAND 版和 NOR 版里是同一个结构；低有效还是高有效、空闲电平停在 1 还是 0、禁用图案是 00 还是 11，只是套在结构外面的输入约定。数据手册里常见同一功能给出两种极性的型号，读表时应当先确认有效电平，再读功能行；顺序反了，就会把保持行误读成禁用行。

到这里，四种保存 1 bit 的结构可以排成一条线：双反相器反馈用两个稳定点保存 1 bit，但没有独立写入控制；SR 锁存器允许用 set/reset 命令改变状态，代价是 set 与 reset 不能给出冲突请求；D 锁存器用一个数据输入和 enable 表达写入，代价是 enable 有效期间输入可以穿透到输出；D 触发器只在时钟边沿采样并保存，代价是必须满足建立时间、保持时间和复位要求。后面的讨论会逐个展开后三种。
\vspace{0.8em plus 0.25em}

单纯反馈还不足以形成可用存储单元。若无法可靠地把它从 0 改成 1，或从 1 改成 0，它只是一个稳定反馈环。因此需要写入控制。锁存器可以理解为带写使能的 1-bit 存储单元：写使能有效时，输入 D 可以影响输出 Q；写使能无效时，反馈通路维持原来的 Q。

\begin{longtable}{L{0.22\linewidth}L{0.24\linewidth}L{0.42\linewidth}}
\toprule
写使能 & 数据输入 & 输出行为 \\
\midrule
有效 & 0 或 1 & 输出跟随输入变化 \\
无效 & 任意 & 输出保持原值 \\
\bottomrule
\end{longtable}
\vspace{1.0em plus 0.35em}
\begin{block}[x=2.0cm,y=1.6cm]
  \node[draw, rectangle, minimum width=2.4cm, minimum height=2.0cm, fill=BookCodeBg, rounded corners=3pt] (sym) at (0,0) {};
  \node[above=0.12cm of sym, font=\small\color{BookNavy}] {D-Latch};
  \node[left] at ($(sym.west)+(-0.12,0.5)$) {D};
  \node[left] at ($(sym.west)+(-0.12,-0.5)$) {LE};
  \node[right] at ($(sym.east)+(0.12,0.5)$) {Q};
  \node[right] at ($(sym.east)+(0.12,-0.5)$) {Q'};
  \node[below=0.3cm of sym, font=\footnotesize, BookAccent] {LE=1: 透明};
  \node[below=0.6cm of sym, font=\footnotesize, BookTeal] {LE=0: 保持};
\end{block}
\diagramnote{D 锁存器方框符号：LE=1 时 Q 跟随 D（透明），LE=0 时反馈保持原值。}
\vspace{0.8em plus 0.25em}

例如一个“按住才能设置”的电路中，enable 为 1 时，D=1 会把锁存器写成 1；enable 回到 0 后，即使 D 变回 0，Q 仍保持 1。这样，过去某一刻的输入就被硬件保存下来了。

把 SR 锁存器变成 D 锁存器的意义，在于把两个可能冲突的命令合并为一个数据输入 D 和一个写使能 enable。enable 有效且 D=1 时，相当于请求 set；enable 有效且 D=0 时，相当于请求 reset；enable 无效时，set 和 reset 都不被触发，反馈保持旧值。这样，外部控制逻辑不再同时生成 set/reset 两条危险命令，而是生成“是否写入”和“写入什么值”两个更清晰的问题。

这个改造在门级只需要一个反相器和两个 NAND：

\begin{circuit}[x=1cm,y=1cm]
  \node[nand gate US, draw=BookAccent, line width=0.55pt, logic gate inputs=nn, scale=1.2] (g1) at (3.0,1.6) {};
  \node[not gate US, draw=BookAccent, line width=0.55pt, scale=1.1] (inv) at (1.5,0.1) {};
  \node[nand gate US, draw=BookAccent, line width=0.55pt, logic gate inputs=nn, scale=1.2] (g2) at (3.0,-0.7) {};
  \node[hw state, minimum width=2.5cm, minimum height=1.9cm] (sr) at (6.5,0.45) {SR 锁存器\\（交叉耦合 NAND）};
  \node[hw label] at (-1.0,1.9) {D};
  \draw[hw bus] (-0.7,1.9) -- (g1.input 1);
  \draw[fill=BookInk] (0.5,1.9) circle (0.05);
  \draw[hw bus] (0.5,1.9) -- (0.5,0.1) -- (inv.input);
  \draw[hw bus] (inv.output) -- (g2.input 1);
  \node[hw label] at (-1.0,-0.8) {LE};
  \draw[hw bus] (-0.7,-0.8) -- (2.25,-0.8);
  \draw[fill=BookInk] (2.25,-0.8) circle (0.05);
  \draw[hw bus] (2.25,-0.8) -- (2.25,1.35) -- (g1.input 2);
  \draw[hw bus] (2.25,-0.8) -- (2.25,-0.95) -- (g2.input 2);
  \draw[hw bus] (g1.output) -- node[hw label, above, pos=0.55]{$\overline{S}$} (5.25,1.1);
  \draw[hw bus] (g2.output) -- node[hw label, below, pos=0.55]{$\overline{R}$} (5.25,-0.2);
  \draw[hw bus] ($(sr.east)+(0,0.5)$) -- ++(0.9,0) node[right]{Q};
  \draw[hw bus] ($(sr.east)+(0,-0.5)$) -- ++(0.9,0) node[right]{$\overline{Q}$};
\end{circuit}
\diagramnote{D 锁存器的门级实现：LE=1 时，D 和它的反相信号分别推开两扇 NAND，一路 $\overline{S}$ 或 $\overline{R}$ 被拉低（写入）；LE=0 时两扇都关闭，$\overline{S}$、$\overline{R}$ 同为高（保持）。0,0 禁行从结构上不可能再出现——这就是“把 set/reset 两条危险命令合并成 D 和 LE”在电路结构上的直接体现。}

真实器件中的锁存器还会暴露输出使能、低有效清零、三态输出等控制脚。以 74HC373/74HC573 这类透明锁存器为例，锁存使能决定内部是否继续跟随输入，输出使能只决定引脚是否驱动外部总线；二者不能混淆。关闭输出并不等于内部状态消失，打开输出也不等于重新采样输入。后面讨论总线、寄存器堆和 I/O 芯片时，这个区分会反复出现：内部状态、外部驱动、写入边界是三件不同的事。

8086 是理解锁存器用途的经典案例。8086 的部分地址线与数据线复用，地址阶段先在总线上给出地址，随后同一组引脚又可能转为数据传输。外部存储器不能在数据阶段继续把这些引脚当作地址，因此系统板需要由 ALE 这类地址锁存允许信号驱动外部锁存器，把地址阶段的值保存下来。8086 时代常见的板级器件包括 8282/8283 或 74LS373；今天做低速教学实验时也常用 74HC373/74HC573 复现同一思想。这里的锁存器不是“多余缓冲器”，而是把一个总线相位中的地址变成后续相位仍然可用的状态。

\begin{longtable}{L{0.2\linewidth}L{0.36\linewidth}L{0.34\linewidth}}
\toprule
总线阶段 & 锁存器动作 & 设计含义 \\
\midrule
地址有效 & 锁存窗口打开，地址进入锁存器 & 允许复用总线先表达地址 \\
窗口关闭 & 内部反馈保持地址 & 外部存储器继续看到稳定地址 \\
数据传输 & 总线引脚改作数据或状态用途 & 地址不会随复用引脚变化而丢失 \\
输出使能控制 & 决定是否驱动外部总线或后级 & 不改变锁存器内部保存的值 \\
\bottomrule
\end{longtable}
\vspace{0.8em plus 0.25em}

\begin{waveform}[x=1cm,y=1cm]
  \draw[BookMuted, dashed, line width=0.4pt] (2.4,-1.6) -- (2.4,3.4);
  \draw[BookMuted, dashed, line width=0.4pt] (5.4,-1.6) -- (5.4,3.4);
  \node[hw label] at (1.2,3.6) {T1：地址阶段};
  \node[hw label] at (3.9,3.6) {T2：转向};
  \node[hw label] at (7.4,3.6) {T3/Tw：数据阶段};
  \draw[BookAccent, line width=1.0pt] (0,2.2) -- (0.4,2.2) -- (0.4,2.9) -- (2.0,2.9) -- (2.0,2.2) -- (9.4,2.2);
  \node[hw label, anchor=east] at (-0.2,2.55) {ALE};
  \draw[BookTeal, line width=0.9pt] (0,1.35) -- (0.7,0.85) -- (2.1,0.85) -- (2.4,1.1);
  \draw[BookTeal, line width=0.9pt] (0,0.85) -- (0.7,1.35) -- (2.1,1.35) -- (2.4,1.1);
  \draw[BookMuted, dashed, line width=0.7pt] (2.4,1.1) -- (5.6,1.1);
  \draw[BookAmber, line width=0.9pt] (5.6,1.35) -- (9.4,1.35) -- (9.4,0.85) -- (5.6,0.85) -- cycle;
  \node[hw label, anchor=east] at (-0.2,1.1) {AD 总线};
  \node[hw label] at (1.4,1.1) {地址};
  \node[hw label] at (4.0,0.75) {高阻/转向};
  \node[hw label] at (7.5,1.1) {数据};
  \draw[BookAccent, line width=1.0pt] (0,-0.5) -- (2.0,-0.5) -- (2.0,0.2) -- (9.4,0.2);
  \node[hw label, anchor=east] at (-0.2,-0.15) {锁存后地址};
  \node[hw label] at (5.4,-0.85) {ALE 下降沿后，地址由锁存器保持到数据阶段结束};
\end{waveform}
\diagramnote{8086 复用总线的一个读周期：T1 时 ALE 打开锁存窗口，地址随 AD 总线进入锁存器；ALE 落下后窗口关闭，AD 总线转向传数据，而锁存器把地址保持到 T3 采样完成。锁存器把一个总线相位的值变成后续相位仍然可用的状态。}

从这个案例可以看出，锁存器适合用在协议已经提供清晰电平窗口的地方，例如地址锁存、外部总线相位分离、显示输出暂存等。若协议没有给出可靠窗口，或者多个状态级之间存在组合反馈，透明锁存器反而会让调试变难。教材后面讲最小 CPU 时，PC、指令寄存器和控制状态一般优先使用边沿触发器；讲总线和 I/O 时，再把锁存器作为相位保持和输出隔离工具使用。

锁存器的问题在于透明窗口。只要 enable 有效，输入变化就可能一路穿透到输出。若多个锁存器共用同一个 enable，并且它们之间还接着组合逻辑，状态可能在一个打开窗口内穿过多级，边界变得模糊。对小电路来说，锁存器易于分析；对大规模同步机器来说，更清晰的做法是把写入集中到离散时钟边界。

透明窗口不是“锁存器不能用”的理由，而是要求设计者明确时间边界。两相锁存器（主、从两级各用半个时钟周期的锁存器）、流水线透明锁存、片内时钟门控单元都可能利用锁存器，但它们依赖严格的相位、非重叠时钟（两相高低时间不重叠的时钟）和时序分析。初学同步机器时，先使用边沿触发器更稳妥，因为每一级状态只在时钟沿提交，调试时可以把每个周期看成独立的一步：上一步的 Q 经过组合逻辑形成本步的 D，本步边界再把 D 提交为下一步的 Q。

\topic{专题：门控 D 锁存器的透明窗口——为什么计数器和状态机不用它}

门控 D 锁存器在 LE 有效期间是透明的：D 的任何变化经过几级门延迟就出现在 Q 上。这不是缺陷描述，而是工作机制描述——锁存器保存的是窗口关闭那一刻的 D，窗口之内它只是一条带延迟的组合通路。由此推出两个直接后果。第一，晚到的数据会穿透：只要 D 在窗口关闭之前变化，新值就来得及进入并被保存。第二，毛刺也会穿透：D 上的短暂干扰会经过同样的延迟出现在 Q 上，下游看到的是一条和输入几乎一样乱的波形，只是晚了一点。

用一组具体数字感受两种元件的差别。设 LE 窗口宽 20ns（t=0 打开，t=20ns 关闭），锁存器传播延迟约 2ns。D 在 t=17ns 变化，t=19ns 新值已经出现在 Q 上，窗口关闭时保存的就是这个新值——离关闭只剩 3ns 的变化仍然穿透。D 上一个 2ns 宽的毛刺出现在 t=10ns，约 t=12ns 起 Q 上跟着出现一个毛刺。换成 D 触发器看同样的输入：触发器只关心采样沿前后那一小段窗口，设建立时间 0.5ns、保持时间 0.3ns，那么 t=10ns 的毛刺和 t=17ns 的变化只要不落在沿前后共 0.8ns 的窗口里，就不会被采进去。锁存器的敏感区是整个 20ns，触发器的敏感区是 0.8ns，相差 25 倍；这个倍数就是“透明”二字的代价。

把这段过程按时间轴画出来，两种元件的敏感区差别就一目了然：

\begin{waveform}[x=0.38cm,y=0.85cm]
  % 参考线：t=0 窗口打开，t=20 窗口关闭（采样沿）
  \draw[BookMuted, dashed, line width=0.45pt] (0,-0.8) -- (0,4.9);
  \draw[BookMuted, dashed, line width=0.45pt] (20,-0.8) -- (20,4.9);
  \node[hw label, anchor=west] at (0.15,5.12) {窗口打开};
  \node[hw label, anchor=east] at (19.85,5.12) {窗口关闭（采样沿）};
  % LE：t=0 打开，t=20 关闭
  \draw[BookAccent, line width=1.0pt] (-0.8,4.0) -- (0,4.0) -- (0,4.6) -- (20,4.6) -- (20,4.0) -- (23,4.0);
  \node[hw label, anchor=east] at (-1.0,4.3) {LE};
  % D：t=10 出现 2ns 毛刺，t=17 晚到变化
  \draw[BookAmber, line width=1.0pt] (-0.8,2.6) -- (10,2.6) -- (10,3.2) -- (12,3.2) -- (12,2.6) -- (17,2.6) -- (17,3.2) -- (23,3.2);
  \node[hw label, anchor=east] at (-1.0,2.9) {D};
  % 锁存器 Q：延迟约 2ns，毛刺与晚到变化全部穿透
  \draw[BookRed, line width=1.0pt] (-0.8,1.2) -- (12,1.2) -- (12,1.8) -- (14,1.8) -- (14,1.2) -- (19,1.2) -- (19,1.8) -- (23,1.8);
  \node[hw label, anchor=east] at (-1.0,1.5) {锁存器 Q};
  \node[hw label, text=BookRed] at (13,2.15) {毛刺穿透};
  % D 触发器 Q：只认采样沿，沿后一步更新
  \draw[BookTeal, line width=1.1pt] (-0.8,-0.2) -- (20.7,-0.2) -- (20.7,0.4) -- (23,0.4);
  \node[hw label, anchor=east] at (-1.0,0.1) {触发器 Q};
  \draw[{Stealth[length=1.6mm]}-{Stealth[length=1.6mm]}, BookRed, line width=0.55pt] (19.6,-0.62) -- (20.4,-0.62);
  \node[hw label, anchor=west] at (20.6,-0.62) {敏感窗仅 0.8ns};
  % 时间轴与刻度
  \draw[hw bus] (-0.8,-1.35) -- (23.5,-1.35);
  \foreach \x in {0,10,20} {
    \draw[BookMuted, line width=0.5pt] (\x,-1.45) -- (\x,-1.25);
    \node[hw label] at (\x,-1.75) {\x};
  }
  \node[hw label, anchor=east] at (23.5,-1.75) {ns};
\end{waveform}
\diagramnote{同一个 D 送到两种元件：LE 窗口（t=0 到 20ns）内，D 上 2ns 宽的毛刺（t=10ns）和晚到的变化（t=17ns）都经过约 2ns 延迟穿透到锁存器 Q；D 触发器只认 t=20ns 采样沿前后共 0.8ns 的窗口，毛刺完全不可见，晚到变化被正常采入，沿后 Q 一步更新。敏感区 20ns 对 0.8ns，相差 25 倍——这就是透明窗口的代价。}

\begin{longtable}{L{0.28\linewidth}L{0.33\linewidth}L{0.31\linewidth}}
\toprule
输入事件 & 门控 D 锁存器（LE 窗口 20ns） & D 触发器（沿采样） \\
\midrule
D 在窗口打开前已稳定 & Q 早已跟随，窗口关闭时保存 & 沿到来时采样，沿后 Q 更新 \\
D 在窗口中段才变化（晚到数据） & Q 随之更新，保存的是晚到的值 & 距沿超过建立时间则正常采样，但 Q 只在沿后才变 \\
D 在窗口内出现毛刺 & 毛刺经传播延迟出现在 Q 上 & 毛刺不落在沿附近窗口内则完全不可见 \\
D 在窗口外任意变化 & 无影响 & 无影响 \\
\bottomrule
\end{longtable}

为什么计数器和状态机必须用触发器，而不能用共用一个 LE 的透明锁存器？关键在于这两类电路存在经过组合逻辑回到自身输入的反馈。计数器的下一值是当前值加一，状态机的下一状态由当前状态和输入共同译码。若状态元件是透明锁存器，LE 打开的整个窗口里，新状态会穿过组合逻辑再次改变 D，D 再穿透到 Q，一个窗口内状态可能连续前进好几次，前进步数由窗口宽度和路径延迟共同决定——这等于把时序问题换成了竞态问题。触发器把这个回路切成离散步：一个沿只提交一次，本沿采到的 D 要等下一个沿之后才可能再次影响 D，于是每台状态机每个周期恰好前进一步。这就是同步设计可分析、可调试的根本原因。

这并不是说透明锁存器一无是处。前面 8086 地址锁存的案例已经说明，协议给出清晰电平窗口的地方正是锁存器的用武之地；下一节还会看到，主从结构正是用两个窗口错开的锁存器拼出一个触发器。工程结论因此是分层而不是否决：跨相位保持用锁存器，状态回路的提交边界用触发器。

\section{二、触发器、时钟周期与采样窗口}

触发器解决的是“什么时候写入”的问题。它不像锁存器那样在一段时间内透明，而是在时钟沿附近采样输入，把采样值保存到输出。

\begin{lstlisting}
on clock edge:
  Q = D
between clock edges:
  Q holds
\end{lstlisting}

这让大电路有了清楚节奏。一个时钟沿之后，上一批触发器输出新状态；这些状态经过组合逻辑传播，形成下一批输入；到下一个时钟沿，目标触发器再统一采样。于是“计算”和“保存”被分开了。

以 74HC74 这类双 D 触发器为例，D 是待采样数据，CP 或 CLK 是采样边界，Q 是保存后的输出，异步置位和异步清零用于把状态强制带到已知值。它的引脚命名迫使设计者区分三件事：数据是否已经准备好，时钟边沿是否合格，异步控制脚是否处于安全状态。若 D 在边沿附近变化，触发器不保证采样结果；若清零脚悬空，触发器可能在电源噪声中被误复位；若时钟来自未消抖按钮，单次按压可能被解释成多个边沿。

\begin{chipcase}
把 74HC74 接成最小实验电路时，应把异步置位和异步清零接到确定无效电平或受控复位，D 端由开关、组合逻辑或前级触发器驱动，CP 端只能接经过整形的时钟。若用按钮作为单步时钟，按钮必须先经过上拉/下拉、消抖和施密特整形；否则一次机械按压可能产生多个有效边沿，使 Q 连续翻转或连续采样。这个案例说明，\hwtiming{时钟边沿质量}与逻辑功能同等重要。
\end{chipcase}

把这只触发器拆开看，每个对象都有明确的读法。D 输入是下一个边界要保存的候选值，但“只要最终正确、何时稳定都无所谓”是误解——它必须在采样窗口前后都稳定。CLK/CP 是状态提交的边界，不是任意跳变都能当时钟，边沿质量和电平窗口都有要求。Q 输出是已提交的当前状态，它和 D 只在边界之后才相等，不能认为随时相等。异步复位/置位不等时钟即可强制状态，但不能悬空或随意接按钮。数据手册时序（建立、保持、脉宽、传播延迟）不是高频设计才需要读的东西，低速实验同样会因为不满足它而失败。
\vspace{1.0em plus 0.35em}
\begin{block}[x=2.0cm,y=1.6cm]
  \node[draw, rectangle, minimum width=2.2cm, minimum height=1.8cm, fill=BookCodeBg, rounded corners=3pt] (ff) at (0,0) {};
  \node[above=0.12cm of ff, font=\small\color{BookNavy}] {D-FF};
  \node[left] at ($(ff.west)+(-0.12,0.5)$) {D};
  \node[left] at ($(ff.west)+(-0.12,-0.5)$) {CLK};
  \node[right] at ($(ff.east)+(0.12,0.5)$) {Q};
  \node[right] at ($(ff.east)+(0.12,-0.5)$) {Q'};
  \draw[BookAccent, thick] ($(ff.west)+(0.03,-0.61)$) -- ($(ff.west)+(0.03,-0.39)$) -- ($(ff.west)+(0.25,-0.5)$) -- cycle;
  \node[below=0.3cm of ff, font=\footnotesize, BookAccent] {上升沿采样};
\end{block}
\diagramnote{D 触发器方框符号：只在时钟上升沿采样 D 并保存到 Q。}

触发器为什么能做到“不像锁存器那样透明”？最常见的实现是主从结构：两个 D 锁存器串联，主锁存器的使能接反相后的时钟，从锁存器的使能接原时钟。CLK=0 时，主锁存器透明、从锁存器保持——D 的变化可以进入主，但到不了 Q；CLK 从 0 变 1 的瞬间，主锁存器关闭、从锁存器打开，主此刻保存的值被交给 Q；CLK=1 时，主保持、从透明——Q 跟随主保存的值，但主不再受 D 影响。所以 Q 只在上升沿那一刻获得新值。边沿触发不是一种新的存储元件，而是两个透明窗口错开半个周期的锁存器组合。

\begin{block}[x=1cm,y=1cm]
  \node[hw state, minimum width=2.7cm, minimum height=1.2cm] (m) at (0,0) {主锁存器\\LE = $\overline{CLK}$};
  \node[hw state, minimum width=2.7cm, minimum height=1.2cm] (s) at (4.4,0) {从锁存器\\LE = CLK};
  \node[not gate US, draw=BookAccent, line width=0.5pt, scale=1.0] (inv) at (0.9,-2.1) {};
  \node[hw label] at (-1.8,0) {D};
  \draw[hw bus] (-1.5,0) -- (m.west);
  \draw[hw bus] (m.east) -- node[hw label, above]{M} (s.west);
  \draw[hw bus] (s.east) -- node[hw label, above]{Q} ++(1.5,0);
  \node[hw label] at (-1.5,-2.1) {CLK};
  \draw[hw bus] (-1.2,-2.1) -- (inv.input);
  \draw[fill=BookInk] (0.0,-2.1) circle (0.05);
  \draw[hw bus] (0.0,-2.1) -- (0.0,-2.7) -- (4.4,-2.7) -- (s.south);
  \draw[hw bus] (inv.output) -- ++(0.5,0) |- (m.south);
  \node[hw label] at (2.2,-3.2) {主锁存器在 CLK=0 时透明，从锁存器在 CLK=1 时透明：两个窗口永远错开};
\end{block}
\diagramnote{主从 D 触发器：D 只能在 CLK=0 期间进入主锁存器；CLK 上升沿把主保存的值移交从锁存器，Q 才更新。}

把主从两级放进同一张波形，"边沿"就不再是断言：

\begin{waveform}[x=1cm,y=1cm]
  \draw[hw bus] (0,-0.4) -- (11.4,-0.4);
  \node[hw label] at (11.4,-0.75) {时间};
  \draw[BookAccent, line width=1.0pt] (0,4.4) -- (2,4.4) -- (2,5.2) -- (4,5.2) -- (4,4.4) -- (6,4.4) -- (6,5.2) -- (8,5.2) -- (8,4.4) -- (10,4.4) -- (10,5.2) -- (11,5.2);
  \node[hw label, anchor=east] at (-0.2,4.8) {CLK};
  \draw[BookAmber, line width=1.0pt] (0,3.0) -- (0.8,3.0) -- (0.8,3.8) -- (2.8,3.8) -- (2.8,3.0) -- (4.6,3.0) -- (4.6,3.8) -- (6.8,3.8) -- (6.8,3.0) -- (11,3.0);
  \node[hw label, anchor=east] at (-0.2,3.4) {D};
  \draw[BookTeal, line width=1.0pt] (0,1.6) -- (0.8,1.6) -- (0.8,2.4) -- (4.0,2.4) -- (4.0,1.6) -- (4.6,1.6) -- (4.6,2.4) -- (8.0,2.4) -- (8.0,1.6) -- (11,1.6);
  \node[hw label, anchor=east] at (-0.2,2.0) {M（主）};
  \draw[BookRed, line width=1.1pt] (0,0) -- (2.0,0) -- (2.0,0.8) -- (10.0,0.8) -- (10.0,0) -- (11,0);
  \node[hw label, anchor=east] at (-0.2,0.4) {Q（从）};
  \draw[BookMuted, dashed, line width=0.4pt] (2,-0.4) -- (2,5.6);
  \draw[BookMuted, dashed, line width=0.4pt] (6,-0.4) -- (6,5.6);
  \node[hw label] at (2,5.85) {上升沿：主关闭、从打开};
  \node[hw label, text=BookAmber] at (3.6,2.78) {CLK=1：主保持，D 变也进不来};
  \node[hw label, text=BookTeal] at (1.0,1.12) {CLK=0：M 跟随 D};
\end{waveform}
\diagramnote{主从两级的工作波形。M 在 CLK=0 期间跟随 D，CLK=1 期间保持——所以 D 在 CLK=1 里的翻转（本例 2.8 处）进不了主；Q 只在上升沿复制 M 当时的值。所谓"上升沿采样"，全部秘密就在这两条错开的透明窗口里。}

\topic{专题：JK 与 T 触发器——从两条命令到一个翻转请求}

D 触发器不是唯一的边沿触发器。中小规模器件里长期流行的还有 JK 触发器：它把 SR 的两条命令搬到时钟沿上，J 对应置位请求，K 对应复位请求，并且顺手解决了 SR 的禁用组合问题——J 和 K 同时有效不再是非法输入，而是表示“翻转”。逐行看它的约定：J=1, K=0 时，下一个有效沿把 Q 写成 1，与当前值无关；J=0, K=1 时写成 0；J=K=0 时没有请求，Q 保持；J=K=1 时，下一值被定义成当前值的反，于是每个有效沿 Q 翻转一次。

\begin{longtable}{L{0.08\linewidth}L{0.08\linewidth}L{0.18\linewidth}L{0.54\linewidth}}
\toprule
J & K & 下一 Q & 含义 \\
\midrule
0 & 0 & Q & 无请求：保持旧值 \\
1 & 0 & 1 & 置位请求：沿后 Q=1 \\
0 & 1 & 0 & 复位请求：沿后 Q=0 \\
1 & 1 & $\overline{Q}$ & 翻转请求：沿后 Q 取反 \\
\bottomrule
\end{longtable}

对照 SR 锁存器就能看出 JK 的价值：SR 的第四行是禁用组合，JK 的第四行是有明确定义的功能。同一个“两条命令同时有效”的情形，前者只能交给门延迟竞争，后者被约定成确定行为。这也说明触发器家族的差异不在存储机制——它们都在边沿保存 1 bit——而在输入命令集如何约定。

把 J 和 K 短接在一起，JK 触发器就退化成 T 触发器：T=0 相当于 J=K=0，保持；T=1 相当于 J=K=1，每个有效沿翻转一次。T 触发器几乎不单独出现，但它是计数与分频的基本单元。把 T 固定接 1，每个时钟沿 Q 翻转一次：输入每过两个完整周期，Q 才完成一次高低循环，所以 Q 的频率恰好是输入的一半——一只 T=1 的触发器就是一个二分频器。

把 JK 的四行命令画成两个状态之间的转移，四种输入组合各就各位：

\begin{pdfdiagram}
  % JK 触发器：左 Q=0，右 Q=1
  \node[hw label] at (1.5,1.95) {JK 触发器};
  \node[hw state, minimum width=1.1cm, minimum height=0.8cm] (z) at (0,0) {Q=0};
  \node[hw state, minimum width=1.1cm, minimum height=0.8cm] (o) at (3.0,0) {Q=1};
  \draw[hw bus] (z.north) -- ($(z.north)+(0,0.35)$) -- ($(o.north)+(0,0.35)$) -- (o.north);
  \node[hw label, anchor=south] at (1.5,0.82) {J=1, K=0：置位\\J=1, K=1：翻转};
  \draw[hw bus] (o.south) -- ($(o.south)+(0,-0.35)$) -- ($(z.south)+(0,-0.35)$) -- (z.south);
  \node[hw label, anchor=north] at (1.5,-0.82) {J=0, K=1：复位\\J=1, K=1：翻转};
  \draw[hw bus] (z.west) -- (-0.85,0) -- (-0.85,0.75) -- (-0.22,0.75) -- ($(z.north)+(-0.22,0)$);
  \draw[hw bus] (o.east) -- (3.85,0) -- (3.85,0.75) -- (3.22,0.75) -- ($(o.north)+(0.22,0)$);
  % T 触发器：J、K 短接的退化
  \node[hw label] at (8.4,1.95) {T 触发器（J、K 短接）};
  \node[hw state, minimum width=1.1cm, minimum height=0.8cm] (tz) at (6.9,0) {Q=0};
  \node[hw state, minimum width=1.1cm, minimum height=0.8cm] (to) at (9.9,0) {Q=1};
  \draw[hw bus] (tz.north) -- ($(tz.north)+(0,0.35)$) -- ($(to.north)+(0,0.35)$) -- (to.north);
  \node[hw label, anchor=south] at (8.4,0.82) {T=1：翻转};
  \draw[hw bus] (to.south) -- ($(to.south)+(0,-0.35)$) -- ($(tz.south)+(0,-0.35)$) -- (tz.south);
  \node[hw label, anchor=north] at (8.4,-0.82) {T=1：翻转};
  \draw[hw bus] (tz.west) -- (6.05,0) -- (6.05,0.75) -- (6.68,0.75) -- ($(tz.north)+(-0.22,0)$);
  \draw[hw bus] (to.east) -- (10.75,0) -- (10.75,0.75) -- (10.12,0.75) -- ($(to.north)+(0.22,0)$);
  \node[hw label] at (1.5,-1.75) {自环：J=K=0 时保持旧值};
  \node[hw label] at (8.4,-1.75) {自环：T=0 时保持旧值};
\end{pdfdiagram}
\diagramnote{JK 触发器（左）与 T 触发器（右）的状态转移。JK 的两条跨边覆盖三种组合：J=1, K=0 把 Q 写 1，J=0, K=1 把 Q 写 0，J=K=1 在两条跨边上都出现——每个有效沿必换边，这就是翻转环；J=K=0 时两个状态都走自环保留旧值。T 触发器是 J、K 短接的退化：T=0 保持，T=1 每沿翻转一次。}

把三级这样的触发器串起来，前一级的 Q 作为后一级的时钟，就得到一条纹波分频链：

\begin{longtable}{L{0.16\linewidth}L{0.2\linewidth}L{0.22\linewidth}L{0.32\linewidth}}
\toprule
级 & 时钟来源 & 输出翻转频率 & 相对输入时钟 \\
\midrule
第 1 级 Q$_0$ & 输入时钟 CLK & $f/2$ & 二分频 \\
第 2 级 Q$_1$ & Q$_0$ & $f/4$ & 四分频 \\
第 3 级 Q$_2$ & Q$_1$ & $f/8$ & 八分频 \\
\bottomrule
\end{longtable}

代入数字：输入时钟 1MHz 时，Q$_0$ 为 500kHz，Q$_1$ 为 250kHz，Q$_2$ 为 125kHz；Q$_2$ 的周期是 8μs，恰好等于 8 个输入周期。换一个角度看同一条链：每经过 8 个输入沿，三级输出（Q$_2$,Q$_1$,Q$_0$）的组合就遍历 8 种编码回到起点，所以三级 T 链同时就是一个模 8 的 3-bit 计数器。分频和计数是同一个电路的两种读法：前者看单路输出的频率，后者看多路输出的编码。本章后面讲计数器时还会看到，纹波结构的代价是各级时钟到输出延迟逐级累加，Q$_2$ 要等三级延迟之后才稳定；位数一多，这条链本身就成为关键路径，这正是同步计数器存在的理由。

有了触发器的采样模型，建立/保持时间窗口就可以对照波形图逐项检查：

\begin{waveform}[x=2pt,y=6ex]
  % CLK — 从低电平开始, 上升沿在 32/96/160/224pt
  \draw[BookAccent,line width=3pt]
    (0,20ex) -- ++(32pt,0) -- ++(0,6ex) -- ++(32pt,0) -- ++(0,-6ex)
             -- ++(32pt,0) -- ++(0,6ex) -- ++(32pt,0) -- ++(0,-6ex)
             -- ++(32pt,0) -- ++(0,6ex) -- ++(32pt,0) -- ++(0,-6ex)
             -- ++(32pt,0) -- ++(0,6ex) -- ++(32pt,0) -- ++(0,-6ex)
             -- ++(20pt,0);
  \node[BookAccent,left] at (0,23ex) {CLK};
  % 建立/保持窗口 — 上升沿前 12pt 建立(红), 上升沿后 12pt 保持(琥珀)
  \foreach \x in {20,84,148,212} {\fill[BookRed!25] (\x pt, 19.5ex) rectangle (\x+12 pt, 26.5ex);}
  \foreach \x in {32,96,160,224} {\fill[BookAmber!25] (\x pt, 19.5ex) rectangle (\x+12 pt, 26.5ex);}
  \node[BookRed,right,font=\small] at (20pt,28.5ex) {$t_{su}$};
  \node[BookAmber,right,font=\small] at (34pt,28.5ex) {$t_{h}$};
  % D — 高低变化都远离上升沿窗口
  \draw[BookAmber,line width=3pt]
    (0,5ex) -- ++(8pt,0) -- ++(0,8ex) -- ++(60pt,0) -- ++(0,-8ex) -- ++(56pt,0)
            -- ++(0,8ex) -- ++(56pt,0) -- ++(0,-8ex) -- ++(80pt,0);
  \node[BookAmber,left] at (0,9ex) {D};
  % Q — 上升沿后 tCQ=8pt 更新
  \draw[BookTeal,line width=3pt]
    (0,-9ex) -- ++(40pt,0) -- ++(0,7ex) -- ++(64pt,0) -- ++(0,-7ex) -- ++(64pt,0)
             -- ++(0,7ex) -- ++(64pt,0) -- ++(0,-7ex) -- ++(28pt,0);
  \node[BookTeal,left] at (0,-5.5ex) {Q};
\end{waveform}
\diagramnote{建立/保持时间窗口都相对于上升采样沿：D 必须在上升沿之前 $t_{su}$ 已经稳定，并在上升沿之后 $t_{h}$ 继续保持；Q 在上升沿后经过 $t_{CQ}$ 才更新。}
\vspace{0.8em plus 0.25em}

最常见的同步路径可以写成：

\begin{lstlisting}
source register -> combinational logic -> destination register
\end{lstlisting}
\vspace{1.0em plus 0.35em}

时钟沿到来后，源寄存器输出旧状态对应的新值。这个值进入组合逻辑，经过门延迟和布线延迟，最终到达目标寄存器输入。下一个时钟沿到来时，目标寄存器采样这个结果。因此频率不是越高越好。频率提高意味着周期缩短，留给组合逻辑稳定的时间变少。若周期短到最慢路径来不及稳定，目标寄存器就可能采到旧值、过渡值或错误值。

同步设计的关键，是区分\hwkey{当前状态}和\hwkey{下一状态}。当前状态存在于源寄存器 Q 端，下一状态存在于目标寄存器 D 端。组合逻辑可以在整个周期内反复变化，直到输入稳定、路径传播完成、D 端满足建立时间；但只有时钟沿到来，D 才会成为新的 Q。用软件赋值类比硬件时，最容易犯的错误是把 D 端的临时变化当成已经写入。硬件没有“变量立即更新”的全局瞬间，只有被状态元件定义出来的提交边界。

阅读触发器数据手册时，至少要把四类时间参数分开。第一，时钟到输出延迟说明源状态何时对后级可见。第二，建立时间和保持时间说明目标触发器允许 D 端怎样靠近采样沿。第三，时钟高电平、低电平和脉冲宽度说明怎样的波形才算合格时钟。第四，异步复位或置位的脉宽、恢复时间和释放时间说明复位不能随意撤销。即使低速面包板实验看起来“肉眼很慢”，按钮抖动、慢边沿和悬空输入仍然可能违反这些时间约定。

\begin{longtable}{L{0.22\linewidth}L{0.38\linewidth}L{0.3\linewidth}}
\toprule
数据手册项目 & 应回答的问题 & 板上故障表现 \\
\midrule
时钟到输出延迟 & Q 何时能作为后级输入使用 & 后级过早采样旧值或过渡值 \\
建立/保持时间 & D 在边沿前后要稳定多久 & 偶发采错、亚稳态或重复触发 \\
最小时钟脉宽 & CP 高低电平是否足够长 & 窄脉冲被忽略或产生异常采样 \\
复位恢复时间 & 复位释放离时钟沿是否过近 & 复位后状态机随机进入非入口状态 \\
\bottomrule
\end{longtable}
\vspace{1.0em plus 0.35em}

触发器采样也不是数学瞬间。输入 D 必须在时钟沿到来前稳定一小段时间，这称为建立时间，违反它会采到错误值或不稳定值；时钟沿之后还要继续稳定一小段时间，这称为保持时间，违反它会让新旧值混入采样窗口。
\vspace{1.0em plus 0.35em}
\begin{block}[x=1.6cm,y=1.6cm]
  \node[hw compute, minimum width=2.0cm] (tcq) {tCQ};
  \node[hw compute, right=0.45cm of tcq, minimum width=3.0cm] (tcomb) {组合逻辑};
  \node[hw compute, right=0.45cm of tcomb, minimum width=1.8cm] (twire) {连线};
  \node[hw compute, right=0.45cm of twire, minimum width=1.6cm] (tsu) {tSU};
  \node[hw compute, right=0.45cm of tsu, minimum width=1.6cm] (tskew) {偏斜+余量};
  \node[hw label, below=0.3cm of tcq] {0.8ns};
  \node[hw label, below=0.3cm of tcomb] {3.6ns};
  \node[hw label, below=0.3cm of twire] {0.7ns};
  \node[hw label, below=0.3cm of tsu] {0.4ns};
  \node[hw label, below=0.3cm of tskew] {0.5ns};
  \draw[hw bus] (tcq.west) -- ++(-0.6,0) node[left,hw label]{CLK};
  \draw[hw bus] (tskew.east) -- ++(0.6,0) node[right,hw label]{CLK};
  \draw[BookMuted, decorate, decoration={brace,amplitude=5pt,raise=2pt}] (tcq.north west) -- (tskew.north east) node[midway,above=12pt,BookAccent,font=\small] {Tclk $\geq$ 6.0ns};
\end{block}
\diagramnote{时钟周期预算：各段时间之和决定最小周期。}
\vspace{0.8em plus 0.25em}

一条路径的周期预算可以写成：

\begin{lstlisting}
clock period >= source clock-to-Q
              + combinational delay
              + wire delay
              + destination setup time
\end{lstlisting}

更严谨地写，还要把时钟偏斜和余量纳入公式。若目标触发器的时钟沿比源触发器更早到达，有效可用时间会变少；若目标时钟更晚到达，建立时间压力可能减轻，但保持时间压力可能增加。为了保持第一轮理解清晰，可以先把建立约束写成：

\begin{lstlisting}
Tclk >= tCQ_max + tcomb_max + twire_max + tsetup + tskew + margin
\end{lstlisting}

其中 \code{tCQ_max} 是源触发器时钟到输出的最大时间，\code{tcomb_max} 是组合逻辑最大延迟，\code{twire_max} 是布线最大延迟，\code{tsetup} 是目标触发器建立时间，\code{tskew} 表示不利时钟偏斜，\code{margin} 是留给电压、温度、工艺、建模误差和老化的余量。建立约束回答“最慢结果能否赶上下一个采样沿”。

\begin{longtable}{L{0.2\linewidth}L{0.2\linewidth}L{0.5\linewidth}}
\toprule
项目 & 示例数值 & 读法 \\
\midrule
\code{tCQ_max} & 1.2ns & 源寄存器输出最晚 1.2ns 后有效 \\
\code{tcomb_max} & 5.8ns & 组合逻辑最坏传播时间 \\
\code{twire_max} & 0.9ns & 远距离连线和负载引入延迟 \\
\code{tsetup} & 0.6ns & 目标寄存器采样前要求 \\
\code{tskew+margin} & 0.5ns & 偏斜和工程余量 \\
合计 & 9.0ns & 时钟周期必须不小于 9.0ns \\
\bottomrule
\end{longtable}
\vspace{1.0em plus 0.35em}
\begin{block}[x=2.0cm,y=1.6cm]
  \node[draw, dashed, rounded corners, inner sep=8pt, label={[hw label]above:异步输入}] (rst_in) {外部复位};
  \node[hw state, right=1.2cm of rst_in, label={[hw label]above:第一级}] (ff1) {FF$_1$};
  \node[hw state, right=1.2cm of ff1, label={[hw label]above:第二级}] (ff2) {FF$_2$};
  \node[hw control, right=1.2cm of ff2, label={[hw label]above:同步输出}] (ctrl) {同步复位};
  \draw[hw bus] (rst_in) -- (ff1) node[below,hw label,pos=0.5]{异步拉入};
  \draw[hw bus] (ff1) -- (ff2) node[below,hw label,pos=0.5]{同步释放};
  \draw[hw bus] (ff2) -- (ctrl) node[below,hw label,pos=0.5]{边界对齐};
  \node[below=0.3cm of ff1, hw label] {D=1};
  \node[below=0.3cm of ff2, hw label] {D=Q$_1$};
\end{block}
\diagramnote{异步拉入、同步释放的复位电路。}
\vspace{0.8em plus 0.25em}

若目标频率要求周期为 8ns，上表路径就不合格。修复方式不是在公式里删除余量，而是缩短组合逻辑、增加寄存器边界、改善布局布线、换用更快单元、降低频率，或重新安排微操作。同步设计的严肃性在于：功能真值表正确，仅能说明最终会算对；时序约束也满足，才说明能在规定边界前算对。

保持时间看起来更小，却同样危险。若源寄存器输出太快、路径太短，目标寄存器可能在同一个时钟沿之后立刻看到新值，从而破坏它本应采旧值的窗口。长路径会限制最高频率，短路径也可能破坏保持约束。

保持约束可以按最小延迟理解：

\begin{lstlisting}
tCQ_min + tcomb_min + twire_min >= thold + unfavorable_skew
\end{lstlisting}

这条式子回答“新值是否来得太早”。例如源触发器最快 0.1ns 就改变 Q，路径中几乎没有组合逻辑，布线也很短，而目标触发器要求保持 0.3ns；若时钟偏斜又让目标采样沿相对更晚，那么新值可能在目标仍需要旧值稳定时抵达。修复保持问题的常见手段是插入小延迟、调整布线、改变时钟树或重新放置寄存器。它和建立问题方向相反：建立嫌路径太慢，保持嫌路径太快。

把这几类约束放在一起看，修复手段完全不同：建立约束失败说明最慢路径赶不上采样沿，要缩短组合逻辑、流水线化、降频或优化布局；保持约束失败说明最快路径太早冲进采样窗口，要增加最小延迟、调整时钟偏斜或改变布线；复位释放失败是状态元件退出复位不一致，要异步拉入、同步释放并限定复位域；跨时钟输入失败是采样沿附近输入不受控，要用两级同步、握手、异步 FIFO 或协议化采样。用建立的修法去修保持，或者用降频去解决复位和跨边界问题，都会把故障掩盖成偶发错误。
\vspace{1.0em plus 0.35em}

\begin{waveform}[x=2pt,y=6ex]
  % CLK — 从低电平开始, 上升沿在 32/96/160/224pt
  \draw[BookAccent,line width=3pt]
    (0,20ex) -- ++(32pt,0) -- ++(0,6ex) -- ++(32pt,0) -- ++(0,-6ex)
             -- ++(32pt,0) -- ++(0,6ex) -- ++(32pt,0) -- ++(0,-6ex)
             -- ++(32pt,0) -- ++(0,6ex) -- ++(32pt,0) -- ++(0,-6ex)
             -- ++(32pt,0) -- ++(0,6ex) -- ++(32pt,0) -- ++(0,-6ex)
             -- ++(20pt,0);
  \node[BookAccent,left] at (0,23ex) {CLK};
  % D — 第一个上升沿 tCQ 后变化, 第二个上升沿 tSU 前已稳定
  \draw[BookAmber,line width=3pt]
    (0,5ex) -- ++(44pt,0) -- ++(0,8ex) -- ++(128pt,0) -- ++(0,-8ex) -- ++(88pt,0);
  \node[BookAmber,left] at (0,9ex) {D};
  % Q — 第二个上升沿 tCQ 后才变化
  \draw[BookTeal,line width=3pt]
    (0,-9ex) -- ++(108pt,0) -- ++(0,7ex) -- ++(128pt,0) -- ++(0,-7ex) -- ++(24pt,0);
  \node[BookTeal,left] at (0,-5.5ex) {Q};
  % 测量标注
  \draw[BookRed,line width=1.2pt,<->]
    (32pt,17ex) -- node[above=1pt,font=\scriptsize\bfseries,BookRed]{$t_{CQ}$} (44pt,17ex);
  \draw[BookAmber,line width=1.2pt,<->]
    (44pt,17ex) -- node[above=1pt,font=\scriptsize\bfseries,BookAmber]{组合逻辑延迟} (84pt,17ex);
  \draw[BookTeal,line width=1.2pt,<->]
    (84pt,17ex) -- node[above=1pt,font=\scriptsize\bfseries,BookTeal]{$t_{SU}$} (96pt,17ex);
\end{waveform}
\diagramnote{同一个上升沿既是源寄存器的更新沿，也是目标寄存器的采样沿：源在 $t_{CQ}$ 后改变 D，组合逻辑在下一上升沿 $t_{SU}$ 前让 D 稳定，目标在该上升沿后 $t_{CQ}$ 更新 Q。$t_{CQ}$ + 组合逻辑延迟 + $t_{SU}$ = 最小时钟周期。}



\section{三、关键路径、复位与跨边界输入}

一块同步电路里会有很多寄存器到寄存器路径。最长、最慢、最限制周期的那条叫关键路径。它可能是一串加法器，也可能是多级选择器，也可能是远距离布线加上复杂译码。

例如一个 32-bit 串行进位加法器接在两个寄存器之间。最低位进位要一位一位传到最高位，最高位结果最后才稳定。如果把这个加法器放在一个周期里完成，那么周期至少要等最坏进位链走完。若把运算拆成更短的几段，中间插入寄存器，单个周期可以变短，但完成一次完整运算需要更多周期。

关键路径不只来自“算术看起来复杂”的地方。一个很宽的 mux、远距离控制线、解码后再进入选择器的路径、寄存器堆读口到 ALU 再到写回选择器的路径，都可能成为限制频率的对象。真实 CPU 中，很多优化不是改变逻辑功能，而是重新安排边界：把一次长动作拆成取指、译码、执行、访存、写回；把全局信号改成本地译码；把长总线切成更短的局部通路。时序分析让这些选择可以相互比较，而不是靠直觉判断“这条线应该很快”。

在整机视角下，8086、80386 和现代处理器都在反复处理同一个问题：哪些动作放在一个边界之间完成，哪些动作必须切成多个边界。8086 的外部总线周期把地址、控制和数据分成若干相位，用总线相位和等待状态分步完成长访问；80386 的 32-bit 数据通路、分页地址转换和总线接口让路径更宽、更复杂，因此必须更严格地区分内部状态、外部总线状态和等待状态；现代 CPU 则进一步把取指、译码、重命名、执行、访存和提交拆成流水阶段，用局部化、cache 层次和旁路网络缩短常见路径。FPGA/ASIC 数据通路里，宽 mux、长连线、加法器链和译码扇出也按同样思路处理：插入寄存器、重定时、局部译码和布局优化。它们规模不同，但工程判断相同：若一条路径太长，就要缩短组合逻辑、增加状态边界或改变协议等待方式。

这也解释了"一周期完成"和"拆段插入寄存器"的取舍：一周期完成长运算，控制简单、单条动作需要的边界少，代价是关键路径长、频率低；拆段插入寄存器让单周期路径变短、频率可能更高，代价是完成一次动作需要更多边界和更复杂的控制。
\vspace{1.0em plus 0.35em}
\begin{block}[x=2.0cm,y=1.6cm]
  \node[draw, dashed, rounded corners, inner sep=8pt, minimum width=1.6cm, label={[hw label]above:异步输入域}] (async) {外部事件};
  \node[hw state, right=1.6cm of async, label={[hw label]above:第一级同步}] (ff1) {FF$_1$};
  \node[hw state, right=1.4cm of ff1, label={[hw label]above:第二级同步}] (ff2) {FF$_2$};
  \node[draw, dashed, rounded corners, inner sep=8pt, right=1.6cm of ff2, minimum width=1.8cm, label={[hw label]above:目标时钟域}] (sync) {同步后信号};
  \draw[hw bus] (async) -- (ff1) node[below,hw label,pos=0.5]{可能亚稳态};
  \draw[hw bus] (ff1) -- (ff2) node[below,hw label,pos=0.5]{恢复时间};
  \draw[hw bus] (ff2) -- (sync) node[below,hw label,pos=0.5]{稳定输出};
  \node[below=0.3cm of ff1, BookRed, font=\small] {meta};
  \node[below=0.3cm of ff2, BookTeal, font=\small] {sync};
\end{block}
\diagramnote{两级触发器同步器结构。}
\vspace{0.8em plus 0.25em}

触发器通电后的值不一定符合设计预期。若一组状态寄存器上电后取值各不相同，后面的组合逻辑可能立刻进入非法状态。复位信号的作用，是把关键状态单元带到已知点，例如计数器清零、状态机回到 idle、控制寄存器进入安全默认值。

并不是所有触发器都必须复位。控制状态、外部使能、写存储器信号、片选、错误标志和协议状态通常需要复位，因为它们会决定机器是否对外产生动作；纯数据缓冲如果有可靠写使能保护，未必需要每一位都在复位时清零。过度复位会增加布线、负载和时序压力，尤其在大规模芯片中会让复位网络本身变成困难对象。因此，复位设计不是“全清零”，而是列出哪些状态必须已知、哪些输出必须安全、哪些数据可以等第一次有效写入后再使用。

按这个原则分类：控制状态机必须复位，它要有唯一入口和安全输出；外部写使能和片选必须复位，上电和复位期间不能误写外设或存储器；计数器和超时器通常要复位，初始计数会影响协议等待和超时判断；流水数据寄存器视情况而定，若有 valid 位保护，数据位可以等首次有效写入；同步器和事件 pending 位必须复位，复位释放不能制造伪事件。
\vspace{1.0em plus 0.35em}
复位本身也要讲边界。

若复位释放时，有的触发器先退出复位，有的后退出复位，状态机可能看到一半新状态、一半旧状态。常见处理是异步拉入复位以便快速进入安全状态，再同步释放，让所有相关触发器在清楚的时钟边界恢复工作。复位不能简化为“清零所有东西”。哪些状态需要复位，哪些数据寄存器可以不复位，哪些输出在复位期间必须安全，都要由硬件需求决定。

\topic{专题：异步复位与同步复位——立即生效与沿上生效的对照}

触发器的复位有两条实现路径。异步复位使用触发器专用的 CLR 引脚，直接强制内部反馈环进入 0 状态，不问时钟：引脚一有效，Q 立刻被拉成 0，时钟是否在运行都无所谓。同步复位只是数据路径上的一个约定：复位有效时，D 端之前的选择逻辑把 0 送到触发器输入，Q 要等下一个有效沿才变成 0。前者改变存储元件本身的状态，后者改变准备提交给存储元件的下一值——这是两种复位最本质的区别。

两条路径各有代价。异步复位的优势是上电即有效：时钟还没起振、时钟被门控关闭、系统卡在未知状态时，只有不问时钟的复位才能保证把状态拉到已知点。它的代价在释放端：撤销复位同样是一个异步事件，若释放时刻落在时钟沿附近的恢复时间窗口内，触发器可能这一拍就响应新的 D，也可能仍停留在复位值，结果不可预知。同步复位的优势恰好相反：生效和释放都以时钟沿为界，普通的建立/保持分析就能覆盖，不会制造半个周期的歧义；它的代价是必须有时钟才起作用，并且复位选择逻辑串在 D 路径上，给关键路径增加一级延迟。

把两条路径画成结构对照，“改存储元件本身”与“改下一值”的区别就写在连线上：

\begin{pdfdiagram}
  % 左：异步复位，CLR 专用引脚
  \node[hw label] at (0,1.2) {异步复位：CLR 引脚直接强制};
  \node[hw state, minimum width=1.9cm, minimum height=1.2cm] (ff1) at (0,0) {D 触发器\\（含 CLR）};
  \draw[hw bus] (-1.8,0.25) -- ($(ff1.west)+(0,0.25)$);
  \node[hw label] at (-1.45,0.48) {D};
  \draw[hw bus] (-1.8,-0.25) -- ($(ff1.west)+(0,-0.25)$);
  \node[hw label] at (-1.45,-0.48) {CLK};
  \draw[hw bus] (ff1.east) -- (1.3,0);
  \node[hw label] at (0.9,0.26) {Q};
  \node[hw label] (rst1) at (0,-1.75) {外部复位 rst};
  \draw[hw return] (rst1.north) -- (ff1.south);
  % 右：同步复位，D 路径上的 mux
  \node[hw label] at (6.6,1.2) {同步复位：0 经 mux 等沿写入};
  \node[hw compute, minimum width=1.6cm, minimum height=1.2cm] (mux) at (5.2,0) {mux\\2 选 1};
  \node[hw state, minimum width=1.9cm, minimum height=1.2cm] (ff2) at (8.2,0) {D 触发器};
  \draw[hw bus] (3.0,0.3) -- ($(mux.west)+(0,0.3)$);
  \node[hw label] at (3.45,0.53) {D 数据};
  \draw[hw bus] (3.0,-0.3) -- ($(mux.west)+(0,-0.3)$);
  \node[hw label] at (3.45,-0.53) {常量 0};
  \node[hw label] (rst2) at (5.2,-1.75) {rst 作选择端};
  \draw[hw return] (rst2.north) -- (mux.south);
  \draw[hw bus] (mux.east) -- (ff2.west);
  \node[hw label] at (6.62,0.26) {D};
  \node[hw label] (clk2) at (8.2,-1.75) {CLK};
  \draw[hw bus] (clk2.north) -- (ff2.south);
  \draw[hw bus] (ff2.east) -- (9.8,0);
  \node[hw label] at (9.45,0.26) {Q};
\end{pdfdiagram}
\diagramnote{异步复位（左）：rst 走触发器专用 CLR 引脚（虚线控制），不问时钟，引脚一有效 Q 立即清零，时钟停振时也成立。同步复位（右）：rst 只是 D 路径上 mux 的选择端，有效时把常量 0 送到 D，Q 要等下一个有效沿才变 0；生效与释放都由时钟沿界定，代价是选择逻辑串在数据路径上，给关键路径加一级延迟。}

\begin{longtable}{L{0.24\linewidth}L{0.3\linewidth}L{0.34\linewidth}}
\toprule
关键问题 & 异步复位 & 同步复位 \\
\midrule
何时生效 & 引脚有效立即生效，不问时钟 & 下一个有效沿才生效 \\
无时钟时能否工作 & 能：上电、停振时都有效 & 不能：沿不来状态不变 \\
数据路径代价 & 不占 D 路径，使用专用引脚 & 选择逻辑串入 D 路径，增加一级延迟 \\
释放时刻 & 仍是异步事件，可能违反恢复时间 & 由建立/保持分析覆盖，时刻确定 \\
引脚毛刺影响 & 毛刺直接清掉状态 & 只有沿附近的毛刺才可能被采入 \\
\bottomrule
\end{longtable}

用一个数量例子看释放风险。设触发器要求的恢复时间是 1.0ns，异步复位在某个沿之前 0.3ns 才撤销——离沿太近，这个触发器本次沿是否退出复位不可预知。若状态寄存器的两个 bit 因复位走线长度不同，一个看到释放离沿 0.3ns，另一个看到 1.2ns，前者行为不定、后者满足恢复时间正常退出，两个 bit 就可能错开一拍退出复位，状态机出现一个周期的混合状态，可能走进任何编码。这就是“异步拉入、同步释放”成为标准做法的原因：拉入用异步路径，保证任何时刻都能进入安全状态；释放经过本节前面的两级触发器同步到本地时钟边界，让所有状态元件在同一拍恢复工作。两种复位不是互相替代的竞争方案，而是分别负责“进得去”和“出得来”两道工序。

跨时钟或来自外部世界的输入还会带来亚稳态。若输入 D 正好在时钟沿附近变化，触发器内部节点可能短时间落在不稳定区域，既不像明确 0，也不像明确 1。多数情况下，它会在一小段时间后恢复到明确的 0 或 1，但恢复到哪一边、需要多久，不能按普通组合逻辑精确保证。因此来自另一个时钟节奏或外部世界的信号，不能直接拿来控制大面积电路，通常要先经过同步结构，降低亚稳态扩散概率。

最常见的单 bit 同步结构是两级触发器串联。第一级直接采样异步输入，可能进入亚稳态；第二级在下一个时钟沿再采样第一级输出，给第一级更多时间恢复到明确 0 或 1。两级同步不能从数学上消灭亚稳态，但能显著降低亚稳态传播到后级控制逻辑的概率。多 bit 数据不能简单地每一位各接两级同步后直接使用，因为各位可能在不同周期被采到，形成混合值；多 bit 跨边界通常要使用握手、有效位、灰码计数器或异步 FIFO。

工程上常用 \term{平均无故障时间}{MTBF, Mean Time Between Failures} 描述亚稳态风险是否可接受。它不是功能真值表，而是可靠性指标：异步输入变化越频繁、目标时钟越快，触发器遇到采样窗口的机会越多；留给第一级同步器恢复的时间越长，亚稳态传播到第二级之后的概率越低。因此，提高同步级数、降低异步事件频率、给同步器留出完整周期、避免在同步器后面接复杂组合逻辑，都能改善可靠性。严肃设计不会声称“两级同步一定安全”，而是说明在目标频率、输入事件频率和器件参数下，失效率是否满足系统要求。

放一个数量级感受“恢复时间买可靠性”的杠杆有多大。亚稳态未恢复概率随恢复时间按指数 $e^{-t/\tau}$ 下降：恢复窗口每多一个器件时间常数 $\tau$，失败概率约下降 2.7 倍；每多约 $2.3\tau$，失败概率大约下降一个数量级。假设某触发器 $\tau=0.1$ns，第一级留约 2ns（约 $20\tau$）恢复窗口，单次采样的失败概率约 $10^{-10}$ 量级；异步输入平均每秒变化 1000 次，则大约 $10^7$ 秒（近 4 个月）才出现一次可见失败。若只留 0.3ns（约 $3\tau$）——比如同步器后面紧跟一条长组合路径再吃一级触发器——恢复窗口少了约 17 个 $\tau$，失败概率抬升约 $e^{17}$（约 7 个数量级），达到约 $10^{-3}$，同样输入下每天失败数万次。这些数字是示意量级，真实参数必须查器件手册；但结论形态是普适的：\hwkey{可靠性按指数购买，时钟越快、事件越频繁、恢复窗口越短，买到的越少}。

\begin{longtable}{L{0.22\linewidth}L{0.36\linewidth}L{0.32\linewidth}}
\toprule
MTBF 影响因素 & 对可靠性的影响 & 设计含义 \\
\midrule
目标时钟频率 & 采样次数越多，遇到危险窗口机会越多 & 高频域更需要规范同步结构 \\
异步事件频率 & 输入变化越频繁，触发亚稳态机会越多 & 高频事件应使用握手或 FIFO \\
恢复时间 & 两级触发器之间时间越充足，传播概率越低 & 同步器后不应接长组合路径再采样 \\
器件特性 & 触发器恢复参数决定概率下降速度 & 应使用目标工艺或芯片数据手册参数 \\
\bottomrule
\end{longtable}
\vspace{0.8em plus 0.25em}

\begin{longtable}{L{0.22\linewidth}L{0.34\linewidth}L{0.34\linewidth}}
\toprule
跨边界对象 & 推荐处理 & 不当处理的风险 \\
\midrule
单 bit 按钮/中断 & 两级同步后再进入状态机 & 亚稳态扩散或一次事件被识别多次 \\
多 bit 数据总线 & 使用握手、valid/ready 或 FIFO & 各 bit 来自不同采样时刻，形成错误编码 \\
复位释放 & 异步拉入、同步释放 & 状态机部分触发器先运行 \\
脉冲信号 & 转换为电平握手或拉宽后同步 & 窄脉冲被漏采或重复采样 \\
\bottomrule
\end{longtable}
\vspace{0.8em plus 0.25em}

\topic{专题：跨时钟域的两种基本做法——单 bit 同步与多 bit 协议}

上表把跨边界对象分了类，这里把背后的两类做法讲透。分类的依据不是信号数量本身，而是“接收方能否容忍采样时刻的不确定”。单 bit 信号只有一个自由度，采早采晚最多差一个周期，两级同步器把亚稳态概率降到可忽略量级之后，剩下的时机不确定性可以靠协议容忍。多 bit 数据的每一位各有自己的采样时刻，逐位同步可能让同一拍采到不同代的 bit，错的就不只是时机，而是数值本身，所以必须引入协议。

单 bit 做法本章已经展开：两级触发器串联，第一级留给亚稳态恢复，第二级输出干净电平，前面也已经算过 MTBF 的账——恢复窗口从约 2ns 缩到 0.3ns，失败间隔就从近 4 个月缩成每天数万次。这里只补一条数量级推论：因为失败概率随恢复时间按指数 $e^{-t/\tau}$ 下降，MTBF 对余量极端敏感，余量每多一个器件时间常数 $\tau$，失败间隔大约拉长 2.7 倍；每多约 $2.3\tau$，大约拉长 10 倍。在两级之后再添第三级，恢复窗口近似翻倍，MTBF 随之再抬高若干个数量级，从“产品寿命内偶见”推到“远远超出产品寿命”；反过来，同步器后面紧接一条长组合路径再吃一级触发器，等于把恢复窗口削掉大半，MTBF 可能从年级跌回天级。结论很朴素：可靠性按指数购买，买多买少由恢复窗口决定，而恢复窗口几乎不花钱——只是触发器和布线。

多 bit 做法先看为什么不能逐位同步。设源域一个 4-bit 计数值从 0111 变成 1000，四个 bit 同时变化；若每位各接一条两级同步器，各路走线和器件延迟略有差异，接收域某一拍可能采到低三位的新值和最高位的旧值——于是接收方依次看到 0111、0000、1000，中间的 0000 是源从未发送过的幻影值。控制逻辑若拿幻影值做判断，机器就按一个不存在的状态行动。逐级加触发器也解决不了这个问题，因为它消除的是亚稳态，不是各位之间的采样时刻分歧。

把这次传送按接收时钟逐拍画出，幻影值的出现过程就清晰可见：

\begin{waveform}[x=0.8cm,y=0.9cm]
  % 参考线：x=3 源值跳变，x=4/7/10 接收时钟采样沿
  \draw[BookMuted, dashed, line width=0.45pt] (3,2.3) -- (3,5.15);
  \draw[BookMuted, dashed, line width=0.45pt] (4,0.95) -- (4,5.15);
  \draw[BookMuted, dashed, line width=0.45pt] (7,0.95) -- (7,5.15);
  \draw[BookMuted, dashed, line width=0.45pt] (10,0.95) -- (10,5.15);
  % 源域数据：0111 -> 1000，四位同时变化
  \draw[BookTeal, line width=0.9pt] (0,2.8) rectangle (3,3.6);
  \node at (1.5,3.2) {0111};
  \draw[BookTeal, line width=0.9pt] (3,2.8) rectangle (11.5,3.6);
  \node at (7.25,3.2) {1000};
  \node[hw label, anchor=east] at (-0.2,3.2) {源域数据};
  % 接收时钟：采样沿在 4/7/10
  \draw[BookAccent, line width=0.9pt] (0,4.4) -- (4,4.4) -- (4,5.0) -- (5.5,5.0) -- (5.5,4.4) -- (7,4.4) -- (7,5.0) -- (8.5,5.0) -- (8.5,4.4) -- (10,4.4) -- (10,5.0) -- (11.5,5.0);
  \node[hw label, anchor=east] at (-0.2,4.7) {接收时钟};
  % 接收端采样值：0111 -> 0000（幻影）-> 1000
  \draw[BookTeal, line width=0.9pt] (0,1.2) rectangle (4,2.0);
  \node at (2,1.6) {0111};
  \draw[BookRed, line width=0.9pt] (4,1.2) rectangle (7,2.0);
  \node[text=BookRed] at (5.5,1.6) {0000};
  \draw[BookTeal, line width=0.9pt] (7,1.2) rectangle (11.5,2.0);
  \node at (9.25,1.6) {1000};
  \node[hw label, anchor=east] at (-0.2,1.6) {接收端采样值};
  \node[hw label, text=BookRed] at (5.5,0.6) {幻影值：源从未发送过 0000};
\end{waveform}
\diagramnote{4-bit 计数值从 0111 跳变到 1000，四个 bit 同时变化；各路两级同步器的走线和器件延迟略有差异，最高位的新值到得比低三位晚。接收端在 x=4 的采样沿采到低三位的新值和最高位的旧值，于是依次看到 0111、0000、1000——中间的 0000（红色）是源从未发送过的幻影值。逐级加触发器消除的是亚稳态，消除不了各位之间的采样时刻分歧；多 bit 跨域必须靠握手、格雷码或异步 FIFO 把数值作为整体传递。}

\begin{longtable}{L{0.24\linewidth}L{0.34\linewidth}L{0.32\linewidth}}
\toprule
跨域对象与做法 & 机制 & 代价与适用 \\
\midrule
单 bit 电平或事件：两级同步器 & 第一级吸收亚稳态，第二级给出稳定电平 & 延迟一两拍；适用按钮、中断、状态标志 \\
多 bit 批量数据：请求/应答握手 & 源先放好数据再发请求，数据保持稳定直到应答返回 & 每次传输数拍；适用低速、非连续传输 \\
多 bit 连续数据流：异步 FIFO & 读写两侧各用自己的时钟，指针用格雷码跨域传递 & 结构复杂、占用资源；适用连续数据流 \\
\bottomrule
\end{longtable}

握手的要点是“数据稳定在先，请求在后，撤销在应答之后”。源域把数据总线放好并保持不变，然后发出 request；request 经两级同步进入接收域时，数据早已稳定了好几个周期，接收方采数据不存在时刻分歧；接收方采完后发 acknowledge，同样同步回源域；源域收到应答之后才允许改变数据。一次传输多花几拍，换来的是多 bit 值作为整体被采走。异步 FIFO 则把这套思想结构化和流水化：写指针在源域递增，读指针在接收域递增，两个指针必须跨域比较才能判断空满；指针采用格雷码，相邻计数值只差一个 bit，于是指针跨域时可以安全地逐位同步——最多把指针读旧一拍，不会读出幻影编码。格雷码在这里的角色值得单独记住：它不是检错手段，而是把“多 bit 同时变化”约束成“一次只变一个 bit”，让逐位同步重新合法化。

把一次四相握手画成波形，数据、请求、应答的先后次序就固定下来了：

\begin{waveform}[x=0.8cm,y=1cm]
  % 数据总线：先放稳，传输全程不动
  \draw[draw=BookMuted, fill=BookCodeBg, line width=0.7pt] (-0.5,3.95) rectangle (0.6,4.45);
  \node[hw label] at (0.05,4.2) {旧};
  \draw[draw=BookTeal, fill=BookCodeBg, line width=0.7pt] (0.6,3.95) rectangle (7.6,4.45);
  \node[hw label] at (4.1,4.2) {新数据 D：传输全程保持稳定};
  \draw[draw=BookMuted, fill=BookCodeBg, line width=0.7pt] (7.6,3.95) rectangle (10.8,4.45);
  \node[hw label] at (9.2,4.2) {下一数据};
  \node[hw label, anchor=east] at (-0.7,4.2) {data};
  % req（源域）
  \draw[BookAccent, line width=1.0pt] (-0.5,3.2) -- (1,3.2) -- (1,3.52) -- (7,3.52) -- (7,3.2) -- (10.8,3.2);
  \node[hw label, anchor=east] at (-0.7,3.36) {req};
  % req 同步入接收域：晚约两拍
  \draw[BookTeal, line width=1.0pt] (-0.5,2.35) -- (3,2.35) -- (3,2.67) -- (9,2.67) -- (9,2.35) -- (10.8,2.35);
  \node[hw label, anchor=east] at (-0.7,2.51) {req 同步后};
  \draw[BookMuted, dashed, line width=0.5pt, -{Stealth[length=1.6mm]}] (1,3.05) -- (3,3.05) -- (3,2.72);
  \node[hw label, anchor=south] at (2,3.08) {两级同步};
  % ack（接收域）
  \draw[BookAccent, line width=1.0pt] (-0.5,1.5) -- (4,1.5) -- (4,1.82) -- (10,1.82) -- (10,1.5) -- (10.8,1.5);
  \node[hw label, anchor=east] at (-0.7,1.66) {ack};
  % ack 同步回源域
  \draw[BookTeal, line width=1.0pt] (-0.5,0.65) -- (6,0.65) -- (6,0.97) -- (10.8,0.97);
  \node[hw label, anchor=east] at (-0.7,0.81) {ack 同步后};
  \draw[BookMuted, dashed, line width=0.5pt, -{Stealth[length=1.6mm]}] (4,1.35) -- (6,1.35) -- (6,1.02);
  \node[hw label, anchor=north] at (5,1.28) {两级同步};
  % 四个相位标注
  \node[hw label, anchor=west] at (1.08,3.72) {① 数据稳定后才发请求};
  \node[hw label, anchor=west] at (4.08,2.06) {② 接收方采数据，回应答};
  \node[hw label, anchor=west] at (7.08,3.72) {③ 撤请求、改数据};
  \node[hw label, anchor=north] at (10,1.36) {④ 撤应答，回到空闲};
  \node[hw label] at (5.15,0.15) {一次传输多花数拍，换来多 bit 值作为整体被采走};
\end{waveform}
\diagramnote{req/ack 四相握手的一次传输：源域先把 data 放稳，再发 req（①）；req 经两级同步晚约两拍进入接收域，此刻数据早已稳定数个周期，接收方采数据不存在时刻分歧，采完回 ack（②）；ack 同步回源域后，源域才撤 req 并允许改变数据（③）；接收方看到 req 落下再撤 ack，双方回到空闲（④）。下一次传输（右侧下一数据）必须在 ③ 之后才开始。}

安全启动控制板中的 \code{start_request} 可以作为例子。按钮原始触点先经过默认电平、保护和消抖，得到板内稳定请求；若后级控制器是同步状态机，这个请求还要在控制器时钟域内同步，状态机只能使用同步后的 \code{start_request_sync}。这样第一章的“可靠 bit”与本章的“可靠采样边界”才真正接上：前者解决电平和路径，后者解决保存时刻。

\begin{chipcase}
用 74HC14 和 74HC74 可以搭出一个完整的单 bit 跨边界入口。

74HC14 负责把按钮或慢边沿信号整形成确定翻转的逻辑电平，74HC74 的两级触发器负责把该电平带入目标时钟域。若后级是计数器或状态机，只允许使用第二级 Q；第一级 Q 是给亚稳态恢复留时间的内部边界，不应直接参与控制。这个小电路比“按钮直接接计数器 CP”更接近真实机器的输入入口。
\end{chipcase}

\topic{专题：从外部事件到同步状态机}

许多同步电路的故障并不发生在核心状态机内部，而是发生在外部事件进入状态机之前。按钮、传感器、限位开关、中断线、另一片芯片的 ready 信号，都可能在本时钟域的任意时刻变化。若把这些信号直接接到状态机输入，设计等于默认外部世界服从本时钟的建立时间和保持时间；这个默认在真实硬件中通常不成立。

可以把外部事件进入同步状态机的链路写成六段：

\begin{lstlisting}
raw pin
  -> electrical default and protection
  -> filtering or debounce
  -> level normalization
  -> clock-domain synchronization
  -> event detection or handshake
  -> FSM consumption
\end{lstlisting}

第一段处理电气边界。原始引脚不能悬空，不能让过压、静电或长线干扰直接进入芯片内部；上拉、下拉、限流、钳位、RC 滤波和施密特输入都属于这一层。第二段处理物理抖动或噪声。按钮按下不是一个理想阶跃，可能在数毫秒内多次跳变；传感器线缆也可能受到干扰。第三段把电平极性转成内部语义，例如把低有效的 \code{start_n} 转成高有效的 \code{start_pressed}。

第四段才进入本章的核心：同步。一个单 bit 电平在消抖后仍然是异步输入，因为它的变化时刻不由目标时钟决定。两级同步器给第一级触发器留下恢复时间，降低亚稳态传播概率。同步后的信号仍然只是“本时钟域看到的电平”，若状态机需要的是一次事件，还必须再做边沿检测、事件保持或握手确认。

这条链上的每一段都不能替代另一段：默认电平与保护解决悬空、过压、静电和异常电流，但不能保证采样时刻安全；滤波与消抖解决机械抖动、窄噪声和线缆干扰，但不能消除跨时钟亚稳态；语义归一化解决低有效、高有效和断线默认的含义问题，但不能证明事件只发生一次；两级同步把单 bit 异步电平带进目标时钟域，但不能直接处理多 bit 一致性；边沿检测把同步电平转换成单周期事件，但不能保证目标状态机一定接受；握手或事件保持让事件持续到被消费并确认，但不能替代前面的电气保护。跳段的代价不是少一个步骤，而是某一层风险被悄悄漏过去。

以安全启动按钮为例，较完整的命名应当区分每一层信号：

\begin{lstlisting}
start_button_raw       // 引脚上的原始电平
start_button_filtered  // 经过默认电平、保护和消抖
start_pressed          // 归一化后的语义电平
start_pressed_meta     // 第一级同步器输出
start_pressed_sync     // 第二级同步器输出
start_event            // 本时钟域检测到的一次请求
start_pending          // 等待状态机消费的保持事件
\end{lstlisting}

这些名字不是形式主义，而是调试边界。若示波器看到 \code{start_button_raw} 抖动，问题属于电气或消抖；若 \code{start_pressed_sync} 偶发异常，重点检查同步器、时钟域和亚稳态扩散；若 \code{start_event} 出现两次，重点检查边沿检测和消抖窗口；若 \code{start_event} 只有一次但状态机没有启动，重点检查 \code{start_pending}、状态转移条件和消费确认。

事件保持尤其重要。假设同步后的 \code{start_event} 只有一个时钟周期，而状态机当前处在 \code{RESET_WAIT} 或 \code{FAULT}，这个事件可能被漏掉。更稳妥的做法是让事件进入 \code{start_pending} 寄存器，直到状态机在允许启动的状态下消费它，再由 \code{start_ack} 清除。

\begin{lstlisting}
on clock edge:
  if reset:
    start_pending = 0
  else if start_event:
    start_pending = 1
  else if start_ack:
    start_pending = 0
\end{lstlisting}

这段逻辑说明，外部事件不是直接驱动危险输出，而是先转化为本时钟域内的状态。状态机可以在 \code{IDLE} 中消费 \code{start_pending}，在 \code{FAULT} 中忽略或保留它，在复位期间清除它。这样，事件、状态和安全策略都由时钟边界定义，而不是由一根不受控的外部线决定。

\begin{warningbox}
两级同步器只适合单 bit 控制电平或事件标志。若外部输入是 8-bit 数据、地址、计数值或编码状态，逐位接两级同步器不能保证各位来自同一个源时刻。多 bit 跨域必须使用握手、有效信号加稳定保持、格雷码计数器或异步 FIFO 等协议。
\end{warningbox}

检查一条外部事件链路，可以分四个方面逐项确认。第一，电气处理：默认电平、保护、电压范围和抖动处理，缺失它就会出现悬空、抖动和多次触发。第二，同步处理：进入哪个时钟域、几级同步、复位后同步器处于什么默认状态，缺失它就会出现亚稳态扩散和复位后伪事件。第三，事件处理：电平如何变成一次事件，事件是否保持到被消费，缺失它就会出现窄脉冲漏采或重复消费。第四，状态机接口：哪些状态接受事件，哪些状态拒绝事件，拒绝时事件清除还是保留，缺失它就会出现故障状态误启动或请求丢失。

这个专题把第一章的电气可靠性、本章的同步边界和后续控制器的状态消费连接在一起。对硬件工程而言，“外部输入已经变成 0/1”不是终点；只有当它被目标时钟域可靠采样、被事件协议保持、被状态机按安全条件消费之后，它才成为可用于控制机器动作的内部事实。

\begin{waveform}[x=2pt,y=6ex]
  % CLK — 从低电平开始, 上升沿在 32/96/160/224/288pt
  \draw[BookAccent,line width=3pt]
    (0,22ex) -- ++(32pt,0) -- ++(0,6ex) -- ++(32pt,0) -- ++(0,-6ex)
             -- ++(32pt,0) -- ++(0,6ex) -- ++(32pt,0) -- ++(0,-6ex)
             -- ++(32pt,0) -- ++(0,6ex) -- ++(32pt,0) -- ++(0,-6ex)
             -- ++(32pt,0) -- ++(0,6ex) -- ++(32pt,0) -- ++(0,-6ex)
             -- ++(32pt,0) -- ++(0,6ex) -- ++(32pt,0) -- ++(0,-6ex)
             -- ++(16pt,0);
  \node[BookAccent,left] at (0,25ex) {CLK};
  % async input — 在采样沿附近变化
  \draw[BookAmber,line width=3pt]
    (0,11ex) -- ++(90pt,0) -- ++(0,6ex) -- ++(130pt,0) -- ++(0,-6ex) -- ++(44pt,0);
  \node[BookAmber,left] at (0,14ex) {async input};
  % FF1 — 上升沿采到变化中的输入, 先亚稳态后恢复
  \draw[BookRed,line width=3pt] (0,-2ex) -- ++(96pt,0);
  \draw[BookRed,dashed,line width=2.6pt]
    (96pt,-2ex) .. controls (108pt,7ex) and (124pt,-5ex) .. (136pt,4ex);
  \draw[BookRed,line width=3pt] (136pt,4ex) -- ++(88pt,0);
  \draw[BookRed,dashed,line width=2.6pt]
    (224pt,4ex) .. controls (236pt,-5ex) and (252pt,7ex) .. (264pt,-2ex);
  \draw[BookRed,line width=3pt] (264pt,-2ex) -- ++(40pt,0);
  \node[BookRed,left] at (0,2ex) {FF$_1$ (meta)};
  % FF2 — 下一个上升沿采到已经恢复的值
  \draw[BookTeal,line width=3pt]
    (0,-14ex) -- ++(160pt,0) -- ++(0,6ex) -- ++(128pt,0) -- ++(0,-6ex) -- ++(16pt,0);
  \node[BookTeal,left] at (0,-11ex) {FF$_2$ (sync)};
  \node[BookRed,font=\small] at (116pt,32ex) {亚稳态};
  \node[BookTeal,font=\small\bfseries] at (180pt,-17ex) {稳定输出};
\end{waveform}
\diagramnote{两级同步器波形：异步输入在上升沿附近变化；第一级在上升沿采到变化中的输入，可能短暂亚稳态（红色虚线），但在下一个上升沿前恢复到明确值；第二级在同一个上升沿采到的是已经稳定的值，输出干净。}


\section{四、寄存器、计数器与移位结构}

一个触发器保存 1 bit。把多个触发器并排，并让它们共享同一个时钟边界，就得到多 bit 寄存器。寄存器可以理解为一组共享时钟和写使能的触发器。若写使能有效，在时钟沿采样输入；若写使能无效，保持原值。

\begin{lstlisting}
on clock edge:
  if write_enable:
    Q = D
  else:
    Q holds
\end{lstlisting}



寄存器的意义不只是“能存数”。它把连续传播的组合逻辑切成离散步骤。上一个周期的结果在寄存器输出端稳定存在，成为本周期组合逻辑的输入；本周期算出的结果，要到下一个时钟沿才成为新的状态。共享同一提交边界是关键：寄存器保存的是同一个边界上的一组 bit，不是一位一位随意变化。

多 bit 寄存器还承担“原子提交”的角色。若一个 8-bit 状态表示 \code{01111111} 到 \code{10000000} 的变化，多个 bit 会同时改变；后级逻辑应在同一个时钟边界之后看到新状态，而不是看到一部分旧 bit、一部分新 bit。若这些 bit 来自同一组触发器，时钟边界和传播延迟可以被统一分析；若它们来自不同异步来源，后级就可能看到不存在于设计语义中的混合值。因此，寄存器不只是若干触发器的并排复制，而是一个“这些 bit 同属一个提交时刻”的约定。

写使能寄存器通常有两种实现方式。第一种是在触发器前放数据选择器：使能为 1 时选择新数据，使能为 0 时选择旧 Q 反馈；第二种是使用器件或标准单元内部提供的 enable 触发器。两者的目标都是让时钟边界仍然存在，只是决定这个边界保存新值还是旧值。把 enable 直接拿去切时钟看起来也能保持状态，但会把 enable 的毛刺和传播延迟带进时钟树，风险远高于在数据路径上做选择。

\begin{block}[x=1.8cm,y=1.5cm]
  \node[hw compute, minimum width=2.2cm] (mux) {MUX};
  \node[hw state, right=1.0cm of mux, minimum width=2.0cm] (ff) {D-FF};
  \draw[hw bus] (ff.east) -- ++(0.8,0) node[right] (q_out) {Q};
  \draw[hw bus] (q_out) |- ++(0,1.2) -| (mux.north) node[pos=0.7, above right, hw label, BookTeal] {反馈};
  \draw[hw bus] (mux) -- (ff) node[above, hw label, pos=0.5]{D};
  \draw[hw bus] ($(mux.west)+(-0.8,0)$) node[left]{D$_{in}$} -- (mux.west);
  \node[hw control, below=0.8cm of mux, minimum width=1.4cm] (we) {WE};
  \draw[hw bus] (we.north) -- ($(we.north)+(0,0.34)$) -| (mux.south) node[pos=0.78, above right, hw label]{选择};
  \node[below=0.3cm of we, hw label] {0=保持, 1=写入};
\end{block}
\diagramnote{写使能寄存器：数据路径决定写入或保持。WE=0 时 mux 选择旧 Q 反馈，时钟照常到达但状态不变；WE=1 时选择新 \code{D\_in}。时钟树全程干净，这就是它优于门控时钟的地方。}

\begin{longtable}{L{0.22\linewidth}L{0.36\linewidth}L{0.32\linewidth}}
\toprule
寄存器动作 & 下一值表达 & 硬件含义 \\
\midrule
保持 & \code{next = current} & 写回旧 Q 或关闭写使能 \\
装载 & \code{next = input} & 在边界保存外部或组合结果 \\
清零 & \code{next = 0} & 同步清零或受控复位路径 \\
置位 & \code{next = all\_ones} & 标志位、掩码或诊断状态写入 \\
条件更新 & \code{next = cond ? input : current} & 由控制线决定是否提交 \\
\bottomrule
\end{longtable}

计数器的下一值来自当前值加一：

\begin{lstlisting}
next = current + 1
on clock edge:
  current = next
\end{lstlisting}



该结构里有反馈，但反馈不是直接绕回去，而是经过加法器，并且被时钟边界切开。若 current 为 3，本周期组合逻辑算出 next 为 4；只有下一个时钟沿到来，寄存器才保存 4。随后组合逻辑再从 4 算出 5。如果没有寄存器边界，\code{current+1} 的结果会立刻回到输入，再加一，再回到输入，电路无法表达“每个周期只增加一次”。

\begin{block}[x=1.8cm,y=1.5cm]
  \node[hw compute, minimum width=2.4cm] (add) {+1};
  \node[hw state, right=1.0cm of add, minimum width=2.2cm] (ff) {4-bit 寄存器};
  \draw[hw bus] (add) -- (ff) node[above, hw label, pos=0.5]{next};
  \draw[hw bus] (ff.east) -- ++(0.6,0) node[right] (q) {Q[3:0]};
  \draw[hw bus] (q) |- ++(0,1.2) -| (add.north) node[pos=0.7, above right, hw label, BookTeal] {当前值};
  \node[hw control, below=0.8cm of add, minimum width=1.4cm] (en) {EN};
  \draw[hw bus] (en.north) -- ($(en.north)+(0,0.34)$) -| (add.south) node[pos=0.78, below right, hw label]{使能};
  \node[below=0.3cm of en, hw label] {0=保持, 1=加一};
\end{block}
\diagramnote{4-bit 同步计数器。}

把 3、4、5、6、7 的计数过程画成波形，"整体提交"和"逐位爬"的区别就一目了然：

\begin{waveform}[x=1cm,y=1cm]
  \draw[BookAccent, line width=1.0pt] (0,3.4) -- (2,3.4) -- (2,4.2) -- (3,4.2) -- (3,3.4) -- (4,3.4) -- (4,4.2) -- (5,4.2) -- (5,3.4) -- (6,3.4) -- (6,4.2) -- (7,4.2) -- (7,3.4) -- (8,3.4) -- (8,4.2) -- (9,4.2) -- (9,3.4) -- (10,3.4) -- (10,4.2) -- (11,4.2);
  \node[hw label, anchor=east] at (-0.2,3.8) {CLK};
  \draw[BookTeal, line width=0.9pt] (0,1.5) -- (10.2,1.5) -- (10.2,2.3) -- (0,2.3) -- cycle;
  \foreach \x in {2.2,4.2,6.2,8.2} {
    \draw[BookTeal, line width=0.9pt] (\x,1.5) -- (\x,2.3);
  }
  \node at (1.1,1.9) {3};
  \node at (3.2,1.9) {4};
  \node at (5.2,1.9) {5};
  \node at (7.2,1.9) {6};
  \node at (9.2,1.9) {7};
  \node[hw label, anchor=east] at (-0.2,1.9) {Q[3:0]};
  \draw[BookAmber, line width=1.0pt] (0,0.9) -- (2.2,0.9) -- (2.2,0.3) -- (4.2,0.3) -- (4.2,0.9) -- (6.2,0.9) -- (6.2,0.3) -- (8.2,0.3) -- (8.2,0.9) -- (10.2,0.9) -- (10.2,0.3) -- (11,0.3);
  \node[hw label, anchor=east] at (-0.2,0.6) {Q0};
  \node[hw label] at (5.6,2.72) {每个上升沿后 $t_{CQ}$，整体值同时 +1（不是一位一位爬）};
\end{waveform}
\diagramnote{4-bit 计数器波形：每个上升沿，寄存器把“当前值+1”整体提交一次，总线带上的值从 3 整齐跳到 4、5、6、7。注意 Q0——它每个周期翻转，正好是 CLK 的二分频；Q1 再慢一半。计数器的每个 bit 都是比它低一位的半速时钟，这是后面分频和定时器的基础。}

74HC161/74HC163 这类 4-bit 计数器把这种结构做成了可直接使用的芯片。它们通常有时钟、计数使能、并行装载、清零、进位输出和 4 个 Q 输出。计数使能决定是否执行 \code{current+1}，并行装载允许在边界写入外部给定值，进位输出用于把多个计数器级联成更宽的计数器。读这类芯片的数据手册时，不能只看“它会加一”，还要查清楚清零是同步还是异步、装载优先级如何、使能脚无效时是否保持、进位输出何时有效。

\begin{longtable}{L{0.2\linewidth}L{0.38\linewidth}L{0.32\linewidth}}
\toprule
计数器引脚/功能 & 对应的状态问题 & 典型用途 \\
\midrule
CP/CLK & 哪个边界提交新计数值 & 单步时钟、系统时钟、分频后时钟 \\
ENP/ENT 等使能 & 本周期是否执行加一 & PC 加一、循环计数、定时等待 \\
并行装载 & 是否用外部值覆盖加一结果 & 跳转地址、预置计数、状态恢复 \\
清零/复位 & 初始状态如何确定 & 上电入口、错误恢复、实验复位 \\
进位输出 & 更宽计数器何时进位 & 级联 8-bit、16-bit 或更宽计数器 \\
\bottomrule
\end{longtable}

把两个 4-bit 计数器级联成 8-bit 计数器时，低 4 位的进位不能被理解为普通数据输出。它是高 4 位是否在同一边界加一的控制条件，必须满足高位计数器的使能时序。若把低位 Q 经过随意组合逻辑再驱动高位时钟，就把同步计数器退化成脉动计数器，时序关系会变得难以分析。同步级联的要点是：所有位共享同一个时钟，进位只作为下一状态选择条件。

\topic{专题：纹波计数器与同步计数器对照}

上一段提醒过：级联计数器时若把低位 Q 经过随意组合逻辑去驱动高位时钟，同步结构就退化成了脉动计数器。脉动计数器也叫\hwrisk{纹波计数器}（ripple counter），它和同步计数器完成同一个功能，但时钟的组织方式完全不同，值得放在一起对照。

纹波计数器的做法极为节省：每个触发器接成翻转模式（把 D 触发器的 $\overline{Q}$ 接回 D），第一级由外部时钟驱动，之后每一级的时钟直接取自前一级的 Q 输出。第 0 位在时钟沿之后先翻转；第 1 位看到第 0 位变化后才跟着翻转；如此逐级传递，第 $n$ 位要等前面各级依次翻完才轮到本级，延迟随位号一级一个 $t_{CQ}$ 累加。变化像波浪一样从低位传到高位，这正是“纹波”二字的来历。

这个逐级延迟不是理论细节。以 4-bit 纹波计数器从 0111 计到 1000 为例，设每级 $t_{CQ}=1$ns：

\begin{longtable}{L{0.14\linewidth}L{0.30\linewidth}L{0.44\linewidth}}
\toprule
时钟沿后时刻 & 已完成的动作 & 总线上读到的值 \\
\midrule
0ns & 外部时钟沿到达，各级尚未更新 & 0111（旧值 7） \\
1ns & 第 0 位翻转完成 & 0110（读数 6，中间值） \\
2ns & 第 1 位跟随翻转 & 0100（读数 4，中间值） \\
3ns & 第 2 位跟随翻转 & 0000（读数 0，中间值） \\
4ns & 第 3 位最后翻转 & 1000（读数 8，最终值） \\
\bottomrule
\end{longtable}

在最终值稳定之前，总线上依次出现了 6、4、0 三个并不存在于计数序列中的中间值。若后级只是低速显示驱动，这几纳秒毫无影响；但若后级是译码逻辑——例如“计数值等于 8 时打开某写使能”——中间值 0 就可能让另一路译码输出意外有效，制造窄毛刺。纹波计数器的风险不在计数错误，而在读取窗口内的值不可信。

把表中的过程按纳秒画成波形，三个中间值各停留 1ns 的读取窗口就一目了然：

\begin{waveform}[x=1.4cm,y=1cm]
  % 参考虚线：1/2/3/4ns 各级翻转时刻
  \foreach \x in {1,2,3,4} {\draw[BookMuted, dashed, line width=0.45pt] (\x,-0.7) -- (\x,5.5);}
  % CLK：0ns 一个上升沿
  \draw[BookAccent, line width=1.0pt] (-0.6,5.0) -- (0,5.0) -- (0,5.35) -- (5.2,5.35);
  \node[hw label, anchor=east] at (-0.8,5.17) {CLK};
  % Q0：沿后 1ns 翻成 0
  \draw[BookTeal, line width=1.0pt] (-0.6,4.35) -- (1,4.35) -- (1,4.0) -- (5.2,4.0);
  \node[hw label, anchor=east] at (-0.8,4.17) {Q$_0$};
  % Q1：沿后 2ns 翻成 0
  \draw[BookTeal, line width=1.0pt] (-0.6,3.35) -- (2,3.35) -- (2,3.0) -- (5.2,3.0);
  \node[hw label, anchor=east] at (-0.8,3.17) {Q$_1$};
  % Q2：沿后 3ns 翻成 0
  \draw[BookTeal, line width=1.0pt] (-0.6,2.35) -- (3,2.35) -- (3,2.0) -- (5.2,2.0);
  \node[hw label, anchor=east] at (-0.8,2.17) {Q$_2$};
  % Q3：沿后 4ns 翻成 1
  \draw[BookTeal, line width=1.0pt] (-0.6,1.0) -- (4,1.0) -- (4,1.35) -- (5.2,1.35);
  \node[hw label, anchor=east] at (-0.8,1.17) {Q$_3$};
  % 延迟逐级累加标注
  \draw[{Stealth[length=1.6mm]}-{Stealth[length=1.6mm]}, BookRed, line width=0.55pt] (0,5.62) -- (4,5.62);
  \node[hw label, anchor=west, text=BookRed] at (4.15,5.62) {Q$_3$ 等满 4 个 $t_{CQ}$};
  % 总线读数行：红格为中间值
  \draw[draw=BookMuted, fill=BookCodeBg, line width=0.7pt] (-0.6,-0.3) rectangle (1,0.3);
  \node[hw label] at (0.2,0) {0111};
  \draw[draw=BookRed, fill=white, line width=0.7pt] (1,-0.3) rectangle (2,0.3);
  \node[hw label, text=BookRed] at (1.5,0) {0110};
  \draw[draw=BookRed, fill=white, line width=0.7pt] (2,-0.3) rectangle (3,0.3);
  \node[hw label, text=BookRed] at (2.5,0) {0100};
  \draw[draw=BookRed, fill=white, line width=0.7pt] (3,-0.3) rectangle (4,0.3);
  \node[hw label, text=BookRed] at (3.5,0) {0000};
  \draw[draw=BookTeal, fill=BookCodeBg, line width=0.7pt] (4,-0.3) rectangle (5.2,0.3);
  \node[hw label] at (4.6,0) {1000};
  \node[hw label, anchor=east] at (-0.8,0) {总线读数};
  % 时间刻度
  \foreach \x/\t in {0/0,1/1,2/2,3/3,4/4} {\node[hw label] at (\x,-0.75) {\t};}
  \node[hw label, anchor=west] at (5.25,-0.75) {ns};
\end{waveform}
\diagramnote{4-bit 纹波计数器从 0111 计到 1000（每级 $t_{CQ}=1$ns）：CLK 沿后 1ns $Q_0$ 先翻成 0，它的变化再触发 $Q_1$，逐级传递到 $Q_3$ 已是沿后 4ns。底部总线行依次读出 0111、0110、0100、0000、1000——红色三格是计数序列中不存在的中间值，后级译码逻辑在这个窗口读数就会看到幻影。同步计数器所有位共享时钟，沿后一步整体更新，没有这些格子。}

同步计数器走另一条路：所有触发器共享同一个时钟，每一位的下一值由组合逻辑根据当前各位统一算出（\code{next = current + 1}），在同一个时钟边界整体提交。前文 4-bit 计数器波形里“整体提交、不是一位一位爬”说的就是这件事。代价是每个触发器前面多了下一状态逻辑，位数越宽进位链越长；收益是时钟沿之后只有一级 $t_{CQ}$ 加组合路径的延迟，各位同时稳定，任何时刻读到的都是真实计数值。

速度差异可以用数字算出来。设 $t_{CQ}=1$ns：纹波结构 8 位时最高位要 $8\times1=8$ns 才稳定，时钟周期必须大于 8ns，上限约 125MHz，还未计任何译码余量；同步结构同样 8 位，若下一状态组合逻辑最大延迟 2ns、目标建立时间 0.4ns，周期下限约为 $1+2+0.4=3.4$ns，5ns 周期（200MHz）仍有富余。位数越多，纹波结构的级数惩罚线性增长，同步结构的组合延迟只缓慢增加。

\begin{longtable}{L{0.12\linewidth}L{0.20\linewidth}L{0.18\linewidth}L{0.18\linewidth}L{0.20\linewidth}}
\toprule
结构 & 时钟接法 & 稳定延迟 & 读数行为 & 适用场合 \\
\midrule
纹波计数器 & 每级时钟取自前一级 Q & 低位 1 个 $t_{CQ}$ 起，逐级累加 & 稳定前出现短暂中间值 & 低速分频链、定时链、读数不被同域逻辑消费的场合 \\
同步计数器 & 所有触发器共享时钟 & 一级 $t_{CQ}$ 加组合路径 & 边界后各位同时稳定 & 计数值要被同域译码、比较、控制的场合 \\
\bottomrule
\end{longtable}

选择原则因此很直接：只把计数器当分频器用、中间值无人读取，纹波结构简单省电；计数值要被同一时钟域的逻辑读取、译码或比较，就必须用同步计数器。这解释了为什么 74HC161/163 这类通用计数芯片全部做成同步结构，而纹波结构更多出现在专用分频链和异步预分频器里。

\topic{专题：任意模数计数器——复位法与预置法}

4-bit 计数器自然按模 16 循环，但十进制显示、BCD 运算和秒分计时需要的是模 10。把模 16 改成任意模 $N$，工程上有两种标准做法，差别在“如何让第 $N$ 个状态消失”。

第一种是复位法：让计数器照常加一，另用一个译码门监视计数值，一旦检出 $N$ 就立刻异步清零，不等下一个时钟沿。以模 10 为例，译码门检出 1010（十进制 10），清零端把各级直接拉回 0000：

\begin{longtable}{L{0.16\linewidth}L{0.24\linewidth}L{0.48\linewidth}}
\toprule
时钟沿序号 & 沿后稳定值 & 说明 \\
\midrule
0 & 0000 & 初始值 \\
1 & 0001 & 正常加一 \\
$\cdots$ & $\cdots$ & 依次递增 \\
9 & 1001 & 最后一个合法值（十进制 9） \\
10 & 0000 & 加一结果 1010 短暂出现，译码门检出后立即异步清零 \\
\bottomrule
\end{longtable}

注意第 10 拍：1010 并非完全不出现，它作为瞬态存在几纳秒——正是这段瞬态触发了清零。这带来复位法的固有弱点：清零脉冲的宽度只有“译码门延迟加清零响应”那么长，若各级触发器清零速度不齐，最快的一级先清掉会使译码条件消失、清零脉冲提前撤除，最慢的一级可能还没清完，计数器从一个错误值重新起步。复位法的优点是只要一个译码门，缺点是把可靠性押在各级清零延迟一致上。

第二种是预置法：不消灭第 $N$ 个状态，而是根本不让它出现。译码门检出的是最后一个合法值 $N-1$（模 10 时为 1001），检出后拉起同步装载控制，下一个时钟沿把初值 0000 装入，覆盖掉本来的加一结果：

\begin{longtable}{L{0.16\linewidth}L{0.24\linewidth}L{0.48\linewidth}}
\toprule
时钟沿序号 & 沿后稳定值 & 说明 \\
\midrule
0 & 0000 & 初始值 \\
1 & 0001 & 正常加一 \\
$\cdots$ & $\cdots$ & 依次递增 \\
9 & 1001 & 译码门检出末态，本拍内拉起装载控制 \\
10 & 0000 & 时钟沿把初值装入，1010 从头到尾不出现 \\
\bottomrule
\end{longtable}

预置法的每一次变化都发生在时钟边界上，没有任何瞬态；装载控制有整整一个周期的时间稳定，只需像普通同步输入一样满足建立时间。代价是每级前面多一级装载选择逻辑，且译码门在整个 1001 周期内都必须稳定有效，不能带毛刺。

把两种做法在回卷附近的几拍画在同一时间轴上，瞬态有无一眼可见：

\begin{waveform}[x=0.75cm,y=0.85cm]
  % 参考线：x=4 回卷沿（全高），x=2/6 采样沿（仅上半部）
  \draw[BookMuted, dashed, line width=0.45pt] (4,0.85) -- (4,6.5);
  \draw[BookMuted, dashed, line width=0.45pt] (2,3.7) -- (2,6.5);
  \draw[BookMuted, dashed, line width=0.45pt] (6,3.7) -- (6,6.5);
  % CLK：上升沿在 2/4/6/8/10/12
  \draw[BookAccent, line width=0.9pt] (0,5.6) -- (2,5.6) -- (2,6.2) -- (3,6.2) -- (3,5.6) -- (4,5.6) -- (4,6.2) -- (5,6.2) -- (5,5.6) -- (6,5.6) -- (6,6.2) -- (7,6.2) -- (7,5.6) -- (8,5.6) -- (8,6.2) -- (9,6.2) -- (9,5.6) -- (10,5.6) -- (10,6.2) -- (11,6.2) -- (11,5.6) -- (12,5.6) -- (12,6.2) -- (13,6.2);
  \node[hw label, anchor=east] at (-0.2,5.9) {CLK};
  % 复位法 Q：1010 作为瞬态出现几纳秒后被异步清零
  \draw[BookTeal, line width=0.9pt] (0,4.0) rectangle (2,4.8);
  \node at (1,4.4) {1000};
  \draw[BookTeal, line width=0.9pt] (2,4.0) rectangle (4,4.8);
  \node at (3,4.4) {1001};
  \draw[BookRed, line width=0.9pt] (4,4.0) rectangle (4.5,4.8);
  \draw[BookTeal, line width=0.9pt] (4.5,4.0) rectangle (6,4.8);
  \node at (5.25,4.4) {0000};
  \draw[BookTeal, line width=0.9pt] (6,4.0) rectangle (8,4.8);
  \node at (7,4.4) {0001};
  \draw[BookTeal, line width=0.9pt] (8,4.0) rectangle (10,4.8);
  \node at (9,4.4) {0010};
  \node[hw label, anchor=east] at (-0.2,4.4) {复位法 Q};
  \node[hw label, text=BookRed, anchor=west] at (4.6,5.12) {1010};
  % 译码门输出：检出 1010 后的窄清零脉冲
  \draw[BookRed, line width=0.9pt] (0,2.9) -- (4.05,2.9) -- (4.05,3.5) -- (4.5,3.5) -- (4.5,2.9) -- (13,2.9);
  \node[hw label, anchor=east] at (-0.2,3.2) {译码门输出};
  \node[hw label, text=BookRed, anchor=west] at (4.65,3.68) {窄脉冲};
  % 预置法 Q：1010 从头到尾不出现
  \draw[BookTeal, line width=0.9pt] (0,1.3) rectangle (2,2.1);
  \node at (1,1.7) {1000};
  \draw[BookTeal, line width=0.9pt] (2,1.3) rectangle (4,2.1);
  \node at (3,1.7) {1001};
  \draw[BookTeal, line width=0.9pt] (4,1.3) rectangle (6,2.1);
  \node at (5,1.7) {0000};
  \draw[BookTeal, line width=0.9pt] (6,1.3) rectangle (8,2.1);
  \node at (7,1.7) {0001};
  \draw[BookTeal, line width=0.9pt] (8,1.3) rectangle (10,2.1);
  \node at (9,1.7) {0010};
  \node[hw label, anchor=east] at (-0.2,1.7) {预置法 Q};
  % 装载控制：1001 整拍内有效
  \draw[BookAccent, line width=0.9pt] (0,0.1) -- (2,0.1) -- (2,0.7) -- (4,0.7) -- (4,0.1) -- (13,0.1);
  \node[hw label, anchor=east] at (-0.2,0.4) {装载控制};
  \node[hw label, anchor=west] at (4.2,1.0) {检出 1001，整拍有效};
\end{waveform}
\diagramnote{复位法（上）：第 10 个时钟沿后，加一结果 1010 作为瞬态出现几纳秒（红色窄格），译码门检出后输出一条窄清零脉冲把各级拉回 0000——脉冲宽度只有“译码门延迟加清零响应”那么长，各级清零速度不齐时可能清不干净。预置法（下）：译码门在 1001 的整个周期内拉起装载控制，下一个时钟沿把 0000 装入，1010 从头到尾不出现，序列里没有任何瞬态。}

\begin{longtable}{L{0.12\linewidth}L{0.14\linewidth}L{0.22\linewidth}L{0.14\linewidth}L{0.26\linewidth}}
\toprule
做法 & 译码的状态 & 动作方式 & 有无瞬态 & 关键问题 \\
\midrule
复位法 & 检出 $N$（如 1010） & 异步清零，不等时钟沿 & 有，$N$ 短暂出现 & 清零脉冲仅数纳秒，各级清零不齐时可能清不干净 \\
预置法 & 检出 $N-1$（如 1001） & 同步装载，等下一时钟沿 & 无 & 装载控制须整拍稳定并满足建立时间，多用一级选择逻辑 \\
\bottomrule
\end{longtable}

两种做法在芯片族里都有对应型号：74HC160/161 用异步清零，74HC162/163 用同步清零。同一族芯片同时提供两种清零方式，正说明这个差别是真实的设计选择，而不是教科书里的文字游戏。BCD 计数器（模 10）只是最常用的一例；秒针的模 60、地址回卷的模 $2^k$ 分区，用的都是同一对方法。

\topic{专题：格雷码、环形与约翰逊计数器}

二进制顺序不是唯一的计数方式。改变“下一状态是什么”的规则，可以得到三类各有专长的计数器，它们的取舍恰好说明：状态编码本身就是设计。

格雷码计数器让相邻状态只差一位。3-bit 格雷码的顺序是 000、001、011、010、110、111、101、100，随后回到 000。普通二进制从 0111 计到 1000 时四位同时变化，异步采样器在变化途中可能读到任意混合值；格雷码每步只变一位，采样器最坏读到旧值或新值，绝不可能读出第三种组合。本章前面讲两级同步器时提过“多 bit 跨域要用格雷码计数器”，原因就在这里：异步 FIFO 的读写指针用格雷码编码，指针跨域采样时的多值风险被结构性消除。代价是下一状态逻辑比二进制加一复杂，且状态顺序不再是数值顺序，不能直接当地址用。

环形计数器把移位寄存器的最后一级接回串行输入，形成一个环。4 级环以种子值 0001 起步，依次经过 0001、0010、0100、1000，再回到 0001。$N$ 个触发器只得到 $N$ 个状态，状态利用率很低，但换来一个独特性质：每个触发器的输出本身就是状态标志，译码连一个门都不用——想知道“是否处于第 2 态”，看第 2 根输出线即可。独热编码的状态机正是这个结构的推广：把带种子值的环形计数器当作状态寄存器，一个触发器对应一个状态；反过来看，环形计数器可以看成转移固定的独热状态机特例。后文状态编码专题里说的“one-hot 译码直接”，说的就是这件事。

约翰逊计数器（扭环计数器）只改一处：接回串行输入的不是最后一级本身，而是它的反相值。4 级约翰逊从 0000 起步，依次经过 0000、0001、0011、0111、1111、1110、1100、1000，再回到 0000——$N$ 个触发器得到 $2N$ 个状态，比环形多一倍。更妙的是译码：每一步仍只有一位变化，且每个状态都可以由唯一一对相邻位（含首尾相接的一对）识别，译码一个状态只需一个两输入门。例如状态 0000 由“$Q_3=0$ 且 $Q_0=0$”唯一确定，0111 由“$Q_3=0$ 且 $Q_2=1$”确定。译码门少、每步只变一位，使约翰逊计数器的译码输出干净无毛刺，常用于多相时钟生成、步进电机相序和简单分频。

\begin{longtable}{L{0.13\linewidth}L{0.11\linewidth}L{0.16\linewidth}L{0.20\linewidth}L{0.28\linewidth}}
\toprule
结构 & 4 级状态数 & 每步变化位数 & 译码代价 & 典型用途 \\
\midrule
二进制 & 16 & 最多 4 位同时变 & 全译码，变化途中可能出中间值 & 通用计数、地址生成 \\
格雷码 & 16 & 恒为 1 位 & 译码无中间态 & 跨时钟域指针、角度与位置编码 \\
环形 & 4 & 2 位（一位退、一位进） & 无需译码门，Q 直接当状态标志 & 独热状态机、简单轮转控制 \\
约翰逊 & 8 & 恒为 1 位 & 每态一个两输入门 & 多相时钟、步进电机相序、分频 \\
\bottomrule
\end{longtable}

三类结构共用同一个提醒：未在序列里的编码必须有恢复路径。环形计数器若因上电扰动出现两个 1，就会在错误序列里一直循环；可靠设计要在上电时置入种子值，或加自检逻辑把非法编码拉回主循环——这和状态机“未使用编码要有默认去向”是同一条约定。

\topic{专题：LFSR 伪随机计数器}

把移位寄存器的某几位抽出来做异或，把异或结果接回串行输入，就得到线性反馈移位寄存器（LFSR）。抽头位置用多项式记号表示：$x^4+x^3+1$ 表示 4 级结构中取第 4 位和第 3 位（即 $Q_3$、$Q_2$）做异或反馈。设状态写作 $Q_3Q_2Q_1Q_0$，每个时钟沿各位左移一格，反馈值 $f=Q_3\oplus Q_2$ 装入 $Q_0$：

\begin{lstlisting}
f = Q3 XOR Q2
on clock edge:
  Q3 = Q2; Q2 = Q1; Q1 = Q0; Q0 = f
\end{lstlisting}

画成结构图，数据在四个触发器之间逐拍左移，唯一的组合逻辑是那只异或门：

\begin{pdfdiagram}
  \node[hw label] at (3.3,0.85) {4 级 LFSR：$x^4+x^3+1$，遍历 15 个非零状态};
  \node[hw state, minimum width=1.1cm, minimum height=0.75cm] (q3) at (0,0) {Q$_3$};
  \node[hw state, minimum width=1.1cm, minimum height=0.75cm] (q2) at (2.2,0) {Q$_2$};
  \node[hw state, minimum width=1.1cm, minimum height=0.75cm] (q1) at (4.4,0) {Q$_1$};
  \node[hw state, minimum width=1.1cm, minimum height=0.75cm] (q0) at (6.6,0) {Q$_0$};
  \draw[hw bus] (q0.west) -- (q1.east);
  \draw[hw bus] (q1.west) -- (q2.east);
  \draw[hw bus] (q2.west) -- (q3.east);
  \draw[hw bus] (q3.west) -- ++(-0.9,0);
  \node[hw label, anchor=south] at (-1.35,0.12) {序列输出};
  % 抽头异或反馈
  \node[hw compute, minimum width=2.3cm, minimum height=1.0cm] (xor) at (1.5,-1.9) {XOR\\$f=Q_3\oplus Q_2$};
  \draw[hw bus] (q3.south) |- (xor.west);
  \draw[hw bus] (q2.south) -- ($(q2.south)+(0,-0.6)$) -| (xor.north);
  \node[hw label, anchor=east] at (-0.12,-1.1) {抽头};
  \node[hw label, anchor=west] at (2.35,-1.0) {抽头};
  \draw[hw return] (xor.east) -- (6.6,-1.9) -- (q0.south);
  \node[hw label, fill=white, inner sep=1pt] at (4.6,-1.62) {反馈 $f$ 装入 $Q_0$};
  \node[hw label] at (3.3,-2.75) {四个触发器共享同一时钟，沿上同时移位};
\end{pdfdiagram}
\diagramnote{4 级 LFSR（$x^4+x^3+1$）：每个时钟沿 $Q_0\to Q_1\to Q_2\to Q_3$ 同时移位，抽头 $Q_3$、$Q_2$ 经唯一一只异或门算出反馈 $f$ 装入 $Q_0$。这组抽头对应本原多项式，从非零种子出发遍历全部 15 个非零状态后精确重复；全零是死状态，上电必须置入非零种子。}

这组抽头选得恰当（$x^4+x^3+1$ 是 4 级的本原多项式之一），LFSR 会遍历全部 $2^4-1=15$ 个非零状态才重复。从种子 0001 起步，前 8 拍为：

\begin{longtable}{L{0.10\linewidth}L{0.22\linewidth}L{0.24\linewidth}L{0.32\linewidth}}
\toprule
拍 & 状态 $Q_3Q_2Q_1Q_0$ & 反馈 $f=Q_3\oplus Q_2$ & 说明 \\
\midrule
0 & 0001 & 0 & 种子值，必须非零 \\
1 & 0010 & 0 & \\
2 & 0100 & 1 & \\
3 & 1001 & 1 & \\
4 & 0011 & 0 & \\
5 & 0110 & 1 & \\
6 & 1101 & 0 & \\
7 & 1010 & 1 & 之后继续遍历其余 7 态，第 15 拍回到 0001 \\
\bottomrule
\end{longtable}

把 15 个状态写成十进制是 1、2、4、9、3、6、13、10、5、11、7、15、14、12、8，随后精确重复。这个序列“伪”在完全确定、可由种子重现，“随机”在局部统计性质接近真随机序列：任一输出位上 0 和 1 的出现次数只差一次，游程分布也与随机序列一致，因此得名伪随机。

LFSR 的价值在于用极少的逻辑产生长序列。与二进制计数器相比，它没有进位链：每个触发器的下一状态只是一根移位线，唯一的一级组合逻辑是那个异或门，因此同样位宽下更快、更省。三类经典用途：一是伪随机测试，芯片内建自测（BIST）用 LFSR 给被测电路灌激励，再压缩响应做签名分析；二是定时器和大模数计数，凡是只问“数满多少个周期”、不问数值顺序的场合，LFSR 都比二进制计数器划算；三是扰码，通信链路把数据流与 LFSR 序列异或，打散长连 0 和长连 1，让频谱和时钟恢复更友好。

必须记住两个约束。第一，全零是死状态：$f=0\oplus0=0$，一旦进入 0000 就永远停在那里，所以上电必须置入非零种子，或加全零检测逻辑把序列拉回。第二，状态顺序不是数值顺序，需要“第几号状态”的场合要查表，不能直接当地址计数器用。另外，LFSR 的内核正是一个移位寄存器——接下来就来正式看移位结构本身。

移位寄存器把一组 bit 向左或向右移动。它可以用于串行输入、串行输出、简单延迟线和位级变换。

\begin{lstlisting}
next[3] = current[2]
next[2] = current[1]
next[1] = current[0]
next[0] = serial_in
\end{lstlisting}



假设 current 为 \code{1010}，\code{serial\_in} 为 1。下一个时钟沿后，寄存器变成 \code{0101}。注意，四个位不是按顺序一个一个搬。组合逻辑同时准备好四个 next 输入，时钟沿到来时一起写入。该例说明状态更新不一定是加法，也可以是重新排列、选择、清零、保持或装载外部输入。

移位结构在整机里非常常见。串行外设把一个引脚上的 bit 流变成内部并行字节，显示驱动把内部字节逐位送到外部，流水线把阶段之间的状态向前推进，调试链把片内状态一位一位移出。74HC595 的移位寄存器与输出锁存器分开，正好说明两级状态边界：移位时钟改变内部移位状态，锁存时钟才改变外部可见输出。这个结构避免了外部 LED 或后级芯片在移位过程中看到中间状态乱跳。

\begin{longtable}{L{0.22\linewidth}L{0.36\linewidth}L{0.32\linewidth}}
\toprule
74HC595 边界 & 发生的状态动作 & 读者应观察到的现象 \\
\midrule
\code{SHCP} 边沿 & 串行 bit 进入内部移位寄存器 & 外部输出可保持不变 \\
连续 8 次移位 & 内部移位状态逐步形成一个字节 & 中间状态不应直接驱动负载语义 \\
\code{STCP} 边沿 & 内部字节提交到输出锁存器 & Q0 到 Q7 一起更新 \\
\code{OE_n}/\code{MR_n} & 控制输出驱动和清零默认状态 & 低有效脚必须处于确定电平 \\
\bottomrule
\end{longtable}

\begin{block}[x=2.0cm,y=1.6cm]
  \node[hw state, minimum width=3.2cm, minimum height=1.0cm, label={[hw label]above:8-bit 移位寄存器}] (shift) {};
  \node[hw memory, right=1.0cm of shift, minimum width=3.2cm, minimum height=1.0cm, label={[hw label]above:8-bit 输出锁存}] (latch) {};
  \node[right=0.5cm of latch] (q) {Q$_0$..Q$_7$};
  \draw[hw bus] (shift) -- (latch) node[above, hw label, pos=0.5]{8-bit};
  \draw[hw bus] (latch) -- (q);
  \draw[hw bus] ($(shift.west)+(-0.6,0)$) node[left]{DS} -- (shift.west);
  \node[below=0.35cm of shift, hw label] {SHCP 移位};
  \node[below=0.35cm of latch, hw label] {STCP 锁存};
\end{block}
\diagramnote{74HC595：移位寄存器与输出锁存器分开。串行数据经 DS 在 SHCP 边沿逐位移入，外部输出全程不动；STCP 边沿才把整字节一次性提交到 Q0..Q7。}

这一类器件提前引出后续流水线寄存器的思想。流水线寄存器保存的不是单个数，而是一组需要一起前进的状态：数据、控制位、valid 位、异常标志和目的寄存器编号。若只有数据向前移动而控制位没有同步移动，后级就会把一个周期的数据解释成另一个周期的操作。移位寄存器是最简单的“状态随边界推进”模型；CPU 流水线是同一思想在更宽状态集合上的应用。

\topic{专题：串并转换与动态扫描显示}

74HC595 的两级边界前面已经画过。把它真正用到板子上，还要把每个引脚的职责和一种典型负载——多位数码管——连起来看，才能理解为什么这颗芯片几乎成了“输出扩展”的代名词。

先看引脚职责。SER（部分手册记作 DS）是串行数据输入，给出当前要移入的那一位；SRCLK（SHCP）是移位时钟，每个有效沿把 SER 的位移入内部移位寄存器，内部状态前进一格，外部引脚纹丝不动；RCLK（STCP）是锁存时钟，有效沿把整个移位寄存器的内容一次性抄进输出锁存器，8 个输出脚同时更新；OE（低有效）是输出使能，决定锁存器内容是否真正驱动引脚，无效时输出呈高阻。此外还有 SRCLR（MR，低有效）异步清空移位寄存器，以及 Q7' 级联输出，把移出的位交给下一片的 SER。

\begin{longtable}{L{0.16\linewidth}L{0.34\linewidth}L{0.40\linewidth}}
\toprule
引脚 & 职责 & 关键问题 \\
\midrule
SER（DS） & 串行数据输入，给出待移入的一位 & 该位必须在 SRCLK 有效沿前稳定 \\
SRCLK（SHCP） & 移位时钟，驱动内部移位寄存器 & 移位期间外部输出不变，这正是设计目的 \\
RCLK（STCP） & 锁存时钟，把整字节提交到输出 & 8 次移位完成后才给，过早会输出半成品字节 \\
OE（低有效） & 输出使能，无效时输出高阻 & 上电期间保持无效，避免显示乱闪 \\
Q7' & 最高位移出，级联下一片 SER & 多片串联时移位次数按总位数计 \\
\bottomrule
\end{longtable}

两级寄存器为什么输出不抖，值得再说透一层。移位过程必然要经历 8 个中间状态，若这些中间状态直接驱动 LED 或数码管，用户会看到一片乱闪。595 把“内部状态推进”和“外部可见提交”拆到两个时钟沿上：SRCLK 管内部怎么动，RCLK 管外部什么时候看。外部负载在任何时刻看到的要么是完整的旧字节，要么是完整的新字节，绝不会看到移了一半的混合字节——这就是“寄存器是原子提交”在接口芯片上的一次应用。

把 8 次移位加一次锁存画到同一条时间轴上，两个时钟各管一段的边界就清楚了：

\begin{waveform}[x=1.0cm,y=1cm]
  % 参考虚线：第一个移位沿与 RCLK 提交沿
  \draw[BookMuted, dashed, line width=0.45pt] (0,-0.6) -- (0,4.6);
  \draw[BookMuted, dashed, line width=0.45pt] (8.2,-0.6) -- (8.2,4.6);
  \node[hw label, anchor=west] at (0.08,4.72) {第 1 个移位沿};
  \node[hw label, anchor=east] at (8.12,4.72) {8 次移位完成，RCLK 提交};
  % SRCLK：8 个移位脉冲
  \draw[BookAccent, line width=1.0pt] (-0.5,4.0) -- (0,4.0) \foreach \i in {0,...,7} { -- (\i,4.32) -- ({\i+0.5},4.32) -- ({\i+0.5},4.0) -- ({\i+1},4.0) } -- (9.5,4.0);
  \node[hw label, anchor=east] at (-0.7,4.16) {SRCLK};
  % SER：逐拍给出 d7..d0
  \foreach \i/\b in {0/7,1/6,2/5,3/4,4/3,5/2,6/1,7/0} {
    \draw[draw=BookTeal, fill=BookCodeBg, line width=0.6pt] ({\i-0.45},2.75) rectangle ({\i+0.55},3.25);
    \node[hw label] at ({\i+0.05},3.0) {$d_{\b}$};
  }
  \node[hw label, anchor=east] at (-0.7,3.0) {SER};
  % RCLK：一次锁存脉冲
  \draw[BookAccent, line width=1.0pt] (-0.5,2.0) -- (8.2,2.0) -- (8.2,2.32) -- (8.7,2.32) -- (8.7,2.0) -- (9.5,2.0);
  \node[hw label, anchor=east] at (-0.7,2.16) {RCLK};
  % OE（低有效）：提交后拉低
  \draw[BookAmber, line width=1.0pt] (-0.5,1.32) -- (8.9,1.32) -- (8.9,1.0) -- (9.5,1.0);
  \node[hw label, anchor=east] at (-0.7,1.16) {OE};
  % 输出锁存器内容
  \draw[draw=BookMuted, fill=BookCodeBg, line width=0.7pt] (-0.5,-0.25) rectangle (8.2,0.25);
  \node[hw label] at (3.85,0) {旧字节（移位期间不变）};
  \draw[draw=BookTeal, fill=BookCodeBg, line width=0.7pt] (8.2,-0.25) rectangle (9.5,0.25);
  \node[hw label, anchor=south] at (8.85,0.32) {新字节};
  \node[hw label, anchor=east] at (-0.7,0) {Q0..Q7};
  \node[hw label, text=BookMuted] at (4.7,-0.62) {OE 低有效：拉低后锁存内容才驱动引脚};
\end{waveform}
\diagramnote{74HC595 串并转换时序（括号内为部分手册别名：SER=DS、SRCLK=SHCP、RCLK=STCP）：SRCLK 的 8 个有效沿把 SER 上的 $d_7$ 到 $d_0$ 逐位移入内部移位寄存器，期间输出锁存器保持旧字节，外部引脚纹丝不动；8 次移位完成后给一个 RCLK 沿，整字节一次性提交到 Q0..Q7；OE 低有效，拉低之前引脚呈高阻。内部状态推进与外部可见提交分开，正是这颗芯片输出不抖的原因。}

最典型的负载是多位数码管。N 位数码管通常段线并联、位线分开：所有位的 a 段接同一根线，每位的公共端各有位选。要显示不同的数字，只能分时复用：每一小段时间只选通一位，同时段码线送出这一位的图案；下一小段时间选通下一位，段码同步换成下一位的图案。只要整帧刷新率超过约 50Hz（工程上常取 100Hz 以上），人眼视觉暂留会把轮流点亮看成同时常亮。

数字算一遍：4 位数码管，每位点亮 2ms，一帧 8ms，帧率 125Hz，每位每秒被刷新 125 次，远在闪烁阈值之上。代价是每位占空比只有 1/4，亮度随之下降，需要在器件允许的脉冲电流范围内提高段电流来补偿。扫描顺序上有一条经验：先消隐、再换码、后点亮——若先开下一位再改段码，上一位的段码会短暂串到新位上，形成鬼影。用 595 驱动段码时这条天然成立：新段码在移位期间被锁存器挡在门内，移完后一拍 RCLK 原子切换，剩下的只是位选时序。

\begin{longtable}{L{0.16\linewidth}L{0.36\linewidth}L{0.36\linewidth}}
\toprule
扫描步骤 & 动作 & 检查点 \\
\midrule
消隐 & 关闭当前位选（或令 OE 无效） & 位选与段码不能同时切换 \\
送段码 & 段码线（或 595）给出下一位图案 & 段码必须在位选打开前稳定 \\
选通 & 打开下一位位选 & 同一时刻只有一位被选通 \\
保持 & 维持点亮约 1 至 2ms & 整帧周期小于 20ms，即帧率大于 50Hz \\
循环 & 推进到下一位，周而复始 & 扫描由定时器驱动，不能被长任务卡住 \\
\bottomrule
\end{longtable}

串并转换和动态扫描合起来，说明同一个原则：外部接口的引脚永远比内部状态少，解决办法是把“哪些位先走、哪些位后走”变成明确的时间约定，并用寄存器边界保证每一拍对外都是完整的。

把本节几种结构放在一起看，各自的设计要点不同：普通寄存器保存 ALU 结果、地址和控制字，要写使能只在目标周期打开；计数器用于 PC 加一、循环计数和超时计数，要使能、清零、装载优先级明确；移位寄存器用于串并转换、延迟线和输出扩展，要移位边界和输出锁存边界分开；标志寄存器保存 zero/carry/error 等条件，要说明哪些指令或事件会覆盖标志。

\section{五、状态机把硬件动作切成阶段}

有限状态机由三部分组成：当前状态、输入条件、下一状态逻辑。当前状态保存在寄存器里，组合逻辑根据当前状态和输入条件决定下一状态。

\begin{block}[x=1cm,y=1cm]
  \node[hw state, minimum width=2.25cm, minimum height=0.9cm] (state) at (0,0.7) {状态寄存器};
  \node[hw compute, minimum width=2.75cm, minimum height=0.9cm] (next) at (3.5,0.7) {下一状态逻辑};
  \node[hw control, minimum width=2.25cm, minimum height=0.9cm] (out) at (7.0,0.7) {输出逻辑};
  \node[hw lane, minimum width=2.0cm] (inputs) at (3.5,-1.15) {输入};
  \node[hw lane, minimum width=1.55cm] (clk) at (-2.2,0.7) {CLK};
  \node[hw lane, minimum width=1.8cm] (outputs) at (9.6,0.7) {控制输出};
  \draw[hw bus] (clk) -- (state);
  \draw[hw bus] (state) -- (next);
  \node[hw label] at (1.75,1.08) {当前状态};
  \draw[hw return] (next.north) |- (1.75,1.55) -| (state.north);
  \node[hw label] at (1.75,1.78) {next state};
  \draw[hw bus] (state.south) |- (1.15,-0.42) -- (5.95,-0.42) -| (out.south);
  \node[hw label] at (5.45,-0.62) {当前状态};
  \draw[hw bus] (inputs.north) -- (next.south);
  \draw[hw bus] (inputs.east) -| (out.south);
  \draw[hw bus] (out) -- (outputs);
  \node[hw label, text width=9.4cm] at (3.8,-2.0) {状态寄存器只在时钟边界提交；下一状态逻辑和输出逻辑是组合网络。图中把反馈放在上方直线，避免弯箭头穿过节点。};
\end{block}
\diagramnote{有限状态机结构：状态寄存器 + 下一状态逻辑 + 输出逻辑。}


设计状态机时，第一步不是写状态名字，而是列出机器必须区分的阶段。若一个控制器要等待请求、装载数据、运行运算、等待外设完成、提交结果，那么这些阶段是否能合并，取决于它们是否需要不同的控制输出和不同的等待条件。第二步才是给阶段编码，并写出每个状态在每种关键输入下的下一状态。第三步是列出输出：哪些输出只由状态决定，哪些输出还依赖输入，哪些输出必须寄存后再送往外部。

一个完整状态机说明至少包含五张表。第一张是状态含义表，说明每个状态代表哪个硬件阶段——状态命名漂亮但动作不明确，等于没有状态。第二张是转移表，说明当前状态遇到关键输入后进入哪里，等待、错误和默认分支都不能遗漏。第三张是输出表，说明每个状态打开哪些控制线，写使能、片选、ALU 选择、驱动使能不能在错误周期打开。第四张是复位和非法状态表，说明上电、复位释放和异常编码如何处理，入口状态、默认安全输出和恢复路径缺一不可。第五张是波形检查表，说明哪些信号必须在同一周期出现——最终结果正确但中间周期错误，同样是失败。缺少任何一张，状态机都容易退化成“名字看起来对，但边界不清楚”的草率描述。

以一个三阶段硬件控制器为例：

\begin{longtable}{L{0.2\linewidth}L{0.32\linewidth}L{0.34\linewidth}}
\toprule
当前状态 & 条件 & 下一状态 \\
\midrule
IDLE & start=1 & LOAD \\
LOAD & ready=1 & RUN \\
RUN & count\_done=1 & DONE \\
DONE & clear=1 & IDLE \\
\bottomrule
\end{longtable}

若当前状态是 IDLE 且 start=0，状态保持 IDLE；若 start=1，下一状态逻辑准备 LOAD，等时钟沿到来后状态寄存器才真正变成 LOAD。状态机因此不会在一个周期里连续越过 IDLE、LOAD、RUN、DONE，它每次只跨过一个明确边界。

上表还不够完整，因为它只写了“去哪里”，没有写“做什么”。若 LOAD 状态要把外部数据装入寄存器，就必须说明 \code{load_reg_we} 在哪个周期为 1；若 RUN 状态要让计数器递减，就必须说明计数使能、ALU 选择和完成条件；若 DONE 状态要向外部报告结果，就必须说明输出是否寄存、ready 信号保持多久、clear 到来前是否允许覆盖结果。状态机设计的核心不是列出状态名称，而是让状态名称、控制输出和被控制的状态元件逐周期对应。

\begin{longtable}{L{0.16\linewidth}L{0.24\linewidth}L{0.26\linewidth}L{0.24\linewidth}}
\toprule
状态 & 本周期动作 & 关键输出 & 离开条件 \\
\midrule
IDLE & 等待同步后的启动请求 & 关闭写使能与外部驱动 & \code{start_pending=1} \\
LOAD & 装载输入寄存器或初始计数 & \code{load_we=1} & 输入 ready 或装载完成 \\
RUN & 执行运算、计数或移位 & \code{run_en=1} & \code{count_done=1} \\
DONE & 保持结果并等待清除 & \code{done=1}，输出稳定 & \code{clear=1} \\
\bottomrule
\end{longtable}

\begin{block}[x=1cm,y=1cm]
  \node[hw state, minimum width=1.6cm, minimum height=0.9cm] (idle) at (0,0) {IDLE};
  \node[hw state, minimum width=1.6cm, minimum height=0.9cm] (load) at (2.8,0) {LOAD};
  \node[hw state, minimum width=1.6cm, minimum height=0.9cm] (run) at (5.6,0) {RUN};
  \node[hw state, minimum width=1.6cm, minimum height=0.9cm] (done) at (8.4,0) {DONE};
  \draw[hw bus] (idle) -- node[above, hw label, pos=0.58, yshift=3pt]{start=1} (load);
  \draw[hw bus] (load) -- node[above, hw label]{ready=1} (run);
  \draw[hw bus] (run) -- node[above, hw label]{count\_done=1} (done);
  \draw[hw return] (done.south) |- (4.2,-1.45) -| (idle.south) node[pos=0.23, below, hw label]{clear=1};
  \draw[BookMuted,line width=0.55pt,-{Stealth[length=2mm,width=1.4mm]}] (idle.north) .. controls +(0,0.72) and +(0.72,0.72) .. node[above,hw label]{start=0} (idle.east);
  \draw[BookMuted,line width=0.55pt,-{Stealth[length=2mm,width=1.4mm]}] (done.north) .. controls +(0,0.72) and +(0.72,0.72) .. node[above,hw label]{clear=0} (done.east);
  \node[hw label, text width=8.8cm] at (4.2,-2.15) {主路径水平推进，返回路径独占下方通道，自环标签放在节点上方，避免和返回线相交。};
\end{block}
\diagramnote{四状态硬件控制器状态图：IDLE 到 LOAD 到 RUN 到 DONE 再回到 IDLE。每个保持条件画成自环——状态机不是只会前进，更多时候是在等条件。}

把这个状态机真正做完一次。第二步是编码：四个状态用两个 bit，\code{IDLE=00}、\code{LOAD=01}、\code{RUN=10}、\code{DONE=11}，编码刚好用完，没有非法空洞。第三步写下一状态逻辑和输出逻辑：

\begin{lstlisting}
case (S):
  IDLE: next = start      ? LOAD : IDLE
  LOAD: next = ready      ? RUN  : LOAD
  RUN:  next = count_done ? DONE : RUN
  DONE: next = clear      ? IDLE : DONE
  default: next = IDLE

load_we = (S == LOAD)
run_en  = (S == RUN)
done    = (S == DONE)
\end{lstlisting}

写成布尔式也能机械推导。以 \code{next\_LOAD} 为例：进入 LOAD 的条件是“当前 IDLE 且 start=1”，留在 LOAD 的条件是“当前 LOAD 且 ready=0”：

\begin{lstlisting}
next_LOAD = (S == IDLE AND start) OR (S == LOAD AND NOT ready)
\end{lstlisting}

第四张表是复位：上电和复位期间状态寄存器强制为 \code{IDLE=00}，\code{load\_we/run\_en/done} 全为 0。第五张表是波形检查：把一次完整运行按时钟对齐画出来，检查每个控制线是否只在规定周期有效。

\begin{waveform}[x=1cm,y=1cm]
  \draw[BookAccent, line width=0.9pt] (0,7.6) -- (2,7.6) -- (2,8.4) -- (3,8.4) -- (3,7.6) -- (4,7.6) -- (4,8.4) -- (5,8.4) -- (5,7.6) -- (6,7.6) -- (6,8.4) -- (7,8.4) -- (7,7.6) -- (8,7.6) -- (8,8.4) -- (9,8.4) -- (9,7.6) -- (10,7.6) -- (10,8.4) -- (11,8.4);
  \node[hw label, anchor=east] at (-0.2,8.0) {CLK};
  \draw[BookAmber, line width=0.9pt] (0,6.4) -- (1.5,6.4) -- (1.5,7.2) -- (2.5,7.2) -- (2.5,6.4) -- (11,6.4);
  \node[hw label, anchor=east] at (-0.2,6.8) {start};
  \draw[BookAmber, line width=0.9pt] (0,5.4) -- (3.5,5.4) -- (3.5,6.2) -- (4.5,6.2) -- (4.5,5.4) -- (11,5.4);
  \node[hw label, anchor=east] at (-0.2,5.8) {ready};
  \draw[BookAmber, line width=0.9pt] (0,4.4) -- (5.5,4.4) -- (5.5,5.2) -- (6.5,5.2) -- (6.5,4.4) -- (11,4.4);
  \node[hw label, anchor=east] at (-0.2,4.8) {\scriptsize count\_done};
  \draw[BookAmber, line width=0.9pt] (0,3.4) -- (7.5,3.4) -- (7.5,4.2) -- (8.5,4.2) -- (8.5,3.4) -- (11,3.4);
  \node[hw label, anchor=east] at (-0.2,3.8) {clear};
  \draw[BookTeal, line width=0.9pt] (0,2.0) -- (10,2.0) -- (10,2.8) -- (0,2.8) -- cycle;
  \foreach \x in {2,4,6,8} { \draw[BookTeal, line width=0.9pt] (\x,2.0) -- (\x,2.8); }
  \node at (1.0,2.4) {00};
  \node at (3.0,2.4) {01};
  \node at (5.0,2.4) {10};
  \node at (7.0,2.4) {11};
  \node at (9.0,2.4) {00};
  \node[hw label, anchor=east] at (-0.2,2.4) {S[1:0]};
  \draw[BookAccent, line width=0.9pt] (0,0.8) -- (2,0.8) -- (2,1.6) -- (4,1.6) -- (4,0.8) -- (11,0.8);
  \node[hw label, anchor=east] at (-0.2,1.2) {load\_we};
  \draw[BookAccent, line width=0.9pt] (0,-0.4) -- (4,-0.4) -- (4,0.4) -- (6,0.4) -- (6,-0.4) -- (11,-0.4);
  \node[hw label, anchor=east] at (-0.2,0.0) {run\_en};
  \draw[BookMuted, dashed, line width=0.4pt] (2,-0.7) -- (2,8.6);
  \draw[BookMuted, dashed, line width=0.4pt] (4,-0.7) -- (4,8.6);
  \draw[BookMuted, dashed, line width=0.4pt] (6,-0.7) -- (6,8.6);
  \draw[BookMuted, dashed, line width=0.4pt] (8,-0.7) -- (8,8.6);
\end{waveform}
\diagramnote{一次完整运行的对齐波形：输入条件在采样沿前后保持稳定，状态只在上升沿跳转；\code{load\_we} 只在 LOAD 的那一个周期为高，\code{run\_en} 也只在 RUN 的那一个周期为高。检查一台状态机，就是对照这张图逐周期核对：控制线有没有早一拍、晚一拍或粘在错误的周期。}

状态机输出可以只由当前状态决定，也可以同时依赖输入。前者输出更稳定，后者反应更快但更容易受输入毛刺影响。基本原则是：状态机不是若干状态名称的集合，而是“状态寄存器 + 下一状态组合逻辑 + 输出组合逻辑”的硬件结构。

Moore 型输出只由当前状态决定，常用于写使能、片选、锁存时钟、功率使能等需要稳定边界的控制线。Mealy 型输出同时依赖状态和输入，响应可能少等一个周期，但输入毛刺也可能被带到控制线。工程上并不是简单规定只能用哪一种，而是按后果分类：危险输出、外部驱动和写存储器信号倾向于寄存或由稳定状态产生；纯内部、已同步且后级还会采样的条件信号，可以在确认时序余量后使用组合依赖。

\topic{专题：Moore 与 Mealy 输出对照}

上一段把输出分成“只由状态决定”和“同时依赖输入”两类，并给出了按后果分类的原则。这里用同一个功能把两类做实：检测已同步输入 \code{sig} 的上升沿，发现后输出一个脉冲 \code{pulse}。

Moore 型实现需要三个状态。\code{S_LOW} 表示上一拍 sig 为 0，\code{S_EDGE} 表示刚采到上升沿（pulse 只在此状态为 1），\code{S_HIGH} 表示 sig 已稳定为 1、等待回落：

\begin{longtable}{L{0.16\linewidth}L{0.28\linewidth}L{0.16\linewidth}L{0.14\linewidth}L{0.14\linewidth}}
\toprule
当前状态 & 含义 & pulse 输出 & sig=1 去向 & sig=0 去向 \\
\midrule
\code{S_LOW} & 上一拍 sig 为 0 & 0 & \code{S_EDGE} & \code{S_LOW} \\
\code{S_EDGE} & 刚采到上升沿 & 1 & \code{S_HIGH} & \code{S_LOW} \\
\code{S_HIGH} & sig 已稳定为 1 & 0 & \code{S_HIGH} & \code{S_LOW} \\
\bottomrule
\end{longtable}

Mealy 型实现只用两个状态，把“是否刚采到边沿”交给输入参与判断：pulse 在“当前状态为 \code{S_LOW} 且 sig=1”的那个周期直接为 1。

\begin{longtable}{L{0.20\linewidth}L{0.14\linewidth}L{0.26\linewidth}L{0.28\linewidth}}
\toprule
当前状态 & sig & pulse（本拍） & 下一状态 \\
\midrule
\code{S_LOW} & 0 & 0 & \code{S_LOW} \\
\code{S_LOW} & 1 & 1 & \code{S_HIGH} \\
\code{S_HIGH} & 0 & 0 & \code{S_LOW} \\
\code{S_HIGH} & 1 & 0 & \code{S_HIGH} \\
\bottomrule
\end{longtable}

对照两种实现的时序。Moore 版在采到边沿的下一拍才进入 \code{S_EDGE}，脉冲晚一个周期出现，但宽度恰为一个完整时钟周期，且输出只是状态译码，触发器输出之后最多一级门，干净稳定。Mealy 版的 \code{pulse = (state == S\_LOW) AND sig}，在采到边沿的同一个周期内就升高，抢回了一个周期；但这条输出路径上有输入 sig 的组合分量：若 sig 未经同步、带有毛刺，毛刺会原样穿到 pulse 上；即使 sig 本身已是寄存器输出，pulse 也要等沿后 $t_{CQ}$ 加一级门延迟才稳定，脉冲的头尾不像 Moore 版那样整齐。

把同一个 \code{sig} 下的两路输出画在一起，当拍与晚一拍、毛刺与干净的差别直接可见：

\begin{waveform}[x=1.6cm,y=1cm]
  % 参考虚线：时钟沿 1/2/3
  \foreach \x in {1,2,3} {\draw[BookMuted, dashed, line width=0.45pt] (\x,0.55) -- (\x,4.6);}
  % CLK
  \draw[BookAccent, line width=1.0pt] (-0.5,4.0) -- (0,4.0) \foreach \i in {0,...,4} { -- (\i,4.32) -- ({\i+0.5},4.32) -- ({\i+0.5},4.0) -- ({\i+1},4.0) };
  \node[hw label, anchor=east] at (-0.7,4.16) {CLK};
  % sig：沿前窄毛刺 + 沿 1 后升高
  \draw[BookAmber, line width=1.0pt] (-0.5,3.0) -- (0.42,3.0) -- (0.42,3.32) -- (0.62,3.32) -- (0.62,3.0) -- (1.08,3.0) -- (1.08,3.32) -- (4.08,3.32) -- (4.08,3.0) -- (5.0,3.0);
  \node[hw label, anchor=east] at (-0.7,3.16) {sig};
  % Moore pulse：沿 2 后一整拍，干净
  \draw[BookTeal, line width=1.1pt] (-0.5,1.9) -- (2.08,1.9) -- (2.08,2.22) -- (3.08,2.22) -- (3.08,1.9) -- (5.0,1.9);
  \node[hw label, anchor=east] at (-0.7,2.06) {Moore pulse};
  \node[hw label, text=BookTeal, anchor=west] at (3.25,2.14) {晚一拍，完整一拍};
  % Mealy pulse：毛刺穿透 + 当拍响应（头尾参差）
  \draw[BookRed, line width=1.1pt] (-0.5,0.9) -- (0.47,0.9) -- (0.47,1.22) -- (0.67,1.22) -- (0.67,0.9) -- (1.16,0.9) -- (1.16,1.22) -- (2.16,1.22) -- (2.16,0.9) -- (5.0,0.9);
  \node[hw label, anchor=east] at (-0.7,1.06) {Mealy pulse};
  \node[hw label, text=BookRed, anchor=east] at (0.32,0.62) {毛刺穿透};
  \node[hw label, text=BookRed] at (1.66,0.32) {当拍响应，头尾随 $t_{CQ}$ 加门延迟};
\end{waveform}
\diagramnote{同一 \code{sig} 下的两种输出：Mealy 在采到上升沿的当拍拉高，抢回一个周期，但 \code{sig} 上的窄毛刺原样穿透（左侧红色小脉冲），且脉冲头尾要等沿后 $t_{CQ}$ 加一级门延迟，不如 Moore 整齐；Moore 晚一个周期进入 \code{S\_EDGE}，输出恰为整一拍宽的状态译码，毛刺无从穿透。驱动外部器件与写使能选 Moore 或再寄存一拍，同域内部条件信号可用 Mealy。}

\begin{longtable}{L{0.20\linewidth}L{0.34\linewidth}L{0.34\linewidth}}
\toprule
维度 & Moore & Mealy \\
\midrule
本例状态数 & 3 & 2 \\
输出决定因素 & 仅当前状态 & 当前状态加输入 \\
响应时刻 & 采到条件后下一拍 & 采到条件的当拍 \\
输出稳定度 & 整拍稳定，沿后至多一级门 & 受输入路径影响，沿后 $t_{CQ}$ 加组合延迟 \\
毛刺风险 & 输出无输入直接路径 & 输入毛刺可穿透到输出 \\
适用输出 & 写使能、片选、外部驱动 & 同域内部、后级会再采样的条件信号 \\
\bottomrule
\end{longtable}

选择原则于是回到前文的按后果分类。输出要驱动外部器件、写存储器、控制功率，选 Moore，或者在 Mealy 输出后再加一级寄存器——后者把 Mealy 的早判断和 Moore 的干净边界拼在一起，代价是又回到晚一拍。输出只在同一时钟域内部被后级再次采样、且功能上必须抢那一拍，Mealy 可以放心用，前提是输入已经同步、组合路径有时序余量。状态数上 Mealy 常常更省，因为它把一部分状态信息转交给了输入条件；但省下的触发器换来的是输出上的组合路径，这笔交换必须用上面这张对照表逐项想清楚。

一个状态表必须覆盖默认情况。若状态是 \code{IDLE}、\code{LOAD}、\code{RUN}、\code{DONE}，而编码使用两个 bit，则四个编码刚好用完；若状态更多或编码留有空洞，未使用编码不能留给综合器或门延迟“自己决定”。安全写法是给下一状态逻辑设置默认恢复路径，例如进入 \code{FAULT}、\code{RESET_WAIT} 或 \code{IDLE}，同时让写使能、外部驱动、片选和功率使能默认为关闭。这样，即使上电扰动、复位释放问题或外部干扰让状态寄存器进入非法编码，电路也不会在非法状态下驱动危险动作。

把状态机连接到后续 CPU 时，可以把每条控制线都问成一个状态问题：PC 什么时候写入，指令寄存器什么时候装载，寄存器堆写使能在哪个周期打开，ALU 选择哪种运算，存储器读写请求何时有效，外设未 ready 时机器停在哪个状态。若这些问题只用“连线会变成正确值”回答，就没有真正形成控制器；只有把它们逐一对应到状态、条件、下一状态和输出时序上，CPU 才能按周期推进。

以最小多周期 CPU 为例，状态机通常至少要区分取指、译码、执行、访存和写回。取指状态打开存储器读请求 \code{mem_read}，选择 \code{pc_to_addr}，并在存储器 ready 后装载指令寄存器（\code{ir_we}）；译码状态读取寄存器堆并准备立即数，但不写回任何目的寄存器；执行状态选择 ALU 运算 \code{alu_op} 或计算地址，分支条件和 ALU 结果必须在采样前稳定；访存状态等待数据存储器 ready，期间保持地址和控制线不变；写回状态打开寄存器堆写使能 \code{rf\_we}，且只在目标指令的目标寄存器上打开。若存储器没有 ready，状态机不能假设数据已经存在，而应保持在等待状态并关闭不该打开的写使能。这样，CPU 的每个周期都遵循本章模型：当前状态和寄存器输出经过组合控制，形成下一状态和下一批写入。

这种写法也能解释为什么现代芯片会引入更复杂的控制状态和流水阶段。它们不是为了把教材概念复杂化，而是为了缩短关键路径、隐藏等待、保持高吞吐，并让异常、中断、cache miss、分支错误等事件有可恢复边界。无论是 8086 的总线等待状态、80386 的系统控制状态，还是现代 CPU 的提交队列和异常入口，本质上都要回答同一个问题：某个动作在哪个状态被允许，结果在哪个边界对外可见，失败时回到哪个安全状态。

现在硬件已经能够保存一组 bit、按周期加一、按周期移动、按条件切换阶段。下一步要把这些能力用于更大的动作组织：哪些寄存器在这个周期写入，ALU 做哪种运算，哪一路数据被选择，下一阶段去哪里。这就是控制器的基础。

\topic{专题：时钟、复位、状态编码与验证闭环}

前面已经把同步边界讲成寄存器到寄存器的路径，但真实工程中还必须处理四个经常被低估的问题：时钟如何送达，复位如何跨越边界，状态编码如何影响安全性，验证如何证明边界确实成立。这四件事不改变“组合逻辑准备下一值、状态元件保存下一值”的基本模型，却决定一个设计能否从纸面状态表进入可靠机器。

首先区分\hwkey{写使能}和\hwrisk{门控时钟}。写使能的做法是让时钟继续到达触发器，但在时钟沿到来时决定是否接收新 D 值；若写使能为 0，寄存器保持原值。门控时钟则是让某段时钟停止翻转，以降低功耗或隔离无用活动。两者在功能描述上都可能表现为“保持不变”，但硬件风险完全不同。普通数据通路优先使用写使能；真正需要门控时钟时，应使用专用 clock gating 单元，并保证 enable 在安全窗口内被锁存。不能随意用一个 AND 门把 \code{clk} 和 \code{enable} 相与后送给触发器，因为 enable 的毛刺或变化时刻可能制造短脉冲、窄时钟或额外时钟沿。

把几种保持状态的写法排开看：写使能让时钟照常到达，触发器按使能选择写入或保持，使能信号必须满足目标触发器时序；数据 mux 保持用 \code{D=mux(we,new,Q)} 写入旧值实现保持，代价是 mux 延迟进入数据关键路径；专用时钟门控用门控单元生成局部时钟，要保证 enable 被锁存、无毛刺、满足时钟树约束；普通门电路直接切时钟最容易产生短脉冲和不可控偏斜，通常禁止使用。

其次，复位也有“域”。

一个时钟域内部可以规定复位异步拉入、同步释放；两个时钟域之间则不能假设复位释放同时发生。若控制状态机已经退出复位，而它依赖的外设状态机、同步器或计数器仍在复位中，就可能读到无效状态。复位域设计应回答三个问题：哪些状态元件属于同一复位域，哪些跨域信号在复位期间必须保持安全，复位释放后第一个有效事件如何产生。

按对象写复位策略：控制状态要复位到 \code{RESET_WAIT} 或 \code{IDLE}，上电后进入未定义状态等于没有控制器；数据寄存器只复位影响控制或外部可见安全的那些，大面积无意义复位只会增加布线和时序压力；跨域复位释放要在每个时钟域内部同步释放，并等待依赖域 ready，否则一个域先运行、另一个域仍无效；复位后的第一个事件要特别小心，同步器和 pending 位要进入不会伪触发的默认值，否则复位释放本身就会产生假启动或假中断。

第三，状态机编码不是纯粹的排版选择。二进制编码使用较少触发器，但下一状态译码可能较复杂；one-hot 编码为每个状态使用一个触发器，译码直接、常适合较高速度或 FPGA 结构，但触发器数量更多；格雷编码让相邻状态尽量只改变一位，常用于跨域计数指针或减少多位同时变化带来的采样风险。无论使用哪种编码，都要定义未使用编码的下一状态和输出默认值。安全相关控制器还应考虑显式安全编码：为非法编码、错误记录和恢复路径预留设计，而不是指望编码本身永远不出错。

第四，验证必须看波形上的边界。只看最终寄存器值，可能错过“某个周期写使能提前打开”“复位释放产生伪事件”“同步器第一级短暂不稳定”“pending 位被错误清除”等问题。一个同步小系统的波形检查至少要同时观察 \code{clk}、\code{reset}、源寄存器 Q、目标寄存器 D、目标写使能、状态编码、同步器两级输出、事件 pending 位和 ack 信号。

每个对象要检查的问题不同：\code{reset} 的释放是否在本时钟域边界生效，否则状态机会半周期启动或产生伪事件；\code{Q -> D} 路径上 D 是否在采样沿前稳定，否则目标寄存器采到过渡值；\code{write_enable} 是否只在允许写入的周期有效，提前一拍就是寄存器被提前覆盖；同步器的第二级输出是否作为唯一后级输入，后级直接使用 raw 或 meta 信号等于没有同步；\code{pending/ack} 的事件是否保持到被消费，窄事件丢失、重复消费和错误清除都是真实失败；非法状态编码是否关闭危险输出并恢复，未使用状态不能驱动写使能或外部输出。

这些检查会直接连接到后续 CPU 章节。PC 是寄存器，指令寄存器是状态，流水线寄存器是阶段边界，控制器是状态机，异常和中断入口依赖同步事件，写回使能必须在正确周期打开。若第二章的边界概念不清楚，后面的数据通路和最小 CPU 就会退化成“线连起来就能跑”的误解。真正的 CPU 不是一张组合逻辑图，而是一组由时钟、复位、状态编码、控制使能和波形验证共同约束的同步系统。

\topic{经典案例：交通灯控制器的完整设计}

前面的状态机例子只有四五个抽象状态，现在用一个带真实时间需求的案例，把“需求、状态列表、转移表、输出表”完整走一遍。设计一个十字路口的交通灯控制器：东西向绿灯亮 30 秒，转黄灯 5 秒，然后双向全红 2 秒；接着南北向绿灯 30 秒、黄灯 5 秒，再全红 2 秒，随后回到东西向绿灯，如此循环。全红段不是装饰：在一个方向放行之前，两个方向都必须先看到红灯，给已经进入路口的车辆留出清空时间。这 2 秒是用一个专门状态换来的安全余量，删掉它，状态机更简单，路口更危险。

设计的第一步是把时间从状态机里剥离出去。30、5、2 都是计数值，而状态机不应该自己数数。让一个独立的定时器——一个带装载值的递减计数器——负责计数：状态机进入每个状态时，把该阶段的秒数装载进去并启动定时器；定时器减到 0，发出一个周期宽的 \code{timer\_done} 信号；状态机看到它就走向下一阶段，并装载下一阶段的秒数。分工之后，状态机只回答“现在在哪个阶段、下一阶段是什么”，定时器只回答“这个阶段结束了没有”。将来把 30 秒改成 40 秒，只动装载值，状态机结构一行不改。若系统时钟是 1Hz，定时器就是一个最大装载 30 的 5-bit 递减计数器；若系统时钟更快，先分频得到 1Hz 使能，定时器每秒钟才减一次。

状态列表按灯色阶段划分，共六个状态。注意两个全红状态不能合并：它们的灯色输出完全相同，但下一状态不同——一个之后放行南北向，一个之后放行东西向。状态划分由“未来行为是否相同”决定，而不是由“当前输出是否相同”决定；输出相同而去向不同的两个阶段，必须是两个状态。

\begin{longtable}{L{0.22\linewidth}L{0.5\linewidth}L{0.16\linewidth}}
\toprule
状态 & 含义 & 装载值 \\
\midrule
\code{EW\_GREEN} & 东西向绿灯，南北向红灯 & 30s \\
\code{EW\_YELLOW} & 东西向黄灯，南北向红灯 & 5s \\
\code{ALL\_RED\_EW} & 东西向刚结束，双向全红 & 2s \\
\code{NS\_GREEN} & 南北向绿灯，东西向红灯 & 30s \\
\code{NS\_YELLOW} & 南北向黄灯，东西向红灯 & 5s \\
\code{ALL\_RED\_NS} & 南北向刚结束，双向全红 & 2s \\
\bottomrule
\end{longtable}

转移表只有一类条件：\code{timer\_done} 是否为 1。为 0 时保持原状态，下表只列出为 1 时的转移行。

\begin{longtable}{L{0.2\linewidth}L{0.2\linewidth}L{0.2\linewidth}L{0.28\linewidth}}
\toprule
当前状态 & 条件 & 下一状态 & 动作 \\
\midrule
\code{EW\_GREEN} & \code{timer\_done}=1 & \code{EW\_YELLOW} & 装载 5s，重启定时器 \\
\code{EW\_YELLOW} & \code{timer\_done}=1 & \code{ALL\_RED\_EW} & 装载 2s，重启定时器 \\
\code{ALL\_RED\_EW} & \code{timer\_done}=1 & \code{NS\_GREEN} & 装载 30s，重启定时器 \\
\code{NS\_GREEN} & \code{timer\_done}=1 & \code{NS\_YELLOW} & 装载 5s，重启定时器 \\
\code{NS\_YELLOW} & \code{timer\_done}=1 & \code{ALL\_RED\_NS} & 装载 2s，重启定时器 \\
\code{ALL\_RED\_NS} & \code{timer\_done}=1 & \code{EW\_GREEN} & 装载 30s，重启定时器 \\
\bottomrule
\end{longtable}

输出是典型 Moore 型：灯色只由当前状态决定，与 \code{timer\_done} 无关，因此在整个状态期间稳定，输入信号上的毛刺不可能让灯闪一下。

\begin{longtable}{L{0.28\linewidth}L{0.3\linewidth}L{0.3\linewidth}}
\toprule
状态 & 东西向 & 南北向 \\
\midrule
\code{EW\_GREEN} & 绿 & 红 \\
\code{EW\_YELLOW} & 黄 & 红 \\
\code{ALL\_RED\_EW} & 红 & 红 \\
\code{NS\_GREEN} & 红 & 绿 \\
\code{NS\_YELLOW} & 红 & 黄 \\
\code{ALL\_RED\_NS} & 红 & 红 \\
\bottomrule
\end{longtable}

还有两处工程细节值得写明。第一，复位状态应选全红状态而不是某个绿灯状态：上电后的第一件事是让所有方向看到红灯，而不是让某个方向侥幸先走。第二，一个完整循环是 74 秒，但状态机的硬件成本与秒数无关——计时全部由定时器承担。这个案例再次印证本章模型：定时器是“寄存器加递减组合逻辑的反馈”，状态机是“寄存器加下一状态逻辑和输出逻辑”，两者靠 \code{timer\_done} 和装载、启动三条线连接，所有变化都发生在时钟边界上。

\topic{专题：状态编码的三种选择}

交通灯控制器有六个状态，状态寄存器该用几个触发器、每个状态编成什么二进制值，并不是排版偏好。编码直接决定三件事：触发器用掉多少、下一状态和输出译码的组合逻辑有多深、状态翻转瞬间译码输出会不会出毛刺。三种经典选择是二进制编码、格雷编码和独热编码，下表给出交通灯六状态的三种编法。

\begin{longtable}{L{0.28\linewidth}L{0.2\linewidth}L{0.2\linewidth}L{0.2\linewidth}}
\toprule
状态 & 二进制（3 位） & 格雷（3 位） & 独热（6 位） \\
\midrule
\code{EW\_GREEN} & 000 & 000 & 000001 \\
\code{EW\_YELLOW} & 001 & 001 & 000010 \\
\code{ALL\_RED\_EW} & 010 & 011 & 000100 \\
\code{NS\_GREEN} & 011 & 010 & 001000 \\
\code{NS\_YELLOW} & 100 & 110 & 010000 \\
\code{ALL\_RED\_NS} & 101 & 100 & 100000 \\
\bottomrule
\end{longtable}

二进制编码按状态顺序编号，用 $\lceil\log_2 6\rceil=3$ 个触发器，最省存储元件；代价是下一状态逻辑和输出译码都要看全部 3 位，逻辑级数最深。格雷编码让循环中相邻的两个状态只差一位：本例刻意选取 000、001、011、010、110、100 这条六状态环，连从末态 100 回到首态 000 的闭环边也只翻最高一位。独热编码为每个状态分配一个触发器，某一位为 1 就表示“正处于这个状态”，译码常常只需看一位，最浅最直接；代价是触发器从 3 个变成 6 个。

\begin{longtable}{L{0.14\linewidth}L{0.12\linewidth}L{0.36\linewidth}L{0.28\linewidth}}
\toprule
编码 & 触发器数 & 译码逻辑 & 适用场合 \\
\midrule
二进制 & 3 & 最深，需综合全部编码位 & 状态很多、触发器昂贵的 ASIC \\
格雷 & 3 & 深度与二进制相当 & 相邻转移无中间码，跨时钟域计数指针常用 \\
独热 & 6 & 最浅，常常一级 & FPGA 触发器丰富，常是默认选择 \\
\bottomrule
\end{longtable}

格雷码的价值要放在“多位同时翻转”的背景下看。状态从 011 变到 100 时三位都要翻，而物理上三位不可能精确同时到达，译码器在翻转窗口里可能瞬间看到 000 或 111；若这个中间码恰好对应某个有效输出，输出线上就冒出一条毛刺。格雷码每次转移只翻一位，不存在中间码，译码窗口天然干净——这也是异步 FIFO 的读写指针跨时钟域传递时一律用格雷码的原因。

独热编码在 FPGA 上流行有结构上的理由：FPGA 的每个逻辑单元都自带触发器，触发器几乎随用随有，而宽输入的译码逻辑要占用查找表层级、拉长关键路径。用 6 个“免费”的触发器换最浅的译码，往往反而更容易满足时序。ASIC 的情形相反，触发器占面积，状态一多就倾向二进制编码。无论选哪种，未使用编码都必须有去处：二进制和格雷各有 110、111 或 101、111 两个空洞，独热更有 $2^6-6=58$ 个非法编码（多位同时为 1 或全为 0），下一状态逻辑应把它们一律导向全红状态并关闭所有绿灯——非法状态恢复是编码设计的收尾，不是附加题。

\topic{经典案例：1101 序列检测器（允许重叠）}

第二类经典状态机不控制外部设备，而是识别数据流：每个时钟沿采样 1 位输入 \code{x}，当最近收到的 4 位恰好是 1101 时，输出 \code{y} 为 1。要求允许重叠：一次匹配的末尾几位如果同时是下一次匹配的开头，必须被复用，不能匹配一次就从头再来。

设计的关键是给每个状态一个明确含义：“到目前为止，输入串的尾巴与模式 1101 的前缀匹配了多长”。状态不需要记住整段历史，只需记住这个长度——未来能否完成匹配，只取决于它。

\begin{longtable}{L{0.16\linewidth}L{0.5\linewidth}L{0.2\linewidth}}
\toprule
状态 & 含义 & 已匹配前缀 \\
\midrule
\code{S0} & 没有可用前缀（起始或刚失配） & （空） \\
\code{S1} & 尾巴是 1 & 1 \\
\code{S2} & 尾巴是 11 & 11 \\
\code{S3} & 尾巴是 110 & 110 \\
\bottomrule
\end{longtable}

采用 Mealy 型，输出标在转移上（写法为“下一状态/输出”）：

\begin{longtable}{L{0.2\linewidth}L{0.34\linewidth}L{0.34\linewidth}}
\toprule
当前状态 & \code{x}=0 & \code{x}=1 \\
\midrule
\code{S0} & \code{S0}/0 & \code{S1}/0 \\
\code{S1} & \code{S0}/0 & \code{S2}/0 \\
\code{S2} & \code{S3}/0 & \code{S2}/0 \\
\code{S3} & \code{S0}/0 & \code{S1}/1 \\
\bottomrule
\end{longtable}

逐行核对这张表，能看出构造方法是完全机械的。在 \code{S2} 收到 1：已匹配串变成 111，它不是 1101 的前缀，但其尾巴 11 仍是最长可用前缀，所以留在 \code{S2}，不退回起点。在 \code{S3} 收到 1：恰好完成 1101，输出 1；同时这一位 1 是下一个匹配的开头，所以退到 \code{S1} 而不是 \code{S0}。在 \code{S3} 收到 0：得到 1100，其尾巴 0、00、100、1100 没有一个是 1101 的前缀，只能回 \code{S0}。一般规则是：失配时把“已匹配前缀加上新输入”的尾巴从长到短逐个试，第一个仍是模式前缀的，就是下一状态。

\begin{lstlisting}
input x : 1  1  0  1  1  0  1
state   : S0 S1 S2 S3 S1 S2 S3 S1
output y: 0  0  0  1  0  0  1
\end{lstlisting}

重叠检测正是靠“匹配后退回 \code{S1} 而非 \code{S0}”实现的。以输入 11011 为例：第 4 位到达时检测到第一个 1101，机器停在 \code{S1}；第 5 位的 1 随即把它推进 \code{S2}——第二个匹配的前缀 11 已经备好。若输入继续 01，即完整序列 1101101，第二个 1101 在第 7 位被认出，它与第一个匹配共享了第 4 位的那个 1。倘若错误地让机器匹配后回 \code{S0}，这个共享位就被白白丢弃，1101101 中的第二个 1101 将被漏检。

最后说明为什么用 Mealy 型：检测输出依赖“状态加当前输入”，比 Moore 型少一个 DONE 状态，同拍给出结果，响应快一个周期。代价是 \code{y} 直接跟着 \code{x} 走，\code{x} 上的毛刺会原样出现在 \code{y} 上。若 \code{y} 只被后级在时钟沿采样，提前一拍是净收益；若 \code{y} 要驱动写使能这类危险控制线，应当再寄存一拍，用稳定性换掉这一拍的提前量。四个状态用 2 位二进制编码恰好用完，没有空洞编码。

\topic{六、从时序约束到检查方法}

同步设计不能仅以“这个状态机功能正确”作为结论。功能正确说明下一状态逻辑在稳定输入下给出正确答案；时序正确说明这些答案能在规定边界前到达，并且不会过早冲入保持窗口。二者缺一不可。分析一条关键路径，应当把它写成四个对象：源状态元件，即哪个触发器在发射沿后改变 Q，不能只写信号名而不写状态边界；数据路径，即结果经过哪些门、mux、加法器、译码和连线，不能只看逻辑级数而忽略布线和负载；目标状态元件，即哪个触发器在捕获沿采样 D，不能没有采样窗口；时钟关系，即周期、偏斜、抖动和工程余量，不能用理想同相时钟代替真实时钟树。

建立约束和保持约束检查的是同一条路径的两个方向。\hwkey{建立约束}关心最慢数据是否赶得上下一个采样沿；\hwtiming{保持约束}关心最快数据是否过早进入同一个采样沿后的保持窗口。工程上常见的错误，是只盯着最高频率而忽略保持时间。事实上，降低频率可以缓解建立问题，却不一定修复保持问题；保持问题往往需要改变最小延迟、时钟偏斜或布局。

可以用一个寄存器到寄存器路径作为分析模板。若源寄存器 \code{Rsrc} 的 \code{tCQ_max}=0.8ns，组合逻辑最大延迟为 3.6ns，连线最大延迟为 0.7ns，目标寄存器建立时间为 0.4ns，不利偏斜和余量合计 0.5ns，则周期至少为 6.0ns。若系统要求 200MHz，即周期 5.0ns，该路径就不满足建立约束。正确处理方式是把组合逻辑拆段、优化 mux 和加法路径、调整布局或降低频率，而不是把余量从时序预算中删除。

\begin{lstlisting}
required period = 0.8ns + 3.6ns + 0.7ns + 0.4ns + 0.5ns
                = 6.0ns
target period   = 5.0ns
result          = setup violation
\end{lstlisting}

保持检查则使用最小延迟。若 \code{tCQ_min}=0.08ns，组合逻辑最小延迟为 0.02ns，连线最小延迟为 0.03ns，目标保持时间为 0.20ns，再叠加不利偏斜 0.08ns，则数据最早 0.13ns 到达，而目标要求至少 0.28ns 后才允许变化。这条路径存在保持违例。修复方法通常是在数据路径上增加受控延迟、调整寄存器相对位置、修改时钟树或让工具插入保持缓冲。它不是“频率太高”的问题。

\begin{longtable}{L{0.22\linewidth}L{0.34\linewidth}L{0.34\linewidth}}
\toprule
现象 & 优先检查 & 不应直接假设 \\
\midrule
高频下偶发错误 & 建立约束、关键路径、温度和电压角落 & 逻辑表达式一定写错 \\
任何频率都偶发错误 & 保持约束、异步输入、复位释放、亚稳态 & 降低频率一定有效 \\
复位后偶发失控 & 复位释放同步、非法状态恢复、安全默认输出 & 只要复位电平正确就足够 \\
外部按钮偶发多次触发 & 消抖、同步、脉冲宽度、状态机边界 & 按钮已经是数字信号 \\
\bottomrule
\end{longtable}

复位要求同样要写成明确的约定。首先列出哪些寄存器必须复位：控制状态、计数器、使能输出、错误状态、协议状态通常需要复位；纯数据缓冲如果写使能受控，未必每一位都需要复位。其次说明复位期间输出是否安全，例如电机使能、写存储器使能、片选信号和外部驱动必须进入禁止态。最后说明复位如何释放：异步拉入可以快速进入安全状态，同步释放可以避免同一状态机内不同触发器在不同边界开始运行。

按对象核对默认值：控制状态机复位到 \code{IDLE} 或 \code{RESET_WAIT}，只能有一个合法入口状态；写使能复位为 0，复位期间不能写任何寄存器或存储器；外部驱动进入禁止输出或高阻安全态，不能误启动执行机构；错误状态被清除或进入诊断态，要能区分复位清除和真实故障；跨域同步器进入不会触发事件的状态，释放后不能产生伪脉冲。

跨时钟边界要按对象分类。单 bit 电平可以使用两级同步；窄脉冲应先转换为持续足够长的电平事件，或使用握手协议；多 bit 数据不能逐位同步后直接合并，因为各位可能来自不同采样边界；计数器跨域常用格雷码，异步 FIFO 则用读写指针、同步后的指针和空满判断来建立跨域约定。这里的关键不是背某一种结构，而是承认两个时钟域之间没有共同的“同时”概念，必须让协议定义何时有效、何时可采样、何时完成。

\begin{chipcase}
74HC74 这类双 D 触发器可以作为本章的门槛案例：它把 D、CLK、Q、复位和置位引脚明确暴露出来，迫使读者区分数据输入、采样边界和异步控制。74HC161 这类同步计数器则展示“当前状态经加法路径反馈到下一状态”的标准结构。74HC164、74HC595 等移位寄存器进一步说明，多 bit 状态可以按固定连线和时钟边界逐步移动。使用这些芯片时，必须检查电源去耦、输入默认电平、时钟边沿质量、复位极性和输出负载；否则纸面上的状态表无法保证板上行为。
\end{chipcase}

\topic{专题：把本章改写成纸面设计审查}

本章概念可以用常见 74HC/74HCT 逻辑芯片搭成低速实验板，但本书不要求读者实际搭建。更适合本书目标的训练，是把这些器件当作可审查的真实案例：根据数据手册写出电源、输入默认值、复位极性、时钟边界、输出负载和故障表现。若读者没有面包板和仪器，也可以完成同样的纸面设计审查；若读者选择实际搭建，则必须限制在低压直流逻辑范围内，例如 3.3V 或 5V 数字电路，不得把这些小信号电路直接接到市电、电机、继电器线圈或其他大功率负载。LED 必须串联限流电阻，CMOS 输入不能悬空，每片芯片的 VCC 与 GND 之间应靠近芯片放置约 0.1uF 去耦电容，实验板和信号源必须共地。

数据手册是搭电路时的依据。正文用信号名说明连接关系，而不依赖某一种封装的脚号；实际搭建时必须以手中芯片和封装的数据手册为准。低有效输入常写作 \code{reset_n}、\code{MR_n}、\code{OE_n} 或 \code{RD_n}，含义是拉低有效、拉高无效。若把低有效复位误当成高有效复位，电路会长期处于复位或完全失去复位保护。

\begin{longtable}{L{0.24\linewidth}L{0.3\linewidth}L{0.36\linewidth}}
\toprule
器件 & 本章对应概念 & 搭建时的重点 \\
\midrule
74HC14 & 施密特输入、慢边沿整形、按钮消抖前后边界 & 输出反相；语义校正信号取自第一级，需要与输入同相的低有效信号时再串联第二级 \\
74HC74 & 正沿 D 触发器、两级同步器、事件保持位 & 异步置位/复位脚必须接到确定无效电平或受控复位 \\
74HC161/74HC163 & 4-bit 同步计数器、当前状态加一、进位级联 & 区分异步复位和同步复位，计数使能与并行装载不能悬空 \\
74HC157 & 数据选择器、写使能保持、\code{D=mux(we,new,Q)} & 让时钟连续运行，用数据路径决定写入或保持 \\
74HC595 & 串入并出移位寄存器、输出锁存、三态输出 & 区分移位时钟和锁存时钟，输出负载要受限 \\
NE555 或晶振模块 & 低速实验时钟 & 时钟边沿应再经过施密特整形，按钮时钟必须消抖 \\
\bottomrule
\end{longtable}

如果只做纸面训练，不需要准备物料；需要准备的是三张表：器件功能表、引脚电平表和故障定位表。器件功能表说明每颗芯片在本章概念中扮演什么角色；引脚电平表说明每个输入脚在复位、空闲和动作时处于什么电平；故障定位表说明某个现象优先怀疑哪条边界。下面保留实际搭建所需物料，是为了让这些纸面表格有真实约束来源，而不是暗示读者必须做实物实验。

具体到本章实验，常用芯片包括 \code{74HC14}、\code{74HC74}、\code{74HC161} 或 \code{74HC163}、\code{74HC157}、\code{74HC595}；常用阻容范围包括 330R 到 1k 的 LED 限流电阻、4.7k 到 10k 的上拉或下拉电阻、100nF 去耦电容，以及 100nF 到 1uF 的按钮滤波电容。逻辑分析仪或示波器不是第一步必需品，但一旦要判断“按一次为什么计了两次”“复位释放为什么产生假事件”“锁存时钟为什么没有更新输出”，它们就比肉眼观察 LED 更可靠。

\begin{longtable}{L{0.22\linewidth}L{0.34\linewidth}L{0.34\linewidth}}
\toprule
物料 & 用途 & 检查要点 \\
\midrule
稳定直流电源 & 给所有逻辑芯片提供同一电源域 & 电压在芯片允许范围内，实验板各处共地 \\
0.1uF 去耦电容 & 抑制芯片开关瞬间的局部电源扰动 & 每片芯片靠近 VCC/GND 放置一颗 \\
上拉/下拉电阻 & 给按钮、复位和未驱动输入确定默认值 & 松开按钮或开关断开时节点不悬空 \\
LED 与限流电阻 & 低速观察 Q、同步级和计数输出 & LED 电流不超过芯片输出能力 \\
按钮与 RC 滤波 & 产生人工事件和复位请求 & 按钮边沿进入 74HC14 前已有默认电平 \\
逻辑分析仪或示波器 & 观察边沿、抖动、短脉冲和时序关系 & 能同时看到 \code{raw/pressed/sync/clk/reset_n} \\
\bottomrule
\end{longtable}

上电前应先按连接检查表逐项确认。许多低速实验失败不是因为状态机概念错误，而是因为 VCC/GND 接反、低有效输入悬空、输出直接短接、不同电压域硬连或 LED 没有限流。对 CMOS 逻辑而言，“未使用输入不接”不是中性状态，而是让输入节点暴露给噪声和漏电流；它可能让芯片耗电增大、输出随机变化，甚至把后级状态机带入错误状态。

\begin{longtable}{L{0.22\linewidth}L{0.38\linewidth}L{0.3\linewidth}}
\toprule
检查项 & 必须确认 & 常见后果 \\
\midrule
电源脚 & 每片芯片的 VCC、GND 与数据手册一致 & 芯片发热、无输出或永久损坏 \\
去耦 & 每片芯片电源脚旁有 0.1uF 电容 & 时钟沿或输出翻转时产生偶发错误 \\
低有效脚 & \code{reset_n/MR_n/OE_n} 等处于确定有效或无效电平 & 电路长期复位、输出高阻或随机启动 \\
未用输入 & 接到确定 0 或 1，不能悬空 & LED 随机闪烁，计数器误跳变 \\
输出负载 & LED 串电阻，不能多个推挽输出硬短接 & 输出电流过大或总线冲突 \\
电压域 & 3.3V 与 5V 芯片互连满足输入阈值和绝对最大额定值 & 高电平不被识别，或 5V 损伤 3.3V 输入 \\
\bottomrule
\end{longtable}

搭建顺序也应分层进行。第一步只接电源、去耦、一个 LED 和限流电阻，确认电源轨、极性和 LED 方向。第二步只接 74HC14 与按钮，验证按钮松开和按下时输出有确定变化。第三步接 74HC74 两级同步器，用低速时钟观察 \code{button_meta} 和 \code{button_sync} 的周期差。第四步再接 74HC161/74HC163 或 74HC595。若一次把所有芯片都插上，故障会混在一起，读者很难判断问题属于电源、按钮、时钟、复位还是状态逻辑。

第一个可搭电路是“按钮进入同步状态机”的前半段。按钮一端接 GND，另一端通过约 10k 电阻上拉到 VCC，形成低有效原始输入 \code{button_raw_n}。在按钮节点到 GND 之间可并联 100nF 左右电容作为低速消抖起点；随后进入 74HC14。由于 74HC14 是反相施密特触发器，按钮按下时 \code{button\_raw\_n}=0，经过第一级反相后即得到极性校正后的 \code{button_pressed}（1 表示按下）；需要保留与输入同相的低有效信号时，再取第二级输出。这个电路对应外部事件链路中的“默认电平、滤波、语义归一化”。

然后用 74HC74 搭两级同步器。第一触发器的 D 接 \code{button_pressed}，CP 接系统时钟，异步置位保持无效，异步复位接受控的 \code{reset_n}；第一触发器的 Q 记为 \code{button_meta}。第二触发器的 D 接 \code{button_meta}，CP 接同一个系统时钟，Q 记为 \code{button_sync}。状态机、计数器或后级逻辑只能使用 \code{button_sync}，不能直接使用 \code{button_raw_n} 或 \code{button_meta}。若要观察现象，可把 \code{button_pressed}、\code{button_meta}、\code{button_sync} 分别通过较大限流电阻接到 LED；在 1Hz 到 5Hz 的慢时钟下，读者能看到同步输出只在时钟边界变化。

\begin{lstlisting}
button_raw_n -> 74HC14 -> button_pressed

74HC74 stage 1:
  D  = button_pressed
  CP = clk
  Q  = button_meta

74HC74 stage 2:
  D  = button_meta
  CP = clk
  Q  = button_sync
\end{lstlisting}

该实验的观察重点不是“按键能点亮 LED”这么简单，而是要记录四个事实。第一，按钮松开时原始节点有确定默认电平。第二，施密特输入之后的边沿不再长时间停留在中间电压。第三，\code{button_meta} 可能比 \code{button_sync} 早一个周期变化，后级不得使用它。第四，复位拉低后两个触发器进入不会伪触发的默认值，复位释放后输出只能在时钟沿重新进入有效状态。

第二个可搭电路是 74HC161 或 74HC163 计数器。把 CP 接低速时钟，计数使能接有效电平，并行装载保持无效，复位接到受控 \code{reset_n}，Q0 到 Q3 分别通过限流电阻接 LED。此时四个 LED 显示的是同一个时钟边界保存的 4-bit 状态，而不是四个彼此独立的灯。每来一个有效时钟沿，内部寄存器一起更新，组合加一逻辑准备下一个值。若使用 74HC161，应特别注意其复位口是异步复位；若教材实验要强调“同步释放”，可以让外部复位先进入本时钟域同步器，再驱动后级控制逻辑，或选用同步复位方式更清楚的计数器。

这个计数器实验可以验证三件事。第一，把时钟停住后，Q0 到 Q3 保持不变，说明状态由触发器保存。第二，把计数使能关掉后，时钟仍然可以继续进入芯片，但状态保持，说明写使能和门控时钟不是一回事。第三，若把并行装载打开并在 D0 到 D3 上给出固定值，计数器会在规定边界装载该值，说明状态元件可以执行“保持、加一、装载、清零”等不同下一状态选择。

第三个可搭电路是 74HC595 输出链路。它内部有移位寄存器和输出锁存器，通常使用 \code{DS} 作为串行数据，\code{SHCP} 作为移位时钟，\code{STCP} 作为输出锁存时钟，\code{MR_n} 作为低有效清零，\code{OE_n} 作为低有效输出使能。每给 \code{SHCP} 一个有效边沿，串行 bit 进入移位寄存器；连续移入 8 bit 后，再给 \code{STCP} 一个边沿，八个输出 Q0 到 Q7 才一起更新。这个器件非常适合说明“内部状态移动”和“外部可见输出提交”是两个边界。

\begin{lstlisting}
for each bit:
  set DS stable
  pulse SHCP

after 8 bits:
  pulse STCP
  Q0..Q7 update together
\end{lstlisting}

若用按钮手动提供 \code{DS}、\code{SHCP} 或 \code{STCP}，这些按钮信号仍然要先经过上拉/下拉、消抖和施密特整形。未经消抖的手动时钟可能一次按压产生多个边沿，使移位寄存器多移几位。若用微控制器驱动 74HC595，要保证电压域兼容；5V 逻辑输出不能随意接到只允许 3.3V 的输入。Q0 到 Q7 可以驱动小电流 LED 作为观察负载，若要控制继电器、电机或更大负载，必须增加合适的驱动器件、续流保护和隔离策略，不能把 74HC595 输出脚当作功率开关。

第四个可搭电路用 74HC157 和 74HC74 说明写使能。把 74HC74 的 Q 反馈到 74HC157 的一个输入，把新数据 \code{new_bit} 接到另一个输入，用 \code{write_enable} 选择哪一路送到 74HC74 的 D。时钟始终送到 74HC74；当 \code{write_enable=0} 时，D 得到旧 Q，下一个边界写回旧值，表现为保持；当 \code{write_enable=1} 时，D 得到 \code{new_bit}，下一个边界写入新值。这个实验应当与“用 74HC08 把时钟和 enable 相与”明确区分。后者即使在低速实验中似乎能工作，也会把 enable 的毛刺、传播延迟和不对称边沿带进时钟路径；它适合用逻辑分析仪作为反例观察，不适合作为后级状态机的正常时钟。

\begin{longtable}{L{0.22\linewidth}L{0.36\linewidth}L{0.32\linewidth}}
\toprule
实验 & 观察信号 & 预期现象 \\
\midrule
按钮同步器 & \code{raw/pressed/meta/sync/reset_n/clk} & 后级只在时钟边界看到同步后的事件 \\
4-bit 计数器 & \code{clk/reset_n/enable/Q0..Q3} & 关使能时状态保持，开使能时每周期加一 \\
74HC595 输出链路 & \code{DS/SHCP/STCP/MR_n/OE_n/Q0..Q7} & 移位期间输出不乱跳，锁存边界后一起更新 \\
写使能寄存器 & \code{new_bit/write_enable/D/Q/clk} & 时钟连续存在，是否更新由数据路径决定 \\
\bottomrule
\end{longtable}

故障定位应从可观察事实开始，而不是先猜逻辑表达式。若计数器按一次按钮跳多格，优先检查按钮抖动、手动时钟是否经过 74HC14、是否把按钮直接当 CP 使用。若 LED 随机闪烁，优先检查 CMOS 输入是否悬空、复位或输出使能是否处于确定电平、面包板电源轨是否断开。若 74HC595 的 Q0 到 Q7 始终不变，应检查 \code{OE_n} 是否允许输出、\code{MR_n} 是否保持无效、\code{STCP} 是否真的收到锁存边沿，而不是只检查 \code{SHCP}。若复位后出现一次假启动，应检查同步器和 pending 位的复位默认值，以及复位释放是否在本时钟域边界发生。

\begin{longtable}{L{0.23\linewidth}L{0.34\linewidth}L{0.33\linewidth}}
\toprule
失败现象 & 优先检查 & 修复方向 \\
\midrule
按一次计多次 & 按钮抖动、手动时钟、74HC14 整形 & 增加 RC 消抖，整形后再进时钟域 \\
LED 随机闪烁 & 悬空输入、复位脚、输出使能脚、电源轨 & 给所有输入确定默认电平并检查共地 \\
复位后假事件 & 同步器复位值、pending 位默认值、释放边界 & 异步拉入、同步释放，事件位复位为 0 \\
74HC595 输出不变 & \code{OE_n/MR_n/STCP/SHCP} 与数据稳定时间 & 区分移位时钟和锁存时钟，固定控制脚电平 \\
写使能保持失败 & mux 选择脚、Q 反馈路径、D 采样边界 & 观察 \code{write_enable/D/Q/clk} 的周期关系 \\
芯片明显发热 & 电源极性、输出短路、总线冲突、负载电流 & 立即断电，按数据手册逐脚复查 \\
\bottomrule
\end{longtable}
\topic{专题：从数据手册到可靠接线}

真实芯片不能只按逻辑符号连接，还要按数据手册给出的电气边界连接。数据手册通常有两类容易混淆的表。\term{绝对最大额定值}{absolute maximum ratings} 是器件不应超过的损伤边界，不是推荐工作点；长期按这个边界运行并不可靠。\term{推荐工作条件}{recommended operating conditions} 才是设计应采用的电源、电压、温度、输入上升下降时间和时钟条件范围。实验报告应记录使用的是哪一种条件，而不能把“没有立刻烧毁”当成合格。

\begin{longtable}{L{0.23\linewidth}L{0.36\linewidth}L{0.31\linewidth}}
\toprule
数据手册项目 & 读法 & 设计用途 \\
\midrule
绝对最大额定值 & 超过可能损伤器件 & 用来划红线，不用来定正常工作点 \\
推荐工作条件 & 厂商承诺正常工作的范围 & 决定 VCC、电压域、温度和输入边沿要求 \\
\code{VIH/VIL} & 输入被承认为 1 或 0 的电压门限 & 判断 3.3V 输出能否驱动 5V 输入 \\
\code{VOH/VOL} & 在给定输出电流下仍能保证的输出电平 & 判断下一级是否能可靠识别高低电平 \\
输出电流 & 单脚和整片芯片的源/灌电流限制 & 决定 LED 电阻、负载数量和是否需要驱动器 \\
传播延迟 & 输入变化到输出有效的最坏时间 & 计入时序预算，不能只看典型值 \\
最小时钟脉宽 & CP、SHCP、STCP 等高/低电平至少持续多久 & 判断按钮、NE555 或逻辑脉冲是否太窄 \\
最高工作频率 & 在规定电压、负载和温度下的频率上限 & 面包板实验应留出明显余量 \\
\bottomrule
\end{longtable}

电压兼容要用最坏条件判断。若一个 3.3V 微控制器输出接到 5V 供电的 74HC 输入，不能只看示波器上有 3.3V 高电平，还要查 5V 条件下该输入要求的 \code{VIH_min}。如果 \code{VIH_min} 高于 3.3V 输出在负载下能保证的 \code{VOH_min}，这个连接就没有静态高电平余量，即使短时间看似能工作，也可能随温度、电源和芯片批次变化而失败。74HCT 系列常用于 5V 系统接收 TTL 电平标准的输入，是因为它的输入门限更适合较低高电平；但 74HCT 输出仍由自身 VCC 决定，不能把 5V 输出直接送入不耐 5V 的 3.3V 输入。

\begin{lstlisting}
static high-level margin:
  driver VOH_min >= receiver VIH_min + margin

static low-level margin:
  driver VOL_max <= receiver VIL_max - margin
\end{lstlisting}

扇出与负载也要按电流和电平一起看。\term{扇出}{fan-out} 指一个输出能可靠驱动多少个输入。CMOS 输入静态电流很小，所以低速实验中一个 74HC 输出常能驱动多个 CMOS 输入；但这不等于它能随意驱动多个 LED、长线或电容负载。LED 会消耗直流输出电流，电流过大时 \code{VOH} 会下降或 \code{VOL} 会升高，后级输入余量随之变小。长杜邦线和多个输入会增加电容负载，使边沿变慢；边沿太慢又可能让普通 CMOS 输入在门限附近停留过久，造成多次翻转或额外功耗。因此，观察 LED 应使用较大限流电阻，真实负载应使用专门驱动器或晶体管/MOSFET 级。

\begin{longtable}{L{0.22\linewidth}L{0.36\linewidth}L{0.32\linewidth}}
\toprule
负载类型 & 风险 & 正确处理 \\
\midrule
多个 CMOS 输入 & 电容变大，边沿变慢 & 降低频率、缩短连线，必要时加缓冲器 \\
LED 观察负载 & 输出电流过大拉低高电平或抬高低电平 & 串较大限流电阻，优先低电流观察 \\
长杜邦线 & 串扰、反射、接触不良和地线回路 & 低速实验、短线、共地、必要时示波器确认 \\
继电器或电机 & 电流和反灌电压远超逻辑芯片能力 & 使用驱动管、续流二极管、独立供电和隔离策略 \\
多个推挽输出相连 & 一个拉高、一个拉低会形成短路 & 只用三态总线或开漏/开集加上拉结构 \\
\bottomrule
\end{longtable}
时钟质量决定状态边界是否真实存在。一个合法时钟不仅要有频率，还要满足最小高电平时间、最小低电平时间、上升/下降时间和边沿单调性。NE555 可以产生低速时钟，但输出边沿、占空比和电源扰动仍需检查；按钮可以产生人工事件，但未经消抖不应直接作为计数器 CP；晶振模块边沿质量较好，但频率可能高到面包板布线、LED 观察和手工调试都跟不上。手工实验建议先用 1Hz 到 5Hz 的可观察时钟验证状态，再逐步提高频率，并在提高频率前去掉重负载 LED 或改用缓冲观察点。

\begin{longtable}{L{0.22\linewidth}L{0.36\linewidth}L{0.32\linewidth}}
\toprule
时钟来源 & 适合用途 & 额外检查 \\
\midrule
按钮单步 & 教学观察、单周期推进 & 必须消抖并经施密特整形，不直接驱动 CP \\
NE555 & 低速自动运行 & 检查占空比、电源去耦、边沿和最小脉宽 \\
晶振模块 & 稳定周期时钟 & 频率、电压域、输出负载和面包板布线余量 \\
微控制器 GPIO & 可编程时钟或测试脉冲 & 电压兼容、脉宽、地线和软件抖动 \\
\bottomrule
\end{longtable}

面包板适合低速验证概念，不等同于可量产硬件。面包板触点有接触电阻和寄生电容，长杜邦线会带来额外电感和串扰，电源轨可能在中间断开，地线回路会让几个芯片看到不同的瞬时参考电位。这些问题在 1Hz LED 实验里可能不明显，在更快边沿和更多芯片同时翻转时就会变成偶发故障。将来进入 CPU 数据通路、总线和存储器章节时，读者应记住：面包板证明“逻辑关系能被观察”，不能证明“任意频率和任意规模都可靠”。

安全边界也要明确。本章实验只验证低压逻辑信号，不验证功率驱动。若要让逻辑电路控制继电器、电机、灯带或其它外部负载，必须增加驱动晶体管或 MOSFET、栅极/基极电阻、续流二极管或 TVS、独立电源退耦、共地或隔离策略，并重新评估电流、热、反灌电压和失效状态。把 74HC595、74HC74 或 74HC161 的输出脚直接当作功率开关，是把逻辑输出误用为能量输出。

下面是一份完整故障报告样例。现象：用按钮给 74HC161 提供单步时钟，按一次按钮时计数器从 3 跳到 5。定位步骤一，逻辑分析仪同时观察按钮原始节点 \code{button_raw_n}、74HC14 输出 \code{button_pressed} 和计数器 \code{CP}；发现一次按压期间 \code{CP} 出现两个上升沿。定位步骤二，断开计数器，只观察按钮链路；发现 RC 时间常数过小且按钮释放也产生窄脉冲。修复动作：增大消抖电容，保证按钮先进入 74HC14，再由整形后的信号产生单步请求；若该请求进入状态机，则继续经过 74HC74 同步。复测结果：一次按压期间 \code{CP} 只有一个有效边沿，计数器从 3 变为 4。结论：故障不是 74HC161 的加一逻辑错误，而是手动时钟没有形成唯一边界。

实验报告应像一份小型工程记录，而不是只写“现象正确”。每个实验至少记录芯片型号和系列、电源电压、时钟来源、复位极性、每个低有效控制脚的默认电平、异步输入如何进入同步器、观察到的信号、失败现象和修复动作。若使用 74HCT 与 74HC 混合，还要记录电平兼容依据。一般而言，5V 输出不应直接接入只允许 3.3V 的输入；3.3V 输出能否可靠驱动 5V 供电的 74HC 输入，也不能凭“看起来能亮”判断，而应查数据手册中的输入高电平阈值。74HCT 的输入阈值更接近 TTL 电平标准，常用于 5V 系统中接收较低高电平，但输出仍受自身供电影响。

\begin{longtable}{L{0.22\linewidth}L{0.34\linewidth}L{0.34\linewidth}}
\toprule
报告项 & 应记录的内容 & 判断依据 \\
\midrule
物料与版本 & 芯片完整型号、逻辑系列、电源电压 & 可回到数据手册复查电气条件 \\
连接约定 & VCC/GND、复位、置位、输出使能、未用输入 & 没有悬空控制脚和危险负载 \\
时钟与复位 & 时钟来源、频率范围、复位拉入和释放方式 & 状态只在边界变化，复位后无伪事件 \\
异步输入 & 原始引脚、消抖、施密特整形、同步级数 & 后级只使用同步后的信号 \\
观测结果 & LED、逻辑分析仪或示波器看到的关键节点 & 现象能解释到具体边界和信号名 \\
故障处理 & 失败现象、定位步骤、修复前后对比 & 修复不是偶然重插线，而有明确原因 \\
\bottomrule
\end{longtable}

这些实验把“状态、时钟与同步边界”的概念应用到真实芯片上。完成实验后，读者应能为每个小电路写出一页完整的实验记录：使用了哪些芯片，电源和默认电平如何保证，哪些输入是异步的，进入了几级同步器，复位默认状态是什么，哪些输出只是观察 LED，哪些信号绝不能直接驱动危险负载。这样的记录比单纯画出逻辑关系更接近真实工程。

因此，本章的最终目标不是让读者背出触发器、寄存器、计数器和状态机的定义，而是能为一个同步小系统留下\hwevidence{可复查的设计记录}：状态表、复位表、时序预算、跨边界处理方式和异常恢复策略。后续 CPU 的 PC、寄存器堆、控制器、异常入口和设备中断，都将重复使用这套记录格式。

\topic{专题：setup/hold 不是公式，而是采样窗口}
触发器的 setup 和 hold 约束经常被背成两个时间参数，但它们本质上是在说明采样窗口。时钟沿到来前，D 输入要提前稳定一段时间；时钟沿之后，D 输入还要继续保持一小段时间。若输入在窗口内变化，触发器内部再生过程可能无法快速决定 0 或 1，输出延迟变长，甚至进入亚稳态。

\begin{longtable}{L{0.18\linewidth}L{0.42\linewidth}L{0.3\linewidth}}
\toprule
时序参数 & 原理含义 & 违反后果 \\
\midrule
\code{t_setup} & 时钟沿前 D 必须稳定的最短时间 & 采样值可能错误或输出延迟异常 \\
\code{t_hold} & 时钟沿后 D 必须保持的最短时间 & 新数据可能穿透到本周期采样 \\
\code{t_clkq} & 时钟沿后 Q 输出稳定所需时间 & 后级组合逻辑可用时间减少 \\
\code{t_pd} & 组合逻辑最坏传播延迟 & 决定下一级 setup 是否满足 \\
\code{skew} & 不同触发器实际看到时钟沿的差异 & 可能同时影响 setup 和 hold 裕量 \\
\bottomrule
\end{longtable}

单周期路径的基本时序关系可以写成：

\begin{lstlisting}
clk_period >= t_clkq_max + t_logic_max + t_setup + skew + margin
hold_margin = t_clkq_min + t_logic_min - skew - t_hold
\end{lstlisting}

第一条约束决定频率上限，第二条约束决定短路径是否太快。很多初学者只关心“逻辑太慢”，却忽略“逻辑太快”也可能破坏 hold。若同一时钟沿从前级触发器出来的数据太快到达后级，而后级仍处在 hold 窗口内，就可能把本该下一拍才采样的数据提前采进去。解决 hold 的方法通常不是降低频率，而是调整时钟树、增加最小延迟或改变寄存器布局。

\topic{专题：时序预算的一个完整数值例}

把采样窗口的两条不等式套到一组具体数字上。某设计目标频率 50MHz，即时钟周期 20ns。关键路径参数：源触发器 \code{tCQ\_max}=3ns、\code{tCQ\_min}=1.5ns，组合逻辑最坏延迟 9ns，目标触发器建立时间 \code{tSU}=2ns、保持时间 \code{tH}=0.5ns，时钟偏斜按最不利方向取 1ns。时序分析没有更多公式，只有一条纪律：建立检查用“最坏组合”，所有延迟取最大；保持检查用“最好组合”，所有延迟取最小——因为两个问题的敌人方向相反，一个怕数据来得太晚，一个怕数据来得太早。

先做建立检查。数据最晚在时钟沿后 3ns 离开源触发器，再花 9ns 穿过组合逻辑，第 12ns 才到达目标 D 端；目标要求它比下一个时钟沿（第 20ns）至少早 \code{tSU} 加偏斜共 3ns 就绪，即第 17ns 之前。12ns 早于 17ns，检查通过，余量 5ns。

\begin{lstlisting}
setup check (worst corner):
  arrival  = tCQ_max + t_comb_max = 3 + 9      = 12ns
  required = T - tSU - skew       = 20 - 2 - 1 = 17ns
  slack    = 17 - 12                           = +5ns  (pass)
\end{lstlisting}

再做保持检查。数据最快在时钟沿后 1.5ns 就可能离开源触发器；若这条路径几乎没有组合逻辑（例如移位寄存器的直接相邻位），新数据第 1.5ns 就到达目标 D 端。而目标的保持窗口要求旧数据至少维持到时钟沿后 \code{tH} 加偏斜共 1.5ns。1.5ns 对 1.5ns，余量恰好是 0：纸面上不违例，工程上叫边缘。

\begin{lstlisting}
hold check (best corner):
  earliest = tCQ_min + t_comb_min = 1.5 + 0    = 1.5ns
  required = tH + skew            = 0.5 + 1    = 1.5ns
  margin   = 1.5 - 1.5                         = 0ns   (marginal)
\end{lstlisting}

余量为 0 意味着任何一点工艺波动、温度漂移，或实际时钟树偏斜略大于估计值，都会把它推成违例。更要紧的是，降频救不了它：把 50MHz 降到 25MHz，建立余量从 5ns 变成 25ns，保持余量却仍是 0ns——保持窗口跟着时钟沿走，不跟着周期走。修复只有三条路：在数据路径上增加受控的最小延迟（插缓冲器）、减小时钟偏斜（修时钟树）、或换用 \code{tCQ\_min} 更大的触发器。这就是“降频治不了保持问题”的数值含义。

\begin{longtable}{L{0.16\linewidth}L{0.3\linewidth}L{0.26\linewidth}L{0.18\linewidth}}
\toprule
检查项 & 代入组合 & 计算 & 结论与余量 \\
\midrule
建立时间 & 最坏组合：\code{tCQ\_max}、组合最大延迟、\code{tSU}、不利偏斜 & 3+9+2+1=15ns，周期 20ns & 通过，余量 +5ns \\
保持时间 & 最好组合：\code{tCQ\_min}、组合最小延迟约 0、\code{tH}、不利偏斜 & 1.5ns 对 0.5+1=1.5ns & 边缘，余量 0ns，建议加最小延迟 \\
降频效果 & 周期从 20ns 改为 40ns & 建立余量升至 25ns，保持余量仍 0ns & 只治建立，不治保持 \\
\bottomrule
\end{longtable}

静态时序分析工具对设计中每一条寄存器到寄存器路径都做这两遍计算，报告里的 setup slack 与 hold slack 就是上面的 +5ns 和 0ns。手工算一遍的价值在于：看到报告时知道每个数字对应哪一段物理延迟，也知道该改逻辑、改时钟树还是改目标频率。

\topic{亚稳态只能降低概率，不能被完全消灭}
异步输入、跨时钟信号和机械按键都可能在触发器采样窗口内变化。触发器内部的再生电路可能暂时停在中间电平，随后随机解析为 0 或 1，这就是亚稳态。同步器不能保证“永不亚稳”，只能让亚稳态有更多时间衰减，使错误传播到后级的概率低到工程可接受。

\begin{longtable}{L{0.18\linewidth}L{0.42\linewidth}L{0.3\linewidth}}
\toprule
设计选择 & 作用 & 仍需注意 \\
\midrule
两级同步器 & 给第一级亚稳态多一个周期衰减 & 只能同步单 bit 电平，不能同步多 bit 总线值 \\
握手协议 & 用 valid/ready 或请求/应答跨边界 & 吞吐降低，但语义清楚 \\
异步 FIFO & 用双端口存储和格雷码指针跨时钟 & 指针同步和空满判断必须严格验证 \\
脉冲拉宽 & 让短事件至少覆盖目标时钟周期 & 过窄脉冲可能完全漏采 \\
\bottomrule
\end{longtable}

多 bit 总线不能简单给每一位各接两级同步器后直接使用，因为各位解析时间可能不同，接收端会看到不存在的混合值。正确做法通常是先在源时钟域保持数据稳定，再用握手同步“数据可用”事件；或者使用异步 FIFO 让写指针和读指针通过受控编码跨域。

\topic{复位要区分上电初始化、错误恢复和运行控制}
复位不只是“把所有寄存器清零”。有些寄存器必须复位到安全值，有些数据寄存器可以不复位，有些外设需要按电源和时钟顺序释放，有些跨时钟域需要各自同步释放。把复位当成普通控制信号乱用，会造成时钟域之间状态不一致、异步释放毛刺或组合路径被复位扇出拖慢。

\begin{longtable}{L{0.2\linewidth}L{0.36\linewidth}L{0.34\linewidth}}
\toprule
复位对象 & 建议策略 & 检查要点 \\
\midrule
控制状态机 & 复位到明确安全状态 & 复位释放后第一拍输出是否安全 \\
PC/入口状态 & 装入复位向量或启动状态 & 第一条取指地址是否确定 \\
数据寄存器 & 可按需要不复位，由控制路径保证无效时不使用 & 是否有 valid 位防止读旧值 \\
外设接口 & 先复位设备内部，再开放总线响应 & 复位期间是否会误写 MMIO \\
跨域复位 & 异步断言、同步释放或每域独立处理 & 释放边界是否满足目标时钟域时序 \\
\bottomrule
\end{longtable}

良好的复位规格应写清楚：复位信号有效电平、最短保持时间、时钟是否必须运行、复位释放后第一个可接受输入的时间、输出是否需要保持安全值。没有这些边界，系统 bring-up 时会把电源、时钟、复位、状态机和总线问题混成一团。

\topic{状态机验证要覆盖非法状态和输出时序}
状态机图只画正常转移是不够的。真实触发器可能因复位不完整、毛刺、单粒子翻转或设计错误进入未列出的编码。若非法状态没有恢复路径，控制器可能卡死或输出危险控制线。状态机输出还要区分 Moore 型和 Mealy 型：前者只依赖当前状态，后者还依赖输入，可能更快但更容易受输入毛刺影响。

\begin{longtable}{L{0.18\linewidth}L{0.42\linewidth}L{0.3\linewidth}}
\toprule
验证项 & 要检查什么 & 失败风险 \\
\midrule
复位覆盖 & 所有关键状态寄存器进入已知状态 & 上电随机进入危险路径 \\
非法状态恢复 & 未使用编码能回到安全状态或错误入口 & 状态机锁死或误写设备 \\
输出默认值 & 每个状态未显式赋值的控制线有安全默认 & 锁存器推断或保持旧控制线 \\
输入同步 & Mealy 输入是否来自同步边界 & 毛刺直接影响输出控制线 \\
进度性 & 等待状态是否有超时或退出条件 & 外设无响应时永久等待 \\
\bottomrule
\end{longtable}

这套检查表会在 CPU 控制器、内存等待、DMA 描述符、总线仲裁和中断控制器中反复出现。第二章的状态机不是孤立数字逻辑练习，而是后面所有“硬件会按阶段推进”的基础。

\topic{去抖和边沿检测要分成两件事}
机械按键输入常常同时暴露三个问题：异步、抖动和事件检测。异步意味着它不按本机时钟变化，需要同步；抖动意味着一次按下会产生多个跳变，需要滤除；事件检测意味着电平稳定后还要决定什么时候产生一个周期宽的“按下事件”。把三者混在一起，常会导致一次按键被计数多次，或按住不放时不断触发。

\begin{lstlisting}
button pipeline:
  raw_pin -> synchronizer -> debounce counter -> stable_level
  stable_level + previous_level -> one_cycle_press_event
\end{lstlisting}

\begin{longtable}{L{0.18\linewidth}L{0.42\linewidth}L{0.3\linewidth}}
\toprule
阶段 & 原理 & 预期现象 \\
\midrule
同步 & 两级触发器降低亚稳态传播概率 & raw 改变时内部不会产生多 bit 混合状态 \\
去抖 & 输入持续稳定 N 个周期才接受新电平 & 抖动期间 \code{stable_level} 不变化 \\
边沿检测 & 当前稳定电平与上一稳定电平比较 & 一次按下只产生一个事件脉冲 \\
事件消费 & 状态机或计数器只看事件脉冲 & 长按不会重复触发，除非规格要求连发 \\
\bottomrule
\end{longtable}

这个小例子把第二章与第一章连接起来：第一章能用示波器看到按键抖动，第二章用同步状态机把抖动变成稳定事件，后面整机章再把该事件映射成 GPIO 状态位或中断请求。

\section*{章末练习}

本节练习要求把状态边界写清楚。答案必须说明哪个值在组合路径上传播，哪个值在时钟沿保存；若涉及外部输入、复位或跨时钟边界，还必须说明同步和默认安全状态。

\begin{longtable}{L{0.18\linewidth}L{0.46\linewidth}L{0.26\linewidth}}
\toprule
任务 & 要完成的产物 & 预期结果 \\
\midrule
入门：画 Q 波形 & 给定一段 D 和 CLK 波形，分别画出 D 锁存器（LE=CLK）和 D 触发器（上升沿）的 Q & 锁存器 Q 在 CLK=1 全程跟随 D，触发器 Q 只在上升沿变化 \\
入门：填 SR 真值表 & 补全 $\overline{S}$/$\overline{R}$ 四种组合下 Q 的值，并标出禁行 & 禁行说明互补被破坏，不是“存了 1”；释放顺序决定恢复结果 \\
入门：画计数波形 & 从初值 9 开始画 6 个周期的 Q[3:0] 总线带和 Q0 波形 & 每个上升沿整体 +1，Q0 恰好是 CLK 二分频 \\
入门：主从结构 & 给定 D 与 CLK，填出主锁存器 M 和输出 Q 在各阶段的值 & M 只在 CLK=0 跟随 D，Q 只在上升沿复制 M \\
触发器路径时序分析 & 为一条寄存器到寄存器路径列出 \code{tCQ}、组合延迟、连线延迟、建立时间、偏斜和余量 & 能判断目标周期是否通过建立约束 \\
保持违例 & 构造一条最小延迟过短的路径，并写出修复方法 & 能区分建立问题和保持问题 \\
锁存器透明窗口 & 说明两个锁存器共用 enable 时可能出现的穿透风险 & 能解释为什么大系统更偏好清晰时钟边界 \\
复位释放 & 为一个三状态控制器写复位表和释放策略 & 复位期间危险输出关闭，释放在同步边界发生 \\
两级同步 & 说明外部按钮或中断线进入时钟域的处理流程 & 原始异步信号不能直接驱动状态机 \\
MTBF 指标 & 说明为什么同步器只能降低亚稳态风险，并列出影响 MTBF 的因素 & 能把时钟频率、事件频率、恢复时间和器件特性联系起来 \\
外部事件链路 & 为一个按钮或传感器写出从原始引脚到状态机消费的完整信号链 & 能区分电气处理、同步、边沿检测和事件保持 \\
事件保持 & 设计 \code{pending/ack} 结构，说明何时置位、何时清除 & 窄事件不能因状态机暂时不可接收而丢失 \\
多 bit 跨域 & 说明为什么 8-bit 总线不能逐位两级同步后直接使用 & 给出握手、FIFO 或格雷码方案 \\
脉冲同步 & 设计一个不会漏采窄脉冲的跨域事件方案 & 能说明事件保持到被确认 \\
时钟门控 & 比较写使能、数据 mux 保持、专用时钟门控和普通门控时钟 & 能说明为什么普通门电路不应直接切时钟 \\
复位域 & 为两个相关时钟域写复位释放和 ready 等待策略 & 能说明复位跨域也需要同步和安全默认值 \\
状态编码 & 比较二进制、one-hot 和格雷编码的用途与风险 & 能为未使用编码写安全恢复规则 \\
波形验证 & 列出验证一个同步小系统必须观察的信号 & 能从波形判断写使能、复位、pending/ack 是否正确 \\
CPU 连接 & 说明 PC、流水线寄存器、控制器和中断入口分别对应本章哪个概念 & 能把第二章边界模型映射到后续 CPU \\
按钮同步器实作 & 用 74HC14 和 74HC74 写出按钮到同步电平的连接方案 & 能区分原始引脚、消抖整形、第一级同步和第二级同步 \\
计数器实作 & 用 74HC161 或 74HC163 设计 4-bit LED 计数器 & 能说明时钟、复位、计数使能、并行装载和输出负载 \\
移位输出实作 & 用 74HC595 说明串行输入、移位时钟、锁存时钟和输出使能 & 能解释为什么八个输出在锁存边界一起更新 \\
写使能实作 & 用 74HC157 和 74HC74 搭一个可保持的 1-bit 寄存器 & 能说明为什么数据 mux 保持优于普通门控时钟 \\
实验板物料 & 为本章低速逻辑实验列出最小 BOM 和用途 & 能覆盖电源、去耦、默认电平、观察负载和测量工具 \\
接线检查 & 写出上电前检查表，覆盖低有效脚、未用输入、电压域和输出负载 & 能防止悬空、误复位、输出短路和电平不兼容 \\
故障定位报告 & 针对一个失败现象写出定位步骤和修复验证 & 能从现象追到具体信号边界，而不是只重插线 \\
数据手册读取 & 给出 74HC 实验前必须查的参数表和用途 & 能区分绝对最大额定值、推荐工作条件、输入输出电平和时序参数 \\
电压与扇出 & 判断一个输出能否可靠驱动后级输入、LED 或多路负载 & 能用 \code{VOH/VOL/VIH/VIL} 和输出电流说明余量 \\
时钟与面包板限制 & 说明按钮、NE555、晶振模块和面包板布线的边界 & 能解释最小脉宽、边沿质量、寄生效应和低速验证范围 \\
非法状态恢复 & 为状态机未使用编码指定下一状态和输出 & 非法状态不产生危险使能 \\
状态机输出 & 比较只依赖状态的输出和依赖输入的输出 & 能说明稳定性与响应速度的取舍 \\
芯片案例 & 用 74HC74、74HC161 或移位寄存器说明一个本章概念 & 绑定到真实引脚、复位、时钟和输出负载 \\
纹波与同步计数器 & 比较 4-bit 纹波计数器与 4-bit 同步计数器的时钟接法、延迟来源和译码风险 & 能说明纹波逐级延迟累积、译码有毛刺窗口，同步计数器把代价移到组合进位链 \\
LFSR 前八态 & 4 位 LFSR 每拍 \code{Q3} 取旧 \code{Q2}、\code{Q2} 取旧 \code{Q1}、\code{Q1} 取旧 \code{Q0}、\code{Q0} 取旧 \code{Q3} 与旧 \code{Q2} 的异或，从 0001 起列出前 8 个状态 & 序列正确，并能说明 0000 为何必须排除 \\
模 12 计数器 & 用 4-bit 寄存器与加一逻辑设计模 12 计数器（0 到 11 循环） & 归零发生在同步边界，无瞬时非法态 \\
Moore 与 Mealy 区分 & 判断交通灯控制器与 1101 序列检测器各属哪类，并说明输出稳定性差异 & 分类依据写成“输出是否依赖当前输入” \\
交通灯转移表 & 按东西绿 30s、黄 5s、全红 2s、南北绿 30s、黄 5s、全红 2s 的需求补全状态转移表与输出表 & 六状态闭环，\code{timer\_done} 为唯一转移条件，输出只依赖状态 \\
格雷码相邻性 & 核对交通灯六状态格雷编码 000、001、011、010、110、100 的每对相邻状态（含首尾闭环）是否只差一位 & 六次转移逐一核对，全部只差一位 \\
\bottomrule
\end{longtable}
\vspace{1.0em plus 0.35em}
\begin{block}[x=1cm,y=1cm]
  \node[hw lane, minimum width=4.6cm, minimum height=3.0cm, label={[hw label]above:推荐：数据端使能}] (leftbox) at (1.9,0) {};
  \node[hw compute, minimum width=1.2cm] (mux) at (0.6,0.35) {MUX};
  \node[hw state, minimum width=1.35cm] (ff) at (2.7,0.35) {D-FF};
  \node[hw control, minimum width=1.0cm] (we) at (0.6,-0.9) {WE};
  \draw[hw bus] (-1.25,0.35) node[left]{D} -- (mux);
  \draw[hw bus] (mux) -- (ff);
  \draw[hw return] (ff.east) -- ++(0.45,0) |- (1.65,1.05) -| (mux.north);
  \draw[hw bus] (we) -- (mux);
  \node[hw label] at (2.7,-0.9) {CLK 不被门控};

  \node[hw lane, minimum width=4.6cm, minimum height=3.0cm, label={[hw label]above:危险：门控时钟}] (rightbox) at (7.1,0) {};
  \node[hw control, minimum width=1.1cm] (andg) at (5.85,0.35) {AND};
  \node[hw state, minimum width=1.35cm] (ff2) at (7.9,0.35) {D-FF};
  \node[hw control, minimum width=1.0cm] (en) at (5.85,-0.9) {EN};
  \draw[hw danger] (4.55,0.35) node[left]{CLK} -- (andg);
  \draw[hw danger] (andg) -- node[above,hw label]{毛刺会变成边沿} (ff2);
  \draw[hw bus] (ff2.east) -- ++(0.45,0) node[right]{Q};
  \draw[hw danger] (en) -- (andg);
  \node[hw label, text width=9.4cm] at (4.5,-2.1) {左边只改变 D 端选择，时钟树保持干净；右边把 EN 放进时钟路径，EN 的毛刺可能被触发器当作真实时钟沿。};
\end{block}
\diagramnote{写使能 vs 门控时钟：写使能在数据路径上选择；门控时钟可能把 EN 的毛刺变成额外时钟沿。}
\vspace{0.8em plus 0.25em}

\section*{参考解答与要点}

\begin{longtable}{L{0.18\linewidth}L{0.46\linewidth}L{0.26\linewidth}}
\toprule
任务 & 参考解答 & 必须保留的要点 \\
\midrule
入门：画 Q 波形 & 锁存器 Q 在 CLK=1 期间完全复制 D 的每次变化；触发器 Q 只在每个上升沿复制当时刻的 D，其余时间保持。 & 触发器不是“更快的锁存器”，采样时刻完全不同。 \\
入门：填 SR 真值表 & 0,0 禁用；0,1 置 1；1,0 置 0；1,1 保持。从 0,0 同时释放到 1,1 时结果不可预测。 & 禁行是非法输入；可靠设计不产生 0,0 或不同时释放。 \\
入门：画计数波形 & Q[3:0] 总线带：9、10、11、12、13、14 逐沿跳变；Q0：1、0、1、0、1、0。 & 总线带整体跳变，不是逐位爬；Q0 每周期翻转。 \\
入门：主从结构 & M：CLK=0 期间等于 D，CLK=1 期间保持上升沿时刻的 D；Q：每个上升沿等于 M 当时值。 & D 在 CLK=1 期间的变化只影响 M 保持的旧值，不影响 Q。 \\
触发器路径时序分析 & 示例：\code{Tclk >= 0.8ns + 3.6ns + 0.7ns + 0.4ns + 0.5ns = 6.0ns}。若目标周期为 5.0ns，则建立约束失败。 & 使用最大延迟和不利偏斜，不得只用典型门延迟。 \\
保持违例 & 示例：最快路径 0.13ns 到达，而目标保持窗口要求 0.28ns 后才允许变化，存在保持违例。修复可插入保持缓冲、调整布局或修正时钟树。 & 降频不能可靠修复保持违例。 \\
锁存器透明窗口 & enable 有效期间，前一级输入变化可能穿过第一级锁存器、组合逻辑和第二级锁存器，使一次更新跨越多个预期阶段。 & 必须指出透明期间没有清晰提交边界。 \\
复位释放 & 控制状态复位到 \code{IDLE} 或 \code{RESET_WAIT}，写使能和危险输出为 0；复位可异步拉入，但在本时钟域用触发器同步释放。 & 复位值、安全输出、释放边界三者都要写明。 \\
两级同步 & 外部按钮先经过电气默认值和消抖，再进入两级触发器，状态机只使用 \code{button_sync}。 & 两级同步降低亚稳态扩散概率，但不替代消抖。 \\
MTBF 指标 & MTBF 表示亚稳态导致可见失败的平均时间尺度。目标时钟越快、异步事件越频繁，风险越高；第一级到第二级之间恢复时间越长，风险越低；器件恢复参数也会影响结果。 & 同步器是可靠性设计，不是绝对功能保证。 \\
外部事件链路 & 示例链路：\code{raw -> filtered -> pressed -> meta -> sync -> event -> pending -> FSM}。其中 raw/filtered 属于电气与消抖，meta/sync 属于时钟域同步，event/pending 属于目标时钟域事件协议。 & 每一层信号名应对应一个可测量、可观察的边界。 \\
事件保持 & \code{start_event} 置位 \code{start_pending}，状态机在 \code{IDLE} 且安全条件满足时产生 \code{start_ack} 清除 pending；复位或故障可按安全策略清除或拒绝事件。 & 单周期事件不应直接作为危险动作使能。 \\
多 bit 跨域 & 8-bit 总线逐位同步会采到混合时刻的 bit。若是数据传输，应使用 valid/ready 握手或异步 FIFO；若是计数指针，可用格雷码减少一次变化的 bit 数。 & 多 bit 数据必须有整体有效性约定。 \\
脉冲同步 & 可把源域窄脉冲转换为保持电平或翻转事件位，目标域同步后检测变化，再用确认信号让源域清除事件。 & 事件必须持续到目标域可见，不应仅依赖单周期脉冲宽度。 \\
时钟门控 & 写使能让时钟继续到达触发器，只控制是否保存新值；专用时钟门控用于功耗控制并要求 enable 无毛刺和安全锁存；普通 AND 门切时钟会产生短脉冲、额外边沿和不可控偏斜。 & 普通数据通路优先用写使能，不应随意组合门控时钟。 \\
复位域 & 每个时钟域内部同步释放复位；跨域依赖通过 ready、同步器或状态握手确认。同步器、pending 位和危险输出在复位期间应进入不会伪触发的默认值。 & 复位释放顺序属于设计约定，不能假设所有域同时恢复。 \\
状态编码 & 二进制编码节省触发器但译码较复杂；one-hot 译码直接但触发器多；格雷编码适合相邻状态或跨域指针。无论哪种编码，未使用编码都应导向安全状态并关闭危险输出。 & 状态编码选择要绑定资源、速度、跨域和安全要求。 \\
波形验证 & 至少观察 \code{clk/reset/source_q/dest_d/write_enable/state/sync1/sync2/pending/ack}，检查 D 是否在采样沿前稳定、写使能是否只在目标周期有效、复位释放是否产生伪事件。 & 只看最终寄存器值不足以证明同步边界正确。 \\
CPU 连接 & PC 是寄存器，流水线寄存器是阶段边界，控制器是状态机，写回使能是状态控制线，中断入口来自同步后的外部或内部事件。 & 后续 CPU 仍遵守“组合准备、边界提交”的模型。 \\
按钮同步器实作 & 按钮节点用上拉或下拉给出默认电平，经过 RC 和 74HC14 得到整形后的 \code{button_pressed}；再接入 74HC74 两级触发器，后级只使用第二级 Q。复位脚进入不会伪触发的默认状态。 & 必须写清楚每一级信号名，不能把 raw 引脚直接送入状态机。 \\
计数器实作 & 74HC161/74HC163 的 CP 接低速时钟，计数使能接受控有效电平，并行装载保持无效或由实验开关控制，Q0 到 Q3 通过限流电阻接 LED。复位极性按数据手册连接。 & 计数器体现同一时钟边界上的 4-bit 状态更新。 \\
移位输出实作 & 74HC595 的 \code{DS} 在 \code{SHCP} 边沿前稳定，8 次移位后再给 \code{STCP} 边沿，Q0 到 Q7 一起更新；\code{MR_n} 和 \code{OE_n} 保持在确定电平。 & 必须区分移位寄存器内部状态和输出锁存边界。 \\
写使能实作 & 74HC157 在旧 Q 和 \code{new_bit} 之间选择，输出送入 74HC74 的 D；时钟持续送入触发器，\code{write_enable=0} 时写回旧 Q，\code{write_enable=1} 时写入新值。 & 不应使用普通与门直接门控时钟作为正式方案。 \\
实验板物料 & 最小 BOM 应包含稳定 3.3V 或 5V 电源、面包板、目标 74HC/74HCT 芯片、0.1uF 去耦电容、上拉/下拉电阻、LED 与限流电阻、按钮、RC 滤波电容、杜邦线和万用表；逻辑分析仪或示波器用于观察边沿和短脉冲。 & 物料清单必须对应实验中的电源、输入、状态输出和测量需求。 \\
接线检查 & 上电前检查 VCC/GND、低有效复位和输出使能、未用输入默认值、LED 限流、输出是否短接、电压域是否兼容。发现芯片发热或输出异常时应立即断电复查。 & 检查表应能解释常见硬件失败，而不是只列器件名称。 \\
故障定位报告 & 示例：按一次计数器跳多格时，先观察按钮 raw，再看 74HC14 输出，最后看计数器 CP；若 raw 抖动而整形后仍有多个边沿，应增加消抖或避免按钮直接作为时钟。报告要写出修复前后观测差异。 & 失败分析必须对应到可观察信号和修复后的波形。 \\
数据手册读取 & 必须先查推荐工作条件确认 VCC、电压和输入边沿，再查 \code{VIH/VIL/VOH/VOL} 判断电平兼容，查输出电流判断 LED 和负载，查传播延迟、最高频率与最小脉宽判断时钟是否合格；绝对最大额定值只用于损伤边界。 & 不得把绝对最大额定值当作正常工作条件。 \\
电压与扇出 & 若驱动器 \code{VOH_min} 小于接收器 \code{VIH_min} 加余量，高电平不可靠；若 LED 电流过大，输出电平会被拉坏；一个输出驱动多个 CMOS 输入时还要考虑电容负载和边沿变慢。 & 判断必须同时看电压余量、电流限制和负载电容。 \\
时钟与面包板限制 & 按钮单步必须消抖和施密特整形，NE555 要检查占空比和脉宽，晶振模块要检查频率和电压域；面包板只适合低速概念验证，不能证明高频 CPU 级可靠性。 & 状态边界必须来自合格边沿，不来自慢爬升或多次抖动。 \\
非法状态恢复 & 未使用编码的下一状态导向 \code{IDLE}、\code{FAULT} 或复位等待态，输出逻辑默认关闭写使能、驱动使能和片选。 & 非法状态的输出也要安全，不应只写下一状态。 \\
状态机输出 & Moore 输出只依赖状态，通常更稳定；Mealy 输出还依赖输入，响应更快但可能把输入毛刺传到控制线。危险输出宜寄存或只由稳定状态产生。 & 取舍必须结合后级是否采样或驱动危险对象。 \\
芯片案例 & 以 74HC74 为例，D 是待保存数据，CLK 是采样边界，Q 是保存结果，异步复位/置位必须按数据手册极性连接；未用输入不能悬空。 & 真实芯片要检查引脚、电源、复位极性、时钟边沿和负载。 \\
纹波与同步计数器 & 纹波计数器用低位 Q 作高位时钟，第 k 位比第 0 位晚 k 个 \code{tCQ} 翻转，最高位在第 4 个 \code{tCQ} 才稳定，与第 0 位相差 3 个 \code{tCQ}；翻转窗口内译码器会看到瞬时错值，输出毛刺真实存在。同步计数器所有位共用同一时钟沿、同时更新，代价是进位链（高位使能等于低位全 1）随位数变长，成为组合关键路径。 & 纹波的风险在时钟不同步，同步的代价在进位逻辑深度。 \\
LFSR 前八态 & 0001、0010、0100、1001、0011、0110、1101、1010。逐拍按规则计算：高三位左移接收邻位，\code{Q0} 取旧 \code{Q3} 与旧 \code{Q2} 的异或。全 0 态的反馈恒为 0，一旦进入便永远锁死，故初态必须非零。 & 每拍只由旧态决定新态；0000 是陷阱态。 \\
模 12 计数器 & 现态为 1011（即 11）时，下一状态逻辑选 0000 而非加一：用比较器译码 1011 控制数据 mux 即可。若改用异步清零计数器（如 74HC161），需译码 1100 去拉清零脚，会产生短暂的 1100 毛刺态，且各位清零时刻不一致；同步清零（如 74HC163）译码 1011 后在时钟沿干净归零。 & 归零必须发生在时钟边界；异步清零方案要指出毛刺态。 \\
Moore 与 Mealy 区分 & 交通灯的灯色只由当前状态决定，属 Moore 型，输出在整个状态期间稳定；1101 检测器在 \code{S3} 且输入为 1 的转移上给出输出，属 Mealy 型，同拍响应但输入毛刺可直达输出。 & 分类看输出函数的自变量，与状态个数无关。 \\
交通灯转移表 & \code{EW\_GREEN} 到 \code{EW\_YELLOW} 到 \code{ALL\_RED\_EW} 到 \code{NS\_GREEN} 到 \code{NS\_YELLOW} 到 \code{ALL\_RED\_NS} 再回到 \code{EW\_GREEN}，转移条件均为 \code{timer\_done}=1，装载值依次为 5、2、30、5、2、30 秒；输出依次为绿/红、黄/红、红/红、红/绿、红/黄、红/红。\code{timer\_done}=0 时保持原状态。 & 状态数、转移条件、输出必须与需求逐条对应。 \\
格雷码相邻性 & 000 到 001 差最低位；001 到 011 差中间位；011 到 010 差最低位；010 到 110 差最高位；110 到 100 差中间位；100 到 000 差最高位。六次转移均只翻转一位。 & 闭环的最后一边（100 回到 000）也必须核对，不能只看前五次。 \\
\bottomrule
\end{longtable}

\begin{classicnote}
本章的标注对象是“边界”和“状态转移表”。仅说明某个值会变化并不充分，必须说明哪个时钟边界提交，提交前后的组合路径是什么。

\begin{itemize}
\item 纸上实验：画出一个 D 触发器前后各接一段组合逻辑，标出当前周期的旧 Q、组合路径上的 D、下一个时钟沿后的新 Q。
\item 时序分析：为一条路径写出 \code{Tclk >= tCQ + tcomb + tsetup + tskew}，再解释每一项对应哪个硬件现象。
\item 复位设计：复位要让控制状态进入已知安全点，但释放复位也要有同步边界，不能让状态机半边先跑。
\item 状态结构：寄存器是一组并排触发器；计数器是寄存器输出经过加法器反馈到输入；移位寄存器是固定连线选择了相邻 bit。
\item 控制结构：状态机输出可以只依赖状态，也可以依赖状态和输入；后者反应快，但更容易把输入毛刺传给控制线。
\item 检查题：若状态编码里出现未使用编码，电路上电或受干扰进入它时，下一状态逻辑应该怎样让机器回到安全状态？参考答案：未使用编码不应保持未定义，也不应产生危险输出；下一状态逻辑应把它导向复位、IDLE、FAULT 或其他安全诊断状态，同时关闭写使能、功率使能等危险控制线。若安全要求更高，还应记录错误或触发复位流程。
\end{itemize}

经典系统教材通常会把组合逻辑和状态逻辑分开建模：组合逻辑计算下一值，状态元件在边界保存下一值。后面 CPU 的 PC、寄存器、状态机和控制器都遵守这个模型。
\end{classicnote}
